% concatenated from template.tex
%\documentclass[lettersize,journal]{IEEEtran}
\documentclass[journal]{IEEEtran}
% \usepackage{amsmath,amsfonts}
% \usepackage{algorithmic}
% \usepackage{algorithm}
\usepackage{array}
% \usepackage[caption=false,font=normalsize,labelfont=sf,textfont=sf]{subfig}
\usepackage{textcomp}
\usepackage{stfloats}
\usepackage{url}
\usepackage{verbatim}
% \usepackage{graphicx}
\usepackage{cite}
\hyphenation{op-tical net-works semi-conduc-tor IEEE-Xplore}
% updated with editorial comments 8/9/2021

\usepackage{wrapfig}
\usepackage{subcaption}
\usepackage{graphicx}
%
\usepackage{mathptmx}      % use Times fonts if available on your TeX system
\usepackage{booktabs} % For formal tables
\usepackage{algorithm}
\usepackage{algpseudocode}
%+++++++++++++++++++++++++++++++++++++++++++++++++++++++++++++++++++++%
%+++ Math Packages and Commands
%+++++++++++++++++++++++++++++++++++++++++++++++++++++++++++++++++++++%
% Packages
\usepackage{latexsym}
\usepackage{amsmath,amssymb} % AMS
\usepackage{mathtools}       % many extension of amsmath
\usepackage{stmaryrd}        % \llbracket, \rrbracket
\usepackage{bm}              % bold font
\usepackage{bbm}              % bold font
\usepackage{amsthm}

% Mathematical commands that are recursively used
\providecommand{\argmax}{\operatornamewithlimits{argmax}} % argmax
\providecommand{\argmin}{\operatornamewithlimits{argmin}} % argmin
\providecommand{\limsup}{\operatornamewithlimits{limsup}} % limsup
\providecommand{\liminf}{\operatornamewithlimits{liminf}} % liminf
\DeclareMathOperator{\Tr}{Tr}     % trace
\DeclareMathOperator{\Var}{Var}   % variance
\DeclareMathOperator{\Cov}{Cov}   % covariance
\DeclareMathOperator{\Cond}{Cond} % condition number
\DeclareMathOperator{\Det}{Det} % determinant
\DeclareMathOperator{\diag}{diag} % diagonal
\DeclareMathOperator{\vect}{vec}  % vectorization
\DeclareMathOperator{\vech}{vech} % half vectorization
\providecommand{\N}{\mathbb{N}} % integer
\providecommand{\R}{\mathbb{R}} % real field
\providecommand{\E}{\mathbb{E}} % expectation
\providecommand{\T}{\mathrm{T}} % transpose
\providecommand{\rmd}{\mathrm{d}} % d for infinitesimal
\providecommand{\ind}[1]{\ensuremath{\mathbbm{1}_{\left\{#1\right\}}}} % indicator function
\renewcommand{\geq}{\geqslant} % geq
\renewcommand{\leq}{\leqslant} % leq

% Paired Delimiter from mathtools
% E.g., \inner{}{}, \inner*{}{}, \inner[Big]{}{}
% Or,   \norm{}, \norm*{}, \norm[Big]{}
\DeclarePairedDelimiterX{\inner}[2]{\langle}{\rangle}{#1, #2}
\DeclarePairedDelimiter{\norm}{\lVert}{\rVert}
\DeclarePairedDelimiter{\abs}{\lvert}{\rvert}

% Cleverref
\usepackage[colorlinks=true,linkcolor=blue,pdfborder={0 0 0}]{hyperref}% Hyperref
\usepackage{cleveref}

% Must be placed after cleveref
\newtheorem{theorem}{Theorem}[]
\newtheorem{proposition}[theorem]{Proposition}
\newtheorem{corollary}[theorem]{Corollary}
\newtheorem{lemma}[theorem]{Lemma}
\theoremstyle{definition}
% \newtheorem{definition}[theorem]{Definition}
% \newtheorem{remark}[theorem]{Remark}
% \newtheorem{example}[theorem]{Example}
% \newtheorem{assumption}[theorem]{Assumption}
\newtheorem{definition}[]{Definition}
\newtheorem{remark}[]{Remark}
\newtheorem{example}[]{Example}
\newtheorem{assumption}[]{Assumption}


%\usepackage{xcolor}
\usepackage{ifthen}

%------- Tool to display colored notes and changes in the draft and remove them easily
% The available commands are \del, \new, \nnew, \note, \nnote, \todo (same as \TODO).
% All commands take as optional argument the author.  The nn versions give more pronounced colors.
% Example:
%   This is\del{ an experimental version of}\note[Niko]{remark the space within the del command}
%   a {simple to use} LaTeX markup.\todo{improve this example}
%------- Meta-command to markup notes and changes and *remove those markups easily*
\newcommand{\markupdraft}[2]{% {#1: {color|display} command}{#2: desired color or text}
%  the next lines can be incommented, if respectively certain notes or coloring should disappear
\ifthenelse{\equal{#1}{display}}{#2}{}%                 % display only in draft version
\ifthenelse{\equal{#1}{color}}{\color{#2}}{}%           % colored only in draft (for \new command)
}
% ------ helper commands using \markupdraft, author is optional argument
\newcommand{\notecolored}[3][]{\markupdraft{display}{{\color{#2}\noindent[Note (#1): #3]}}}
\newcommand{\newcolored}[3][]{{\markupdraft{color}{#2}#3}%  % kept in the final print
\ifthenelse{\equal{#1}{}}{}{\markupdraft{display}{{\color{yellow!70!black}[#1]}}}}
% ------ end user markupdraft-commands, optional argument is the author:
\newcommand{\del}[2][]{{\markupdraft{display}{{\color{orange}[removed: ``#2''[#1]]}}}} % (to be) removed
%\newcommand{\new}[2][]{\newcolored[#1]{blue}{#2}}%  % kept in the final print
%\newcommand{\new}[1]{#1}%  % kept in the final print
\newcommand{\nnew}[2][]{\newcolored[#1]{red}{#2}}%  % kept in the final print
%\newcommand{\note}[2][]{\notecolored[#1]{green}{#2}}
\newcommand{\nnote}[2][]{\notecolored[#1]{magenta}{#2}}   % note with Anne's preferred default color
\newcommand{\future}[2][]{{\markupdraft{display}{{\color{magenta}[[future:]#2[#1]]}}}}  % don't keep in current final print, but for the future
\newcommand{\TODO}[2][]{\markupdraft{display}{{\color{red}~\\\noindent== TODO: #2 {\color{yellow}(#1)} ==\\}}}
\newcommand{\todo}[2][]{\markupdraft{display}{{\color{red}\noindent++TODO: #2 {\color{yellow}(#1)}++}}}
% ------ Final version
\renewcommand{\del}[2]{}  % make removed sentences invisible
%\renewcommand{\markupdraft}[2]{}  % remove all todo's, notes and coloring of changes
                            % like for the final version
% ------ Usage
\newcommand{\yohe}[1]{{\color{cyan}Y.A.: #1}}
\newcommand{\mori}[1]{{\color[rgb]{0.22, 0.71, 0.55} MORINAGA :  #1}}
% ============================================================================================ %
% ============================================================================================ %

\newcommand{\indup}{\mathbb{I}_{\uparrow}}
\newcommand{\inddown}{\mathbb{I}_{\downarrow}}
\newcommand{\inds}{\mathbb{I}_{s}}
\newcommand{\indl}{\mathbb{I}_{\ell}}

\newcommand{\aup}{\alpha_{\uparrow}}
\newcommand{\adown}{\alpha_{\downarrow}}


\begin{document}

\title{
Convergence Rate of the (1+1)-ES on Locally Strongly Convex and Lipschitz Smooth Functions
%and their Monotonic Transformations
}%


\author{
Daiki Morinaga,
\thanks{D. Morinaga is with the Department of Computer Science, University of Tsukuba, Tsukuba, Japan; RIKEN AIP, Japan (e-mail: morinaga@bbo.cs.tsukuba.ac.jp).}
\and Kazuto Fukuchi,
\and Jun Sakuma,
\and Youhei Akimoto\thanks{
K. Fukuchi, J. Sakuma, and Y. Akimoto are with the Institute of Systems and Information Engineering, University of Tsukuba, Tsukuba, Japan; RIKEN AIP, Japan
(e-mail: fukuchi@cs.tsukuba.ac.jp; jun@cs.tsukuba.ac.jp; akimoto@cs.tsukuba.ac.jp).}
}


% The paper headers
%\markboth{Journal of \LaTeX\ Class Files,~Vol.~14, No.~8, August~2021}%
%{Shell \MakeLowercase{\textit{et al.}}: A Sample Article Using IEEEtran.cls for IEEE Journals}

%\IEEEpubid{0000--0000/00\$00.00~\copyright~2021 IEEE}
% Remember, if you use this you must call \IEEEpubidadjcol in the second
% column for its text to clear the IEEEpubid mark.

\thispagestyle{empty}
\noindent\textbf{Bibliographic information:}
{\small
\begin{verbatim}
@ARTICLE{morinaga2024ieeetevc,
  author={Morinaga, Daiki and Fukuchi, Kazuto and Sakuma, Jun and Akimoto, Youhei},
  journal={IEEE Transactions on Evolutionary Computation}, 
  title={Convergence Rate of the (1+1)-ES on Locally Strongly Convex and Lipschitz Smooth Functions}, 
  year={2024},
  volume={28},
  number={2},
  pages={501-515},
  doi={10.1109/TEVC.2023.3266955}}
\end{verbatim}
}

\noindent\textbf{Erratum:}
The version of this paper published at IEEE TEVC (2024) contains some technical errors. The errors do not affect the correctness
of the theorems. They are corrected in this version. For clarity, the changes are marked in {\color{blue}
blue}.


\maketitle
\setcounter{page}{1}
\begin{abstract}
Evolution strategy (ES) is one of the promising classes of algorithms for black-box continuous optimization.
Despite its broad successes in applications, 
theoretical analysis on the speed of its convergence is limited on convex quadratic functions and their monotonic transformation.%theoretically how fast it converges to a optima on convex functions is still vague.
In this study, an upper bound and a lower bound of the rate of linear convergence of the (1+1)-ES on locally $L$-strongly convex functions with $U$-Lipschitz continuous gradient are derived as $\exp\left(-\Omega_{d\to\infty}\left(\frac{L}{d\cdot U}\right)\right)$ and $\exp\left(-\frac1d\right)$, respectively.
Notably, any prior knowledge on the mathematical properties of the objective function, such as Lipschitz constant, is not given to the algorithm, whereas the existing analyses of derivative-free optimization algorithms require it.
\end{abstract}

\begin{IEEEkeywords}
Derivative-free optimization, black-box optimization, evolution strategy, linear convergence, convergence rate, convex optimization.
\end{IEEEkeywords}

\sloppy


\section{Introduction}
\label{intro}
\subsection{Background}


Black-box optimization (BBO) is one of the classes of optimization problem settings in which only the query of the objective function value $x \mapsto f(x)$ is available.
As the computation performance is rapidly developing, practical applications of BBO problems, such as simulation-based optimization, emerge in extensive fields, that is, geoscience, biology, topology optimization, and machine learning \cite{dong2019efficient, fujii2018cma, kriest2017calibrating, uhlendorf2012long}.

Evolution strategy (ES) is regarded as a promising class of continuous BBO algorithms,
and particularly, variants of the covariance matrix adaptation evolution strategy (CMA-ES)\cite{hansen2014principled,hansen2003reducing,hansen2001completely} are in particular demonstrated to achieve prominent performance in benchmark problems and applications \cite{dong2019efficient, fujii2018cma, kriest2017calibrating, uhlendorf2012long, varelas2018comparative}.
While the practical success of ES has been widely recognized, its theoretical foundation is not significantly robust.

The oldest variant of ES is the (1+1)-ES, which was first proposed in \cite{rechenberg1973evolution} and sophisticated in \cite{kern2004learning}.
Although the mechanism of the (1+1)-ES is simpler than that of the state-of-the-art CMA-ES, both share the essence of algorithms (e.g., the randomness of the sampling of the candidate solution and the adaptation of internal parameters), which makes mathematical analysis considerably complicated.
Therefore, one of the goals of this study is to develop a mathematical technique to analyze the class of algorithms with those properties through the analysis of the (1+1)-ES.

\emph{Linear convergence} is a class of convergence and approximately refers to the process whereby, given an objective function, the algorithm finds an $\epsilon$-neighborhood of an optimum in {$\Theta\left(\log\left(1/\epsilon\right)\right)$} number of $f$-calls.
%,although a precise definition of the class of convergence of non-deterministic one can differ according to what the analysis is focusing on.
As the previous works show that the sequences generated by the (1+1)-ES behave in a manner that can be regarded as the linear convergence \cite{akimoto2018drift,akimoto2020global,auger2013linear,jagerskupper2003analysis,jagerskupper20061+,jagerskupper2007algorithmic,morinaga2019generalized},
this study focuses on the rate of linear convergence of the (1+1)-ES.

% While qualitative result, linear convergence, on such function is atained, the quantitative convergence analysis of the (1+1)-ES, and BBO algorithm furthermore, on this class of function is still not.
% The current study pursue quantitative analysis of \emph{convergence rate}, that is, investigate how fast the algorithm converges in the view of constants which parameterize mathematical features of the function, such as the Lipschitz constant.

% The mathematical properties of the objective function we focuses on are strong convexity and Lipschitz smoothness.
% The functions with these two properties are often studied broadly in the field of convex optimization\todo{citation}, and regarding the (1+1)-ES, linear convergence on the class of function which includes the strongly convex and Lipschitz smooth functions has been proved\cite{morinaga2019generalized}.

% Last but not least in this section is about derivative-free optimization (DFO).
% DFO is a similar setting to BBO in that any derivative of the objective function is not available, and there exist theoretical studies of DFO algorithms which proves linear convergence\todo{citation}.
% However, as discussed in the next section, DFO settings often assumes that mathematical properties of the objective function are explicitly available anlike the BBO setting.



% \mori{
% Strong convexity and Lipschitz smoothness is important because...
% \begin{itemize}
%     \item many previous theoretical works study it.
%     \item it parameterizes the degree of convexity.
%     \item other classes of convex functions are...
% \end{itemize}
% }

\subsection{Related Work}

% Existing theoretical researchs on convergence of the (1+1)-ES are first presented to see how far the result in this paper goes in the view of the same algorithm.
% Then, researchs on convergence of derivative-free optimization (DFO) algorithms on strongly convex and Lipschitz smooth functions follow, to clarify the contribution of this paper among DFO studies.

\paragraph{Convergence analysis of the (1+1)-ES}

We analyze the (1+1)-ES with success-based step-size adaptation, which slightly generalizes the 1/5-success rule proposed in \cite{kern2004learning}. The same variant is analyzed in  \cite{akimoto2018drift,auger2013linear,morinaga2019generalized,glasmachers2020global}, and a further generalized variant with a covariance matrix adaptation is analyzed in \cite{morinaga2021convergence}. 
A slightly different variant is analyzed in \cite{jagerskupper2003analysis, jagerskupper20061+, jagerskupper2007algorithmic}.

On the spherical functions $x\in\R^d\mapsto g(\norm{x}^2)$, where $g:\R\to\R$ is an arbitrary strictly increasing transformation, Akimoto et al.~\cite{akimoto2018drift} derive an expected number of function evaluations for the (1+1)-ES to find an $\epsilon$-neighborhood of the optimum, namely the \emph{first hitting time} (FHT), which is in $\Theta\left(d\cdot\log(\norm{m_0 - x^*}/\epsilon)\right)$, where $m_0\in\R^d$ is an arbitrary initial point of the (1+1)-ES and $x^*$ is the optimum. Prior to \cite{akimoto2018drift}, J\"{a}gersk\"{u}pper~\cite{jagerskupper2003analysis,jagerskupper2007algorithmic} also analyze a slightly different variant on the spherical functions.
%In \cite{akimoto2018drift}, they propose a novel \emph{potential function} for unbounded continuous domain which estimates optimation progress, and derive the results with a \emph{additive drift analysis} approach.

On a subset of convex quadratic functions that can be written as
$f(x_1,\ldots,x_d) = \xi\sum_{i=1}^{d/2} x_i^2 + \sum_{j=d/2+1}^d x_j^2$, where $1/\xi\to 0$ as $d\to\infty$,
J\"{a}gersk\"{u}pper~\cite{jagerskupper20061+} show that the number of function evaluations to halve the function value is in $\Theta(d \cdot \xi)$ with an overwhelming probability.
%Loosely speaking, this result translates to the convergence rate of $\Theta\left(\exp\left(-\frac{1}{d \cdot \xi}\right)\right)$.

On the general convex quadratic functions $x\mapsto g(x^\mathrm{T} H x)$ where { $H\in\R^{d\times d}$ is symmetric positive-definite and is the Hessian matrix if $g(y)=y/2$}, Morinaga et al.~\cite{morinaga2021convergence} derive an upper bound of a convergence rate of the worst case in $O\left(\exp\left(-L/\Tr(H)\right)\right)$ ($L$ is the smallest eigenvalue of $H$ and $\Tr(\cdot)$ is the trace), and also derive a lower bound of the best case in $\Omega\left(\exp\left(-1/d\right)\right)$. 
The current study is a pure expansion of \cite{morinaga2021convergence} to strongly convex and Lipschitz smooth functions, including the general convex quadratic functions.\footnote{An example of strongly convex and Lipschitz smooth functions except for a convex quadratic function is $f(x) = g(x) + h(x)$, where $g(x)$ is a convex quadratic function with a Hessian matrix whose eigenvalues are bounded in $[L+M, U-M]$ and $h(x)$ is a twice continuously differentiable function with a Hessian matrix whose eigenvalues are bounded in $[-M, M]$.}
Note that loosely speaking, if the order of the necessary number of function evaluations to converge is represented as $\Theta(\psi(N))$ for a function $\psi$ and a parameter of interest $N$, it suggests that the convergence rate is in $\Theta\left(\exp(-1/\psi(N))\right)$.

On the \emph{homogeneous} functions with all the shapes of levelsets being similar, Auger and Hansen~\cite{auger2013linear} show that the class of convergence is in \emph{linear convergence}. 

%Linear convegence approximately refers that given a objective function, the number of function evaluations necessary to halve a residual error is at least constant over an optimization process, while it is under the stochastic nature of the algorithm.

On a class of functions that includes the strongly convex and Lipschitz smooth functions and a portion of non-convex functions, Morinaga and Akimoto~\cite{morinaga2019generalized} derive a constant upper bound of the FHT and prove linear convergence.
%They modify the potential function proposed in~\cite{akimoto2018drift} and also explots the additive drift analysis approach.
This result is extended by Y. Akimoto et al.~\cite{akimoto2020global} to the (1+1)-ES with the adaptation of a bounded covariance matrix.
{The dependency of the FHT or the convergence rate on $d$, $L$, and $U$ is not discussed in \cite{akimoto2020global} and \cite{morinaga2019generalized}}.
    
    Golovins et al.\cite{2019gradientlesgolovins} show that on a $L$-strongly convex and $U$-Lipschitz smooth function $f$, their proposed DFO method GLD-Fast finds a solution $x$ such that $f(x) - f(x^*)\leq \epsilon$ after a number of function evaluations in $O\left(\frac{d\cdot U}{L}\cdot\log^2\left(\frac{U}{L}\right)\cdot\log\left(\frac{R}{\epsilon}\right)\right)$ with high probability.
    
    These theoretical analyses assume that the algorithm can explicitly exploit parameters of the functions, such as the Lipschitz constant $U$, while they guarantee that the class of convergence of their algorithms is in linear convergence and also offer the runtime order with dependencies of the function parameters $L$ and $U$.
    As the knowledge of such mathematical properties of the objective function cannot be expected in every case of practice when the derivative is not available,
    the current study focuses on the genuine black-box optimization algorithm, which assumes the situation that only the query of the function is available.


%\yohe{See \url{https://www.cambridge.org/core/journals/acta-numerica/article/derivativefree-optimization-methods/84479E2B03A9BFFE0F9CD46CF9FCD289} for overview of DFO approaches. I found some linear convergence results there. Search for \emph{log(} in the pdf and you will find the linear convergence results.}

% \mori{ some related works :

%     \begin{itemize}
%         \item "On Averaging the Best Samples in Evolutionary Computation"(PPSN in 2020) : Find that the averaging strategy is better on the view of the regret than the random search on quadratic functions (sphere? check later if needed).
%         \item "Asymptotic Convergence Rates for Averaging Strategies"(GECCO'21) : Extend the above. Convergence rate of a continuous BBO algorithm with $\lambda$-parallel evaluation and $\mu$-averaging strategy, which is asymptotic to $\lambda$, on a function class which includes three times continuously differentiable functions with unique optimum and is wider than quadratic functions. Uniform averaging, equal weights for the $\mu$ best. Remark : $\mu$ depends on the dimension.
%         \item "Global linear convergence of Evolution Strategies with recombination on scaling-invariant functions"(Stochastic Processes and their Applications in 2021) : Linear convergence of a non-elitist and population-besed ES on the scaling-invariant functions and their strictly increasing transformation with a unique global optimum, although whether it converges or diverges is found only numerically, and its rate of convergence cannot be evaluated quantitatively. While our approach considers an optimization in the worst case, their theoreticcal result guarantees linear convergence both in the worst case and the best case.
%         \item Akimoto, Auger, Hansen. SPA 2022. (mu, lambda)-ES on spherical. Linear. without estimation of CR.
%     \end{itemize}

% }

\subsection{Contributions}

We prove the upper and lower bound of the convergence rate of the (1+1)-ES with success-based step-size adaptation on functions that can be expressed as a composite of a strictly increasing function and a locally $L$-strongly convex function with $U$-Lipschitz continuous gradients. 
Informally, the result is stated as follows: For a sufficiently large dimension $d$, the sequence of the candidate solutions, $\{m_t\}_{t \geq 0}$, generated by the (1+1)-ES satisfies with probability one.
\begin{multline}
    \frac{1}{d} 
    \geq - \liminf_{t\to\infty} \frac{1}{t}\log\left(\frac{\norm{m_t - x^\mathrm{opt}}}{\norm{m_0 - x^\mathrm{opt}}}\right) 
    \\
    \geq - \limsup_{t\to\infty} \frac{1}{t}\log\left(\frac{\norm{m_t - x^\mathrm{opt}}}{\norm{m_0 - x^\mathrm{opt}}}\right) 
    \in \Omega\left(\frac{L}{d\cdot U}\right). 
\end{multline}
Namely, the candidate solution converges linearly (i.e., geometrically) toward $x^\mathrm{opt}$ with convergence rate (i.e., the factor of decrease of $\norm{m_t - x^\mathrm{opt}}$ in each step) no smaller than $\exp\left(-\frac{1}{d}\right)$ and no greater than $\exp\left(-\Omega\left(\frac{L}{d\cdot U}\right)\right)$. The upper bound is refined to $\exp\left(-\Omega\left(\frac{L}{\Tr(H)}\right)\right)$ if the objective function is a composite of a strictly increasing function and convex quadratic function $\frac{1}{2} x^\mathrm{T} H x$ with Hessian matrix $H$ whose eigenvalues are bounded in $[L, U]$. To the best of our knowledge, this is the first analysis to explicitly derive the convergence rate of the BBO algorithm on $L$-strongly convex function with $U$-Lipschitz continuous gradients. These results are formally stated in \Cref{thm:lowerconvergencerate,theorem:lower,cor:upper}. We evaluate the tightness of the derived bounds in numerical experiments on convex quadratic problems in \Cref{sec:exp} and we empirically observe for a large $T$
\begin{multline}
    \min\left\{ 0.1 \cdot \frac{1}{d} , 10 \cdot \frac{L}{\Tr(H)} \right\}
    \\
    \gtrapprox - \frac{1}{0.1 \cdot T}\log\left(\frac{\norm{m_{T} - x^\mathrm{opt}}}{\norm{m_{0.9 \cdot T} - x^\mathrm{opt}}}\right) 
    \gtrapprox 0.1 \cdot \frac{L}{\Tr(H)} 
\end{multline}
on three types of $H$ with varying $d \in [1, 10000]$ and varying condition number $\frac{U}{L} \in [1, 10^6]$. 

\section{Formulation}

In this section, we introduce the algorithm to be analyzed and the class of objective functions to be analyzed.
The definitions of the linear convergence and the convergence rate are formally stated.

% \subsection{Notation}
% \yohe{I think they don't need to be introduced here. Just explain them when they first appears.}
% \begin{itemize}
%     \item $\ind{\cdot}$ : the indicator function. 
%     \item $\norm{\cdot}$ : the Euclidean norm 
%     \item $\inner{\cdot}{\cdot}$ : the inner product
%     \item $\mathcal{N}(0, I)$ : the standard normal distribution 
%     \item $\Phi$ : the cumulative density function induced by $\mathcal{N}(0, I)$
% \end{itemize}

% % \begin{wrapfigure}[]{R}{0.45\textwidth}
% % \begin{minipage}{0.45\textwidth}
\begin{algorithm}[H]
\caption{(1+1)-ES with success-based $\sigma$-adaptation}
\label{algo}
\begin{algorithmic}[1]
\State \textbf{input} $m_0 \in \mathbb{R}^d$, $\sigma_0 > 0$, $f: \R^d \to \R$
\State \textbf{parameter} $\aup > 1$, $\adown < 1$
\For {$t = 0,1,\dots$}%\textit{until a stopping criterion is met}}
\State $x_t = m_t + \sigma_t \cdot z_t$ \text{where} $z_t \sim \mathcal{N}(0,  I)$
\If {$f\big(x_t\big) \leq f\big(m_t\big)$}
\State $m_{t+1} \leftarrow x_t$
\State $\sigma_{t+1} \leftarrow \sigma_t \cdot \aup$
\Else
\State $m_{t+1} \leftarrow m_t$
\State $\sigma_{t+1} \leftarrow \sigma_t \cdot \adown$
\EndIf
\EndFor
\end{algorithmic}
\end{algorithm}
% % \end{minipage}
% % \end{wrapfigure}
% \begin{algorithm}[H]\color{blue}
% \caption{(1+1)-ES with success-based $\sigma$-adaptation}
% \label{algo}
% \begin{algorithmic}[1]
% \State \textbf{input} $m_0 \in \mathbb{R}^d$, $\sigma_0 > 0$, $f: \R^d \to \R$
% \State \textbf{parameter} $\aup > 1$, $\adown < 1$
% \For {$t = 0,1,\dots$}%\textit{until a stopping criterion is met}}
% \State $z_t \sim \mathcal{N}(0,  I)$, $x_t = m_t + \sigma_t \cdot z_t$
% \If {$f\big(x_t\big) \leq f\big(m_t\big)$}
% \State $m_{t+1} \leftarrow x_t$, $\sigma_{t+1} \leftarrow \sigma_t \cdot \aup$
% \Else
% \State $m_{t+1} \leftarrow m_t$, $\sigma_{t+1} \leftarrow \sigma_t \cdot \adown$
% \EndIf
% \EndFor
% \end{algorithmic}
% \end{algorithm}

\subsection{Algorithm}


We analyze the (1+1)-ES with success-based step-size adaptation \cite{kern2004learning} presented in \Cref{algo}. 
Given the initial search point, $m_0 \in \R^d$, and the initial step-size, $\sigma_0 > 0$, it repeats the following steps until a stopping criterion is satisfied. At each iteration $t \geq 0$, it samples a candidate solution, $x_t \in \R^d$, by adding to the current search point, $m_t$, a Gaussian noise $\sigma_t \cdot z_t$ with standard deviation $\sigma_t$. Here, $z_t \sim \mathcal{N}(0, I)$ indicates that $z$ is a $d$-dimensional random vector drawn from the standard normal distribution. If the candidate solution has a better or equally well objective function value than the current search point (i.e., $f(x_t) \leq f(m_t)$) the search point is updated by the candidate solution, and the step-size is increased by multiplying with $\aup > 1$. Otherwise, the current search point is taken over, and the step-size is decreased by multiplying with $\adown < 1$. 


Adaptation of $\sigma$ is designed to maintain the probability of sampling a better or equally good candidate solution than the current search point, $\Pr_{z_t}[f(x_t) \leq f(m_t)]$, to the target probability $p_\mathrm{target}$ defined as
\begin{equation}
p_\mathrm{target} := \frac{\log(1/\adown)}{\log(\aup / \adown)} 
\label{eq:p_target-def}
.
\end{equation}
The factors of the step-size change, $\aup$ and $\adown$, are the hyper-parameters of this algorithm but not intended to be tuned for each problem. A typical choice is $\adown = \aup^{-1/4}$, so that $p_\mathrm{target} = 1/5$; hence, this approach is originally proposed as \emph{1/5-success rule} \cite{rechenberg1973evolution,kern2004learning}. The increase factor, $\aup$, is typically set to either $\exp(c)$, $\exp(c/\sqrt{d})$, or $\exp(c/d)$, where $c > 0$ is a constant.

Optimization with \Cref{algo} is formulated mathematically as the sequence of the parameters, $\{m_t\}_{t\geq 0}$ and $\{\sigma_t\}_{t\geq 0}$, and the sequence of the random vectors, $\{z_t\}_{t\geq 0}$.
Subsequently, $\ind{A}$ is the indicator function of an event $A$ (i.e., $\ind{A} = 1$ if $A$ occurs) and otherwise $\ind{A} = 0$.
\begin{definition}[Markov chain of the (1+1)-ES]\label{def:algorithm}
Let $\Theta = \R^{d+1}$ be the state space, and a state of \Cref{algo} at iteration $t$ is defined as $\theta_t = (m_t, \log(\sigma_t))$.
Let $\{z_t\}_{t \geq 0}$ be the sequence of independent and $\mathcal{N}(0, I)$-distributed random vectors.
For a measurable function $f : \R^d \to \R$, let $\theta_0 = (m_0, \log(\sigma_0))$ and $\theta_{t+1} = \theta_{t} + \mathcal{G}(\theta_t, z_t; f)$, where
\begin{equation}
\begin{split}
\mathcal{G}(\theta, z; f) := (\sigma \cdot z, \log(\aup)) \cdot &\ind{f(m + \sigma \cdot z) \leq f(m)}\\
+ (0, \log(\adown)) \cdot &\ind{f(m + \sigma \cdot z) > f(m)} .
\end{split}
\end{equation}
Let $\{\mathcal{F}_t\}_{t \geq 0}$ be the natural filtration of $\{\theta_t\}_{t \geq 0}$.
We write
\begin{equation}
\{(\theta_t, \mathcal{F}_{t})\}_{t\geq 0} = \texttt{ES}(f, \theta_0, \{z_t\}_{t \geq 0}) .
\end{equation}
\end{definition}

Because of the properties of the (1+1)-ES, we can generalize the analysis on a relatively simple convex objective function to a function class including non-convex functions.
First, the (1+1)-ES is invariant against an arbitrary strictly increasing transformation of the objective function because it is comparison-based, indicating that it requires only comparing two objective function values and the function values themselves are not used. 
Second, it is invariant against the parallel translations of the coordinate system of the search space. This allows us to assume that the optimum is located at the origin without loss of generality.
Third, because the (1+1)-ES is an elitist strategy (i.e., $f(m_{t+1}) \leq f(m_t)$ for all $t \geq 0$) the function landscape outside $S(m_0) = \{x \in \R^d : f(x) \leq f(m_0)\}$ does not matter. 
Therefore, we can assume a simple structure outside $S(m_0)$ for convenience.
These properties are formally stated in the following proposition. 
Its proof is provided in \Cref{apdx:prop:invariance}.


\begin{proposition}[Properties of the (1+1)-ES]\label{prop:invariance}
Given $\{z_t\}_{t \geq 0}$ and $\theta_0 \in \Theta$, the following hold.

(1) Let $g:\R \to \R$ be an arbitrary strictly increasing function (i.e., $g(x) < g(y) \Leftrightarrow x < y$). Subsequently, $\tilde{\theta}_t = \theta_t$ for all $t \geq 0$ for $\{(\theta_t, \mathcal{F}_{t})\}_{t\geq 0} = \texttt{ES}(f, \theta_0, \{z_t\}_{t \geq 0})$ and $\{(\tilde{\theta}_t, \tilde{\mathcal{F}}_{t})\}_{t\geq 0} = \texttt{ES}(g \circ f, \theta_0, \{z_t\}_{t \geq 0})$.

(2) Let $T: x \mapsto x - x^*$ for an arbitrary $x^* \in \R$ and define $S_T: (m, \log(\sigma)) \to (T^{-1}(m), \log(\sigma))$. Subsequently, $\tilde{\theta}_t = S_T(\theta_t)$ for all $t \geq 0$,
for $\{(\theta_t, \mathcal{F}_{t})\}_{t\geq 0} = \texttt{ES}(f, \theta_0, \{z_t\}_{t \geq 0})$ and $\{(\tilde{\theta}_t, \tilde{\mathcal{F}}_{t})\}_{t\geq 0} = \texttt{ES}(f \circ T, S_T(\theta_0), \{z_t\}_{t \geq 0})$. 

(3) Let $\tilde{f}:\R^d\to\R$ be a function such that its restriction to $S(m_0) = \{x \in \R^d : f(x) \leq f(m_0)\}$ (i.e., $\tilde{f}\rvert_{S(m_0)}$) is equivalent to $f\rvert_{S(m_0)}$ and $\tilde{f}(x) > f(m_0)$ for all $x \in \R^d\setminus S(m_0)$. Subsequently, $\tilde{\theta}_t = \theta_t$ for all $t \geq 0$ for $\{(\theta_t, \mathcal{F}_{t})\}_{t\geq 0} = \texttt{ES}(f, \theta_0, \{z_t\}_{t \geq 0})$ and $\{(\tilde{\theta}_t, \tilde{\mathcal{F}}_{t})\}_{t\geq 0} = \texttt{ES}(\tilde{f}, \theta_0, \{z_t\}_{t \geq 0})$.
\end{proposition}

\subsection{Convergence Rate}

%In this section, a property of convergence to be analyzed is defined, and an analysis method is then introduced.

A deterministic real-valued sequence $\{X_t\}_{t\geq 0}$ is said to convergence linearly with convergence rate $r \in (0, 1)$ if
$\lim_{t\to\infty} \abs*{X_{t+1} - X^\mathrm{opt}}/\abs*{X_t - X^\mathrm{opt}} = r$. A slightly more applicable definition is $\lim_{t\to\infty} \frac{1}{t}\log (\abs*{X_{t} - X^\mathrm{opt}} /\abs*{X_0 - X^\mathrm{opt}}) = \log(r)$.
For a stochastic real-valued sequence, various generalizations are considered. 
In this study, the definition of the convergence rate is the limit of the slope of the logarithmic distance between the current solution and the optimum averaged over time.
\Cref{def:convergence} formally defines the convergence rate of the (1+1)-ES.
In what follows, $\norm{x}$ denotes the Euclidean norm of $x \in \R^d$. 
\begin{definition}[Convergence rate]\label{def:convergence}
Let $f: \R^d \to \R$ be a measurable function with the unique global optimum at $x^\mathrm{opt}\in\R^d$, and $\{\theta_t\}_{t \geq 0}$ be the sequence of the state vectors defined in \Cref{def:algorithm} with $m_0 \neq x^\mathrm{opt}$.
If there exists a constant $\mathrm{CR}>0$ satisfying
\begin{align}
\mathrm{Pr}\left[\lim_{t\rightarrow \infty} \frac{1}{t} \log \left( \frac{\norm{ m_t - x^\mathrm{opt} }}{\norm{m_0 - x^\mathrm{opt}}} \right) = -\mathrm{CR} \right] = 1
,
\end{align}
then $\exp(-\mathrm{CR})$ is called the convergence rate of the (1+1)-ES on $f$.
The upper convergence rate $\exp(-\mathrm{CR}^\mathrm{upper})$ on $f$ and the lower convergence rate $\exp(-\mathrm{CR}^\mathrm{lower})$ on $f$ are defined as constants satisfying
\begin{align}
&\mathrm{Pr}\left[\limsup_{t \to \infty} \frac{1}{t} \log \left( \frac{\norm{ m_t - x^\mathrm{opt} }}{\norm{m_0 - x^\mathrm{opt}}} \right) = -\mathrm{CR}^\mathrm{upper} \right] = 1
,
\\
&\mathrm{Pr}\left[\liminf_{t \to \infty} \frac{1}{t} \log \left( \frac{\norm{ m_t - x^\mathrm{opt} }}{\norm{m_0 - x^\mathrm{opt}}} \right) = -\mathrm{CR}^\mathrm{lower} \right] = 1
,
\end{align}
respectively.
\end{definition}

Our main mathematical tool to derive bounds for the upper and lower convergence rate of the (1+1)-ES is the strong law of large numbers on martingales \cite{chow1967strong}.
It reduces the analysis of the limit behavior to the analysis of the expected sequence change in one step. 
The following proposition is the central tool for our analysis.
Its proof is provided in \Cref{apdx:prop:convergence_rate}.
\begin{proposition}\label{prop:convergence_rate}
Let $\{\mathcal{F}_t\}_{t \geq 0}$ be a filtration of a $\sigma$-algebra, and $\{X_t\}_{t \geq 0}$ be a Markov chain adapted to $\{\mathcal{F}_t\}_{t \geq 0}$.
Consider the following conditions:
\begin{enumerate}
\item[C1] $\exists B^\mathrm{upper} > 0$ such that $\E[X_{t+1} - X_{t} \mid \mathcal{F}_t] \leq - B^\mathrm{upper}$ for all $t \geq 0$;
\item[C2] $\exists B^\mathrm{lower} > 0$ such that $\E[X_{t+1} - X_{t} \mid \mathcal{F}_t] \geq - B^\mathrm{lower}$ for all $t \geq 0$;
\item[C3] $\sum_{t=1}^{\infty}{\color{blue}\Var[X_{t}]} / t^2 < \infty$.
\end{enumerate}
(I) If C1 and C3 hold, then we have
$\Pr\left[\limsup_{t \to \infty} \frac{1}{t} (X_t - X_0) \leq - B^\mathrm{upper} \right] = 1$.
% \begin{align}
% \Pr\left[\limsup_{t \to \infty} \frac{1}{t} (X_t - X_0) \leq - B^\mathrm{upper} \right] = 1 \label{eq:prop:upper}
% .
% \end{align}
(II) If C2 and C3 hold, then we have 
$\Pr\left[\liminf_{t \to \infty} \frac{1}{t} (X_t - X_0) \geq - B^\mathrm{lower} \right] = 1$.
% \begin{align}
% \Pr\left[\liminf_{t \to \infty} \frac{1}{t} (X_t - X_0) \geq - B^\mathrm{lower} \right] = 1\label{eq:prop:lower}
% .
% \end{align}
\end{proposition}

% \subsection{First hitting time}

% In the current work, an expectation of \emph{first hitting time} (FHT) is also analyzed.
% While the convergence rate defined in \Cref{def:convergence} describes the slope of the values of a sequence $X_t$ in the limit of the runtime $t$ to infinity,
% FHT is a finite runtime necessary for $X_t$ to attain a specific target value.
% \begin{definition}[First hitting time]\label{def:FHT}
%     Let $\{X_t\}_{t \geq 0}$ be a sequence of real value given $X_0\in\R$.
%     Then, \emph{first hitting time} $T_\beta^X$ of $\{X_t\}_{t\geq 0}$ for a target value $\beta\leq X_0$ is defined as
%     \begin{equation}
%         T_\beta^X = \min\{t\geq 0\mid X_t\leq \beta\}
%         .
%     \end{equation}
%  \end{definition}

% The following theorem offers a sufficient condition for an upper bound of an expectation of FHT.
%  \begin{theorem}[Additive Drift Theorem, Theorem 2.3 in \cite{hajek1982hitting}]\label{theo:additive_drift}
%     Let $\{X_t\}_{t\geq 0}$ be a stochastic process over a state space $S\subseteq\R$, adapted a filtration $\{\mathcal{F}_t\}_{t\geq 0}$,
%     and let $\T_\beta^X$ be a first hitting time of $\{X_t\}_{t\geq 0}$ for $\beta$ where $\beta\leq X_0$.
%     If there exists $\eta>0$ and $0<\rho<1$ such that for any $t\geq 0$,
%     \begin{equation}
%         \E[\exp(\eta(X_{t+1} - X_t))\mid \mathcal{F}_t]\leq \rho
%         ,
%     \end{equation}
%     then, for $s\in( 1, \rho^{-1} )$,
%     \begin{equation}
%         \E[T_\beta^X] \leq (\log(s))^{-1}\left(\eta(X_0 - \beta) + \log\left(1+\frac{s-1}{1-s\rho}\right)\right)
%         .
%     \end{equation}
% \end{theorem}

\subsection{Problem}\label{subsec:problem}

We analyze the Markov chain of the (1+1)-ES defined in \Cref{def:algorithm} for the objective function falling into the following class of functions, $\mathcal{P}_{d, L, U}^{\gamma}$.
Generally, locally strongly convex and Lipschitz smooth functions around the global optimum and their strictly increasing transformations are in the focus of our analysis.

\begin{definition}[$\mathcal{P}_{d, L, U}^\gamma$]\label{def:problem}
For a measurable function $f$ with the unique global minimum at the origin, let $\mathcal{T}^{\gamma}(f)$ be the set of all functions, $\{h: \R^d \to \R\}$, that can be expressed as $h(x) = (g \circ f)(x - x^\mathrm{opt})$ for $x \in \{\tilde{x} \in \R^d : h(\tilde{x}) \leq \gamma\}$, where $g: \R\to\R$ is strictly increasing and $x^\mathrm{opt} \in \R^d$ is the unique global minimum of $h$.
Let $\mathcal{S}_{d,L,U}$ be the set of all functions satisfying \Cref{asm:L_U_2diff} and \Cref{asm:tr_cond} below. 
The functions to be analyzed in this study are $\mathcal{P}_{d, L, U}^{\gamma} = \mathcal{T}^{\gamma}(\mathcal{S}_{d, L, U}) = \{h \mid h \in \mathcal{T}^\gamma(f) \text{ for } f \in \mathcal{S}_{d, L, U}\}$.
\end{definition}

Two important properties are introduced, \emph{strong convexity} and \emph{Lipschitz smoothness}, which are often assumed to analyze the convergence rate of gradient-based approaches.
In what follows, $\inner{a}{b} = a^\T b = \sum_{i=1}^{d} a_i b_i$ denotes the inner product of $a, b \in \R^d$. 
\begin{definition}[Strong convexity]
A differentiable function $f:\R^d\to \R$ is $L$-strongly convex over $S\subseteq\R^d$ if for any $x, y \in S$,
\begin{equation}
f(x) + \langle y-x, \nabla f(x) \rangle + \frac{L}{2}\norm{x-y}^2 \leq f(y)
    .
\end{equation}
\end{definition}
\begin{definition}[Lipschitz smoothness]
A differentiable function $f:\R^d\to\R$ is $U$-Lipschitz smooth over $S\subseteq\R^d$ if for any $x, y \in S$,
\begin{equation}
f(x) + \langle y-x, \nabla f(x) \rangle + \frac{U}{2}\norm{x-y}^2 \geq f(y)
    .
\end{equation}
\end{definition}
In light of \Cref{prop:invariance}, the Markov chain of the (1+1)-ES on $h \in \mathcal{T}^{\gamma}(f)$ is identified by that on $f$.
As is formally stated in \Cref{prop:f-norm} below, the analysis of the convergence rate of the (1+1)-ES on $h$ reduces to that on $f$. 
Namely, by analyzing the (1+1)-ES on $f \in \mathcal{S}_{d,L,U}$, for which we assume a relatively simple structure, we can generalize the results to a class of objective functions, $\mathcal{P}_{d, L, U}^{\gamma}$.

The first assumption on $f$ is stated below.
\begin{assumption}\label{asm:L_U_2diff}
A function $f:\R^d \to \R$ is measurable, continuously differentiable, $L$-strongly convex, and $U$-Lipschitz smooth in $\R^d$. 
The unique global minimum of $f$ is located at the origin, $x^* := \argmin_{x\in \R^d} f(x) = 0$, and the minimum value is $f(x^*) = 0$.
\end{assumption}
Here, we remark on the assumption.
If $f$ is $L$-strongly convex over $\R^d$, then there exists a unique critical point, which is the unique global minimum. 
That is, $\norm{\nabla f(x)}=0 \Leftrightarrow x=\argmin_{y \in \R^d} f(y)$. 
The existence of the global optimum $x^*$ at the origin and $f(x^*) = 0$ simplify our analysis but do not restrict the class of functions to be analyzed because of the arbitrarity of $g$ and $x^\mathrm{opt}$ in \Cref{def:problem}.
Moreover, the $L$-strong-convexity implies the Polyak-{\L}ojasiewicz condition \cite{Karimi2016}. Together with $f(x^*) = 0$, this indicates $\sqrt{2 L f(x)} \leq \norm{\nabla f(x)}$.
If $f$ is $U$-Lipschitz smooth over $\R^d$, its gradient $\nabla f$ is $U$-Lipschitz continuous. That is, $\norm{\nabla f(x) - \nabla f(y)}\leq U\cdot \norm{x-y}$ for all $x, y\in \R^d$.

% \Theta \setminus \{ \theta | m = 0 \}
\providecommand{\Thetazero}{\Theta_{0}}
The second assumption on $f$ is a technical one, which is proved later in \Cref{cor:tr_cond} to hold for sufficiently large $d$ under \Cref{asm:L_U_2diff}.
In the following, the cumulative density function induced by the $1$-dimensional standard normal distribution, $\mathcal{N}(0, 1)$, is denoted by $\Phi$. 
{See also \Cref{subsec:First-order expansion} for further description of $\mathcal{Q}_z(\theta)$.}
\begin{assumption}\label{asm:tr_cond}
Define
\begin{equation}
\mathcal{Q}_z(\theta) = \frac{2}{\sigma^2} \left(f(m + \sigma z) - f(m) - \inner{\nabla f(m)}{ \sigma z} \right) \enspace,\label{eq:q-def}
\end{equation}
and
\begin{equation}
    \mathcal{V}_\mathrm{std} = 
    \Var[\mathcal{Q}_z(\theta)]/(\E[\mathcal{Q}_z(\theta)])^2
    \enspace.
\end{equation}
Then, for $\mathcal{N}(0, I)$-distributed random vector $z \in \R^d$, 
\begin{multline}
\underset{\theta\in\Thetazero}{\sup}
\left\{\mathcal{V}_\mathrm{std}  \right\}
\\
< \frac14 \cdot \min \left\{ \Phi\left(\frac{\kappa_{\inf}}{2}\cdot\frac{1}{\sqrt{2\pi}}\right) - \frac12 , 1 - \Phi\left( \frac{\kappa_{\inf}}{2}\cdot\frac{3}{\sqrt{2\pi}}\right)\right\} 
\label{eq:tr_cond}
\end{multline}
holds with $\kappa_{\inf} := 
\underset{\theta\in\Thetazero}{\inf}
\left\{\frac{\E[\mathcal{Q}_z(\theta)]}{\E[\mathcal{Q}_z(\theta)\cdot\ind{z_e\leq 0}]}\right\} \geq 1$, where $\Thetazero = \Theta \setminus \{\theta \mid m = 0\}$ and $z_e = \inner{z}{\nabla f(m)}/\norm{\nabla f(m)}$.
\end{assumption}



Although we actually analyze the (1+1)-ES on objective function $f$ satisfying \Cref{asm:L_U_2diff} and \Cref{asm:tr_cond} (i.e., $f \in \mathcal{S}_{d, L, U}$), we can generalize the results to that on the function in $\mathcal{P}_{d, L, U}^{\gamma}$ because of the invariance properties of the (1+1)-ES summarized in \Cref{prop:invariance}.
The following proposition states that the convergence rate of the (1+1)-ES on $h \in \mathcal{T}^{\gamma}(f)$ is derived from the convergence rate of the (1+1)-ES on the corresponding $f \in \mathcal{S}_{d,L,U}$. Moreover, the convergence rate of the (1+1)-ES is equivalent to the limit of the slope of $\log(f(m_t))$ averaged over $t$.
Its proof is provided in \Cref{apdx:prop:f-norm}.
\begin{proposition}\label{prop:f-norm}
For measurable function $f$ satisfying \Cref{asm:L_U_2diff}, let $h \in \mathcal{T}^{\gamma}(f)$ and $x^\mathrm{opt}$ be the unique global minimum of $h$.
Let
$S: (m, \log(\sigma)) \mapsto (m + x^\mathrm{opt}, \log(\sigma))$.
We define $\{(\tilde{\theta}_t, \tilde{\mathcal{F}}_{t})\}_{t\geq 0} = \texttt{ES}(h, S(\theta_0), \{z_t\}_{t \geq 0})$ and $\{(\theta_t, \mathcal{F}_{t})\}_{t\geq 0} = \texttt{ES}(f, \theta_0, \{z_t\}_{t \geq 0})$.
Subsequently,  
if $\tilde{\theta}_0 \in \Theta\setminus\{\tilde{\theta}\mid \tilde{m}= x^\mathrm{opt}\}$ satisfies $h(\tilde{m}_0) \leq \gamma$,
we obtain, with probability one (the probability is taken for $\{z_t\}_{t \geq 0}$), 
\begin{align}
\limsup_{t \to \infty} \frac{1}{t} \log \left( \frac{\norm{ \tilde{m}_t - x^\mathrm{opt} }}{\norm{\tilde{m}_0 - x^\mathrm{opt}}} \right)
&=
\limsup_{t \to \infty} \frac{1}{2t} \log \left( \frac{f(m_t)}{f(m_0)} \right)
,
\label{eq:limsup}
\\
\liminf_{t \to \infty} \frac{1}{t} \log \left( \frac{\norm{ \tilde{m}_t - x^\mathrm{opt} }}{\norm{\tilde{m}_0 - x^\mathrm{opt}}} \right)
&=
%\liminf_{t \to \infty} \frac{1}{t} \log \left( \frac{f(m_t)}{f(m_0)} \right)
\liminf_{t \to \infty} \frac{1}{t} \log \left( \frac{\norm{m_t} }{\norm{m_0}} \right)
\label{eq:liminf}
.
\end{align}
\end{proposition}




% An analysis of FHT on a function in $\mathcal{P}_{d, L, U}$ is also completed by that on $f$.
% \begin{proposition}\label{prop:FHT_mnorm}
%   Let $h \in \mathcal{P}_{d, L, U}$ and $x^\mathrm{opt}$, as given in \Cref{def:problem}.
%   Let $f$ be a function which satisfies \Cref{asm:L_U_2diff}, and $S: (m, \log(\sigma)) \mapsto (m + x^\mathrm{opt}, \log(\sigma))$.
%   We define $\{(\tilde{\theta}_t, \tilde{\mathcal{F}}_{t})\}_{t\geq 0} = \texttt{ES}(h, \tilde{\theta}_0, \{z_t\}_{t \geq 0})$ and $\{(\theta_t, \mathcal{F}_{t})\}_{t\geq 0} = \texttt{ES}(f, S(\tilde{\theta}_0), \{z_t\}_{t \geq 0})$.
%   Subsequently, the first hitting time of $\{\log\norm{m_t - x^\mathrm{opt}}\}_{t\geq 0}$ for $\beta\leq\log\norm{m_0 - x^\mathrm{opt}}$ is equal or smaller than to that of $\{\log(f(m_t))\}_{t\geq 0}$ for $\beta + \log\left(\frac{L}{2}\right)$.
% \end{proposition}
% \begin{proof}
%     In light of \Cref{prop:f-norm}, it is trivial as $f(x) \geq \frac{L}{2}\cdot\norm{x}$ for any $x\in\R^d$ when $f$ satisfies \Cref{asm:L_U_2diff}.
%     
% \end{proof}

\section{Lemmas}


In this section, several mathematical tools to analyze a one-step behavior of the (1+1)-ES on objective function $f$ satisfying \Cref{asm:L_U_2diff} and \Cref{asm:tr_cond} are presented. 
In the following, we let $\Thetazero = \Theta \setminus \{\theta \mid m = 0\}$. 

\subsection{First-order expansion}\label{subsec:First-order expansion}
Consider $\mathcal{Q}_z(\theta)$ defined in \eqref{eq:q-def}. 
Because $f$ is strongly convex under \Cref{asm:L_U_2diff}, we have $\mathcal{Q}_z(\theta) > 0$. We can consider that $\frac{\sigma^2}{2} \mathcal{Q}_z(\theta)$ is the remainder term of the first-order Taylor expansion of $f$ around $m$. 
If we let $m = m_t$, $\sigma = \sigma_t$ and $z = z_t$, we have $f(x_t) = f(m_t) + \inner{\nabla f(m_t)}{ \sigma_t z_t} + \frac{\sigma_t^2}{2} \mathcal{Q}_{z_t}(\theta_t)$.
This quantity is essential for analyzing the one-step behavior of the algorithm.
The following lemma shows some properties of $\mathcal{Q}_z(\theta)$ under \Cref{asm:L_U_2diff}.
Its proof is provided in \Cref{apdx:lemma:variance}.
% \yohe{\Cref{lemma:variance} seems to be not used for drift proof. Move it later to the convergence rate evaluation part. Keep \eqref{eq:q-def}.}
% \mori{As we use the integrability of $\E[\mathcal{Q}]$ and $\Var[\mathcal{Q}]$, it's not very bad to put this lemma here.}
\begin{lemma}\label{lemma:variance}
Suppose that $f$ satisfies \Cref{asm:L_U_2diff}. 
Let $\mathcal{Q}_z(\theta)$ be defined in \eqref{eq:q-def}. 
Let $z \in \R^d$ be $\mathcal{N}(0, I)$-distributed and $z_e = \inner{z}{\nabla f(m)}/\norm{\nabla f(m)}$. 
Then, for any $\theta\in\Thetazero$,
(1) $d L \leq \E[\mathcal{Q}_z(\theta)] \leq d U$;
(2) $\Var[\mathcal{Q}_z(\theta)] \leq 4 d U^2$;
(3) $\abs{ \E[\mathcal{Q}_z(\theta) \cdot \ind{z_e \leq 0}] - \E[\mathcal{Q}_z(\theta)] / 2 } \leq (2 / d)^{1/2} (U/L) \cdot \E[\mathcal{Q}_z(\theta)]$.
% Moreover, if $f$ is twice continuously differentiable, then for any $\bar{m} \in \R^d$
% \begin{enumerate}
% \item[4.] $\lim_{(m,\sigma) \to (\bar{m}, 0)} \E[\mathcal{Q}_z(\theta)] = \Tr(\nabla^2 f(\bar{m}))$;
% \item[5.] $\lim_{(m,\sigma) \to (\bar{m}, 0)} \Var[\mathcal{Q}_z(\theta)] = 2 \Tr((\nabla^2 f(\bar{m}))^2)$;
% \item[6.] $\lim_{(m,\sigma) \to (\bar{m}, 0)} \E[\mathcal{Q}_z(\theta) \cdot \ind{z_e \leq 0}] = \Tr(\nabla^2 f(\bar{m})) / 2$,
% \end{enumerate}
% where the limits are taken for any sequence of $\theta = (m, \sigma) \in\Theta\setminus\{\theta|m\neq x^*\}$.
\end{lemma}

In light of \Cref{lemma:variance}, we can show that \Cref{asm:tr_cond} is satisfied for a sufficiently large $d$ under \Cref{asm:L_U_2diff}. The proof of the following result is provided in \Cref{apdx:cor:tr_cond}.
\begin{corollary}\label{cor:tr_cond}
Given $U \geq L > 0$, there exists an integer $D$ such that \Cref{asm:tr_cond} is satisfied for all $f$ satisfying \Cref{asm:L_U_2diff} and for all $d \geq D$. 
\end{corollary}



\subsection{Log-progress}

The log-progress of the (1+1)-ES is defined as the single step decrease of the logarithm of the objective function value, that is,
\begin{equation}
\log\left(\frac{f(m_{t+1})}{f(m_t)}\right) = \log\left( \frac{ f(m_t + \sigma_t z_t)}{ f(m_t) } \right) \ind{f(m_t + \sigma_t z_t) \leq f(m_t)}.
\end{equation}
Along with inequality $\log(x) \leq x - 1$ for $x > 0$, the next lemma offers an upper bound of the expected log-progress.
The proof is provided in \Cref{apdx:lemma:qualitygain}
\begin{lemma}[Upper bound of expected log-progress]\label{lemma:qualitygain}
Let $f$ be a function which satisfies \Cref{asm:L_U_2diff}. Let $z \in \R^d$ be $\mathcal{N}(0, I)$-distributed and $z_e = \inner{z}{\nabla f(m)}/\norm{\nabla f(m)}$.
%Let $z_e = (\nabla f(m)^T z)/\norm{\nabla f(m)}$, $\delta\in(0, 1)$, and $T = \E\left[z^\mathrm{T} H_z(\theta) z \cdot \ind{z_e \leq 0}\right]$.
For any $\theta\in\Thetazero$,
\begin{multline}
\E\left[\left( \frac{ f(m + \sigma z) - f(m) }{ f(m) } \right) \ind{f(m + \sigma z) \leq f(m)}\right]
\leq \\
\frac{\sigma \norm{\nabla f(m)}}{f(m)}\cdot\left(
\frac{\sigma \cdot \E\left[\mathcal{Q}_z(\theta)\cdot \ind{z_e\leq 0} \right] }{2\norm{\nabla f(m)}}
- \frac{1}{\sqrt{2\pi}}\right)
\\
\cdot \Pr[f(m + \sigma z) \leq f(m) ]
.
\label{eq:qualitygain}
\end{multline}
%The RHS is negative if and only if $\sigma<\sqrt{\frac{2}{\pi}}\cdot \frac{\norm{\nabla f(m)}}{\E\left[\mathcal{Q}_z(\theta) \right]}$. %where \E\left[\mathcal{Q}_z(\theta) \right] > 0$.
%\mori{Is $T=\E\left[z^\mathrm{T} H_z z \cdot \ind{z_e \leq 0}\right]=\frac12 \E[z^\mathrm{T} H_z z]$?}
\end{lemma}

Using inequality $x \geq 1 - \exp(\abs{x})$ for $x \leq 0$, the following lemma argues that the log-progress is finite in expectation. Moreover, with the fact that $\Var[x] \leq \E[x^2] \leq 2 \E\left[\exp(\abs{x})\right]$, it also offers an upper bound of the variance of the log-progress.
Its proof is provided in \Cref{apdx:lemma:variancebound}.
\begin{lemma}[Variance bound of log-progress]\label{lemma:variancebound}
Suppose $f:\R^d\to\R$ satisfies \Cref{asm:L_U_2diff} and $d>3$. 
Let $z \in \R^d$ be $\mathcal{N}(0, I)$-distributed.
For any $\theta\in\Thetazero$,
\begin{multline}
\E\left[\exp\left( \abs*{\log\left( \frac{ f(m + \sigma z) }{ f(m) } \right) \ind{f(m + \sigma z) \leq f(m)} }\right)\right]
\\
\leq
\frac{U}{L}\cdot\left(1+\frac{1}{d-3}\right)
.
\label{eq:lemma:variancebound}
\end{multline}
\end{lemma}
%        \mori{Trial for variance bound.
%        
%        (Convex quadratic case : $f(x) = \frac12 x^\mathrm{T} H x$)
%        
%            We have $\mathcal{Q}_z(\theta) = z^\mathrm{T} H z$.
%            If $\mathcal{Q}_z(\theta)=0$,
%            $f(m+\sigma z) - f(m) = 0$.
%            If $\mathcal{Q}_z(\theta)\neq 0$,
%            \begin{align}
%                f(m+\sigma z) - f(m)
%                =&
%                \frac{\sigma^2}{2} \mathcal{Q}_z(\theta) + \sigma \norm{\nabla f(m)} z_e
%                \\
%                =&
%                \frac{\mathcal{Q}_z(\theta)}{2}\cdot\left(
%                \left(\sigma + \frac{\norm{\nabla f(m)} z_e}{\mathcal{Q}_z(\theta)}\right)^2
%                - \left(\frac{\norm{\nabla f(m)} z_e}{\mathcal{Q}_z(\theta)}\right)^2
%                \right)
%                \\
%                \geq&
%                - \frac{\norm{\nabla f(m)}^2 z_e^2}{2\mathcal{Q}_z(\theta)}
%                .
%            \end{align}
%            Let $\{\lambda_i\}_{i=0, 1, \ldots, d}$ be the eigenvalues of $H$ and $\lambda_i\geq \lambda_j$ for any $i, j = 0, 1, \ldots, d$.
%            hen,
%            \begin{align}
%                &\exp\left(\abs*{\log\left(\frac{f(m+\sigma z)}{f(m)}\right)\cdot\ind{f(m+\sigma z)\leq f(m)}}\right)
%                \\
%                \leq&
%                \exp\left(
%                -
%                \log
%                \left(1
%                - \frac{\norm{\nabla f(m)}^2 z_e^2}{2f(m) \mathcal{Q}_z(\theta)}
%                \right)
%                \right)
%                \\
%                =&
%                \left(1
%                - \frac{\norm{\nabla f(m)}^2 z_e^2}{2f(m) \mathcal{Q}_z(\theta)}
%                \right)^{-1}
%                \\
%                =&
%                \mathcal{Q}_z(\theta)
%                \left(
%                \mathcal{Q}_z(\theta)
%                - \frac{\norm{\nabla f(m)}^2 z_e^2}{2f(m)}
%                \right)^{-1}
%                \\
%                \leq&
%                \mathcal{Q}_z(\theta)
%                \left(
%                \mathcal{Q}_z(\theta)
%                - \lambda_1 z_e^2
%                \right)^{-1}
%                (\because \norm{\nabla f(m)} \leq \sqrt{2 \lambda_1 f(m)} )
%                \\
%                =&
%                z^\mathca{T} H z
%                \left(
%                z^\mathca{T} H z
%                - \lambda_1 z_e^2
%                \right)^{-1}
%                 \\
%                =&
%                1
%                + \lambda_1 z_e^2
%                \cdot
%                \left(
%                z^\mathca{T} H z
%                - \lambda_1 z_e^2
%                \right)^{-1}               
%                .
%            \end{align}
%            The 2nd term of the RMS is stochastically dominated by
%            \begin{equation}
%                \frac{
%                \lambda_1 \mathcal{N}_1^2
%                }{
%                \lambda_d \sum_{i=2}^d \mathcal{N}_i^2
%                }
%                ,
%            \end{equation}
%            where $\{\mathcal{N}_i\}_{i=0, 1, \ldots, d}$ are i.i.d random variables which follow the standard normal distribution.
%            Moreover, $
%                \frac{
%                \mathcal{N}_1^2
%                }{
%                \sum_{i=2}^d \mathcal{N}_i^2
%                }
%            $ follows F-distribution.
%        
%        
%        (Trial 1)
%            Let $G_\mathrm{gain}(\sigma) := f(m+\sigma z) - f(m)$.
%            Then, 
%            \begin{equation}
%                &\frac{d G_\mathrm{gain}(\sigma)}{d\sigma}
%                =
%                \frac{d f(m+\sigma z)}{d\sigma}
%                =
%                \langle \nabla f(m+\sigma z), z\rangle
%                .
%            \end{equatio}      
%            Because
%            \begin{equation}
%                \langle \nabla f(m+\sigma z), z\rangle
%       .         =
%                \langle \nabla f(m+\sigma z) - \nabla f(m),z\rangle
%                + \langle \nabla f(m), z\rangle
%                ,
%            \end{equation}
%            and as $f$ is $L$-strongly convex and $U$-Lipschitz smooth,
%            \begin{equation}
%                \sigma^2\cdot L \norm{z}^2
%                \leq
%                \langle \nabla f(m+\sigma z) - \nabla f(m), \sigma z\rangle
%                \leq
%                \sigma^2\cdot U \norm{z}^2
%                ,
%            \end{equation}
%            then, $\frac{d G_\mathrm{gain}(\sigma)}{d\sigma}$ is positive for any $\sigma>0$ if $\langle \nabla f(m), z\rangle > 0$, which leads that $G_\mathrm{gain}(\sigma)$ is strictly increasing for $\sigma>0$.
%            If $\langle \nabla f(m), \sigma z\rangle \leq 0$, $\frac{d G_\mathrm{gain}(\sigma')}{d\sigma}$ can be zero or negative.
%            Further, since $\frac{d G_\mathrm{gain}(\sigma)}{d\sigma}\to \infty$ as $\sigma\to\infty$, $\argmin_\sigma\{G_\mathrm{gain}(\sigma)\} = \sigma^*$ where $\sigma^*\in\left\{\sigma\mid\frac{d G_\mathrm{gain}(\sigma)}{d\sigma} = 0\right\}$.
%            %Therefore, as $\frac{d G_\mathrm{gain}(\sigma)}{d\sigma} = \langle \nabla f(m+\sigma z), z\rangle$,
%            %\begin{align}
%            %    \min_\sigma\{G_\mathrm{gain}(\sigma)\}
%            %    =&
%            %    \min_\sigma\left\{
%            %    f(m+\sigma z) - f(m)
%            %    - \langle \nabla f(m+\sigma z), z\rangle
%            %    \right\}
%            %    .
%            %\end{align}
%            
%            (Trial 2)
%            In light of mean-value theorem,
%            there exists $\delta\in[0, 1]$ such that
%            \begin{equation}
%                f(m+\sigma z) - f(m) = \langle \nabla f(m+\delta\sigma z), \sigma z\rangle
%                .
%            \end{equation}
%            Assumen. $\delta \neq 0$.
%            As $f$ is $L$-strongly convex,
%            \begin{align}
%                \langle \nabla f(m+\delta\sigma z), \sigma z\rangle
%                =&
%                \frac{1}{\delta} \cdot \left(
%                \langle \nabla f(m+\delta\sigma z) - \nabla f(m), \delta \sigma z\rangle + \langle \nabla f(m), \delta \sigma z\rangle
%                \right)
%                \\
%                \geq&
%                \frac{1}{\delta} \cdot \left(
%                L\delta^2\sigma^2\norm{z}^2 + \langle \nabla f(m), \delta \sigma z\rangle
%                \right)
%                \\
%                =&
%                L\delta\sigma^2\norm{z}^2 + \langle \nabla f(m), \sigma z\rangle
%                \\
%                =&
%                L\delta\left(
%                \norm{\sigma z + \frac{1}{2L\delta}\nabla f(m)}^2 - \frac{1}{4L^2\delta^2}\norm{\nabla f(m)}^2
%                \right)
%                \\
%                \geq&
%                - \frac{1}{4L\delta}\norm{\nabla f(m)}^2
%                .
%            \end{align}
%            Note : $\delta$ is a random variable.
%        }

\subsection{Success probability}

The success probability is the probability of sampling a candidate solution $x_t$ with $f(x_t) \leq f(m_t)$, that is, 
\begin{equation}
\Pr[f(x_t) \leq f(m_t)] = \Pr[f(m_t + \sigma_t \cdot z_t) \leq f(m_t)].
\end{equation}
Analyzing the success probability is important to evaluate the upper bound of the expected log-progress provided in \Cref{lemma:qualitygain}.

The following lemma derives an upper bound and a lower bound of the success probability.
The proof is provided in \Cref{apdx:lemma:successprobability}.
\begin{lemma}[Bounds of success probability]\label{lemma:successprobability}
Define the normalized step-size: 
\begin{equation}
\bar\sigma = \sigma\cdot\E[\mathcal{Q}_z(\theta)]/\norm{\nabla f(m)}
\enspace.
\end{equation}
Let $f$ be a function that satisfies \Cref{asm:L_U_2diff}.
Let $z \in \R^d$ be $\mathcal{N}(0, I)$-distributed.
Then, for any $\theta\in\Thetazero$, and $\epsilon > 0$,
\begin{multline}
\Phi\left( - \frac12 \bar\sigma\cdot(1+\epsilon)\right)
  - \frac{1}{\epsilon^2} \cdot \mathcal{V}_\mathrm{std}
\leq 
\Pr\left[f(m + \sigma z) \leq f(m) \right]
\\
\leq 
\Phi\left( - \frac12 \bar\sigma\cdot(1-\epsilon)\right)
  + \frac{1}{\epsilon^2} \cdot \mathcal{V}_\mathrm{std}
.
\label{eq:lemma:successprobability}
\end{multline}
\end{lemma}

The main message of \Cref{lemma:successprobability} is that the success probability is approximated by $ \Phi\left( - \frac12 \bar\sigma\right)$ if $\mathcal{V}_\mathrm{std} \ll 1$.
Later in \eqref{eq:tr_cond_asymptotic}, $\mathcal{V}_\mathrm{std} \ll 1$ is proved for the limit $d \to \infty$ for functions of interest.
Therefore, we can bound the success probability by using the normalized step-size $\bar\sigma$. 
The following corollary provides sufficient conditions on the normalized step-size to bound the success probability from above and below.
Its proof is provided in \Cref{apdx:cor:successprobability}.
\begin{corollary}%[Normalized $\sigma$ for a constant bound of success probability]
\label{cor:successprobability}
Let $f$ be a function which satisfies \Cref{asm:L_U_2diff}.
Let $z \in \R^d$ be $\mathcal{N}(0, I)$-distributed.
Suppose that $\theta \in \Thetazero$.

Define $B_\theta^\mathrm{high} : \left(0, \frac12\right) \to \R_{> 0}$ as
\begin{equation}
B_\theta^\mathrm{high}(q) := \sup_{
{\epsilon\in\left(\sqrt{ \frac{ 2 }{1 - 2q} \cdot \mathcal{V}_\mathrm{std} }, \infty\right)}
}
\frac{ 2 \Phi^{-1} \left(1 - \left(q + \frac{1}{\epsilon^2} \cdot \mathcal{V}_\mathrm{std}\right) \right)}{ 1 + \epsilon} .
\label{eq:cor:b_high}
\end{equation}
Then, 
$B_\theta^\mathrm{high}(q)$ is positive for any $q\in(0, \frac12)$ and $\theta\in\Thetazero$, 
and $B_\theta^\mathrm{high}(q) \leq 2 \Phi^{-1}(1-q)$ for all $q \in \left(0, \frac12\right)$.
Moreover, 
\begin{equation}
\bar\sigma
\leq B_\theta^\mathrm{high}(q)
\Rightarrow
\Pr\left[f(m + \sigma z) \leq f(m) \right] > q
.
\label{eq:cor:successprobability:large}
\end{equation}

Suppose that $\mathcal{V}_\mathrm{std}<1/2$.%$\Var[\mathcal{Q}_z(\theta)] < \E[\mathcal{Q}_z(\theta)]^2 / 2$.
Define $B_\theta^\mathrm{low}: \left( \mathcal{V}_\mathrm{std} , \frac{1}{2}\right) \to \R_{> 0}$ as
\begin{equation}
B_\theta^\mathrm{low}(q) := \inf_{
{\epsilon\in\left(\sqrt{ \frac{ 1 }{q} \cdot \mathcal{V}_\mathrm{std} }, 1\right)}
}
\frac{ 2 \Phi^{-1} \left(1 - \left(q - \frac{1}{\epsilon^2} \cdot \mathcal{V}_\mathrm{std}\right) \right)}{ 1 - \epsilon} .
\label{eq:cor:b_low}
\end{equation}
Then,
$B_\theta^\mathrm{low}(q)$ is positive for any $q\in\left( \mathcal{V}_\mathrm{std} , \frac{1}{2}\right)$ and $\theta\in\Thetazero$,
and $B_\theta^\mathrm{low}(q) \geq 2 \Phi^{-1}(1-q)$ for all $q \in \left( \mathcal{V}_\mathrm{std} , \frac{1}{2}\right)$.
Moreover,
\begin{equation}
\bar\sigma
\geq B_\theta^\mathrm{low}(q)
\Rightarrow
\Pr\left[f(m + \sigma z) \leq f(m) \right] < q
.
\label{eq:cor:successprobability:small}
\end{equation}
\end{corollary}

Finally, we arrange a tool to bound the success probability by some constants. The state space $\Thetazero$ is divided into three non-exclusive subsets based on the normalized step-size $\bar\sigma$.
The proof is provided in \Cref{apdx:lemma:cases}.
\begin{lemma}[Three cases of $\sigma$]\label{lemma:cases}
Let $f$ be a function which satisfies \Cref{asm:L_U_2diff}.
Let $z \in \R^d$ be $\mathcal{N}(0, I)$-distributed.
Let $B_\theta^\mathrm{high}$ and $B_\theta^\mathrm{low}$ be defined in \Cref{cor:successprobability}.
Define $B_{\inf}^\mathrm{high}(q) = \underset{\theta\in\Thetazero}{\inf} \left\{ B_\theta^\mathrm{high}(q) \right\}$ and $B_{\sup}^\mathrm{low} = \underset{\theta\in\Thetazero}{\sup} \left\{ B_\theta^\mathrm{low}(q) \right\}$.
Then, the following statements hold.
%   Subsequently, there exists the following non-empty open interval $I_q\subset\R$
%   \begin{multline}
%     I_q := \left(
%     \max\left\{
%     \underset{\theta\in\Theta\setminus\{\theta|m\neq x^*\}}{\sup}
%     \left\{\mathcal{V}_\mathrm{std}\right\}
%     ,
%     \right.
%     \right.
%     \\
%     \left.
%     \left.
%     \inf\left\{
%     q\in\left(0, \frac12\right):
%     \kappa_{\inf}
%     \cdot
%     \sqrt{\frac{2}{\pi}} > 
%     \underset{\theta\in\Theta\setminus\{\theta|m\neq x^*\}}{\sup}
%     \left\{
%     B_\theta^\mathrm{low}(q) 
%     \right\}
%     \right\}
%     \right\}
%     , \frac{1}{2} 
%     \right)
%     .\label{eq:def_Iq}
%   \end{multline}
%   Given a constant $q^\mathrm{low}\in I_q$,
%   there exist two constants
%   \begin{equation}
%     q^\mathrm{high} := 
%     \underset{\theta\in\Theta\setminus\{\theta|m\neq x^*\}}{\sup}
%     \left\{
%     \inf
%     \left\{
%     q \in \left(q^\mathrm{low}, \frac12 \right) :
%     B_\theta^\mathrm{low}(q^\mathrm{low}) 
%     > \frac{\aup}{\adown}\cdot B_\theta^\mathrm{high}(q)
%     \right\}
%     \right\}
%     ,
%     \label{eq:def_qhigh}
%   \end{equation}
%   and
%     \begin{align}
%     Q &:= 
%     \sup \left\{ \bar{Q} : \forall q \in (0, \bar{Q}], 
%     %Q_\theta : 
%     \underset{\theta\in\Theta\setminus\{\theta|m\neq x^*\}}{\inf}
%     \left\{
%     B_\theta^\mathrm{high}(q)
%     %B_\theta^\mathrm{high}(Q_\theta) 
%     \right\}
%     \geq 
%     \underset{\theta\in\Theta\setminus\{\theta|m\neq x^*\}}{\sup}
%     \left\{
%     B_\theta^\mathrm{low}(q^\mathrm{low}) \right\}
%     \right\}
%     > 0,
%     \label{eq:def_Q}
%   \end{align}
%   %
%Then, the following statements hold.
%\begin{enumerate}
%\renewcommand{\labelenumi}{(\roman{enumi})}

(i) The open interval $I_q\subseteq (0, 1/2)$ defined as 
\begin{multline}
I_q := \left(
\max\left\{
\underset{\theta\in\Thetazero}{\sup}
\left\{\mathcal{V}_\mathrm{std}\right\}
,
\right.
\right.
\\
\left.
\left.
\sup\left\{
q\in\left(0, \frac12\right):
\kappa_{\inf}
\cdot
\sqrt{\frac{2}{\pi}} \leq 
B_{\sup}^\mathrm{low}(q) 
\right\}
\right\}
, \frac{1}{2} 
\right)
\label{eq:def_Iq}
\end{multline}
is nonempty.

(ii) For any $q^\mathrm{low}\in I_q$, the open interval $I_q^{\mathrm{high}}(q^\mathrm{low}) \subseteq (0, 1/2)$ defined as
\begin{multline}
I_q^{\mathrm{high}}(q^\mathrm{low})
\\
:=
\left(
\underset{\theta\in\Thetazero}{\sup}
\left\{
\sup
\left\{
q :
B_\theta^\mathrm{low}(q^\mathrm{low}) 
\leq \frac{\aup}{\adown}\cdot B_\theta^\mathrm{high}(q)
\right\}
\right\}, \frac12 \right)
\label{eq:def_qhigh}
\end{multline}
is nonempty.
% {\color{red}TBD: Old version below:
% For any $q^\mathrm{low}\in I_q$, 
%   \begin{equation}
%     q^\mathrm{high} := 
%     \underset{\theta\in\Theta\setminus\{\theta|m\neq x^*\}}{\sup}
%     \left\{
%     \inf
%     \left\{
%     q \in \left(q^\mathrm{low}, \frac12 \right) :
%     B_\theta^\mathrm{low}(q^\mathrm{low}) 
%     > \frac{\aup}{\adown}\cdot B_\theta^\mathrm{high}(q)
%     \right\}
%     \right\}
%     \label{eq:def_qhigh}
%   \end{equation}
% exists in $\left(q^\mathrm{low}, \frac12 \right)$.}

(iii)
\begin{equation}
Q := 
\inf \left\{ q : 
%Q_\theta : 
B_{\inf}^\mathrm{high}(q)
%B_\theta^\mathrm{high}(Q_\theta) 
< 
B_{\sup}^\mathrm{low}(q^\mathrm{low})
\right\}
>0
\label{eq:def_Q}
\end{equation}
exists
for any $q^\mathrm{low}\in I_q$.

%\end{enumerate}
Let $q^\mathrm{low} \in I_q$ and $q^\mathrm{high} \in I_q^\mathrm{high}(q^\mathrm{low})$. Then, the following statements hold.
%
\begin{align}
  \text{(iv) }
  &\Pr\left[f(m + \sigma z) \leq f(m) \right] \geq q^\mathrm{high} 
  \ \text{if}\ \bar\sigma < B_{\inf}^\mathrm{high}(q^\mathrm{high}). \label{eq:lemma:toosmall}  
  \\
  \text{(v) }
  &\Pr\left[f(m + \sigma z) \leq f(m) \right] \leq q^\mathrm{low} 
  \ \text{if}\  \bar\sigma > B_{\sup}^\mathrm{low}(q^\mathrm{low}). \label{eq:lemma:toolarge}
  \\
  \text{(vi) }
  &\Pr\left[f(m + \sigma z) \leq f(m) \right] \geq Q 
  \ \text{if}\ \bar\sigma \leq B_{\sup}^\mathrm{low}(q^\mathrm{low}). \label{eq:lemma:reasonable}
\end{align}

% % \begin{enumerate}
% % \renewcommand{\labelenumi}{(\roman{enumi})}
% % \setcounter{enumi}{3}
% %$\sigma < B_{\inf}^\mathrm{high}(q^\mathrm{high}) \cdot \norm{ \nabla f(m) } / \E[\mathcal{Q}_z(\theta)]$,
% \begin{equation}
% \text{(iv) If }
% \bar\sigma < B_{\inf}^\mathrm{high}(q^\mathrm{high}),
% \Pr\left[f(m + \sigma z) \leq f(m) \right] \geq q^\mathrm{high} . \label{eq:lemma:toosmall}
% \end{equation}
% %%$\sigma < B_{\inf}^\mathrm{high}(q^\mathrm{high}) \cdot \norm{ \nabla f(m) } / \E[\mathcal{Q}_z(\theta)]$,

% %$\sigma > B_{\sup}^\mathrm{low}(q^\mathrm{low})\cdot \norm{\nabla f(m)} / \E[\mathcal{Q}_z(\theta)]$,
% \begin{equation}
% \begin{aligned}
% \MoveEqLeft[4]
% \text{(v) If }
% \bar\sigma > B_{\sup}^\mathrm{low}(q^\mathrm{low}),
% \Pr\left[f(m + \sigma z) \leq f(m) \right] \leq q^\mathrm{low} . \label{eq:lemma:toolarge}
% \end{aligned}
% \end{equation}

% %$\sigma \leq B_{\sup}^\mathrm{low}(q^\mathrm{low}) \cdot \norm{\nabla f(m)} / \E[\mathcal{Q}_z(\theta)]$,
% \begin{equation}
% \begin{aligned}
% \MoveEqLeft[5]
% \text{(vi) If }
% \bar\sigma \leq B_{\sup}^\mathrm{low}(q^\mathrm{low}), 
% \Pr\left[f(m + \sigma z) \leq f(m) \right] \geq Q . \label{eq:lemma:reasonable}
% \end{aligned}
% \end{equation}
% Moreover, 
% \begin{equation}
%     \E\left[\log\left( \frac{ f(m + \sigma z) }{ f(m) } \right) \ind{f(m + \sigma z) \leq f(m)}\right]
%     \leq w
%   \end{equation}
%   with $E_\mathcal{Q} = \underset{\theta\in\Theta\setminus\{\theta|m\neq x^*\}}{\sup} \left\{\E[\mathcal{Q}_z(\theta)]\right\}$ and
% \begin{equation}
%     w = 
%     \frac{L}{E_\mathcal{Q}}
%     \cdot
%     \frac{
%     B_{\inf}^\mathrm{high}(q^\mathrm{high})
%     }{\kappa_{\inf}}
%     \cdot
%     \left(
%     \sqrt{\frac{2}{\pi}}\cdot\kappa_{\inf}
%     -
%      B_{\sup}^\mathrm{low}(q^\mathrm{low})
%     \right)
%     \cdot Q
%     >0
%   .
%   \label{eq:psi}
% \end{equation}
%     \todo{I feel that this is not what we want.} \mori{I have a feeling that it is clearer to argue the three lemmas completely separately, that is, the lemma about the success probability, that about the potential function, and that about the upper drift. That is why I didn't put the log-progress bound in this lemma. What do you think?}\yohe{It is okay to split the lemma. But this is not what I meant. My question is whether these cases corresponds to the cases we split to show the potential decrease? It is not too trivial to see if (iv--vi) corresponds to the cases in Corollary 2.}
%   \mori{
%   I see.
%   Let me explain in detail.
%   The cases in this lemma does not correspond to the cases to show the potential decrease in \Cref{lemma:potential_decrease}, but is a subset of the cases to show the expected potential decrease in \Cref{cor:expected_potential_decrease} as follows.

%   \begin{itemize}
%       \item 
%   The reasonable range for success probability we define in \Cref{lemma:cases} : 
%   \begin{equation}
%   \frac{B_{\inf}^\mathrm{high}(q^\mathrm{high}) \cdot \norm{ \nabla f(m) }}{ \E[\mathcal{Q}_z(\theta)]}
%   \leq \sigma \leq  
%   \frac{B_{\sup}^\mathrm{low}(q^\mathrm{low})\cdot \norm{\nabla f(m)}}{ \E[\mathcal{Q}_z(\theta)]}
%   \end{equation}

%   \item
%   The reasonable range for potential decrease we define in \Cref{lemma:potential_decrease} (while not using the constants $s$ and $\ell$) : 
%   \begin{equation}
%   \frac{B_{\inf}^\mathrm{high}(q^\mathrm{high}) \cdot \sqrt{ 2 L f(m) }}{ E_\mathcal{Q}} 
%   \leq \sigma \leq  
%   \frac{B_{\sup}^\mathrm{low}(q^\mathrm{low})\cdot \sqrt{2 L f(m)}}{E_\mathcal{Q}}
%   \label{eq:range_potential_decrease}
%   \end{equation}

%   \item
%   The reasonable range for expected potential decrease in \Cref{cor:expected_potential_decrease} :
%     \begin{equation}
%   \frac{B_{\inf}^\mathrm{high}(q^\mathrm{high}) \cdot \sqrt{ 2 L f(m) }}{ E_\mathcal{Q}}
%   \leq \sigma \leq  
%   \frac{B_{\sup}^\mathrm{low}(q^\mathrm{low})\cdot \norm{\nabla f(m)}}{ \E[\mathcal{Q}_z(\theta)]}
%   \label{eq:range_expected_potential_decrease}
%   \end{equation}
%   \end{itemize}

%   Note that $E_\mathcal{Q}\geq \E[\mathcal{Q}_z(\theta)]$ and $\norm{\nabla f(m)}\geq \sqrt{2L f(m)}$ .
%   Therefore, the range in \Cref{cor:expected_potential_decrease} is a superset of both of the ranges in \Cref{lemma:cases} and \Cref{lemma:potential_decrease}.
%   Then, we exploit the following facts to derive the upper drift in \Cref{cor:expected_potential_decrease}.
%   \begin{enumerate}
%       \item The upper bounds of the $V(\theta_{t+1}) -V(\theta_t)$ we derive in \Cref{lemma:potential_decrease} outside \eqref{eq:range_potential_decrease} (too small/large $\sigma$ case), are still the upper bounds of that outside \eqref{eq:range_expected_potential_decrease}.
%       \item The success probability is bounded by $q^\mathrm{low}$ or $q^\mathrm{high}$ outside \eqref{eq:range_expected_potential_decrease}.
%       \item The success probability is bounded from below by $Q$ inside \eqref{eq:range_expected_potential_decrease}.
%       \item Let us see that the upper bound of the potential decrease \eqref{eq:general_potential_decrease}, which is used to derive $E[V(\theta_{t+1}) - V(\theta_t)\mid\mathcal{F}_t]\leq -w/4$ if \eqref{eq:range_expected_potential_decrease} holds (\Cref{cor:expected_potential_decrease}), is a general bound, that is, it holds for any $\sigma>0$ . Then, we take a $\sup$ of \eqref{eq:qualitygain} over \eqref{eq:range_expected_potential_decrease}, and define it as $-w$ (so I didn't put $w$ in this lemma).
%   \end{enumerate}

%   I will clarify the range \eqref{eq:range_expected_potential_decrease} in \Cref{cor:expected_potential_decrease} to get the reader clearly to see these idea.
%   }
% \end{enumerate}
\end{lemma}
% The condition on $f$ \eqref{eq:tr_cond} is prepared for the existence of $I_q, I_q^\mathrm{high}(q^\mathrm{low})$, and $Q$ defined in (i), (ii), and (iii) respectively.
% Subsequently, $q^\mathrm{low}, q^\mathrm{high}$ is prepared to control the success probability under or above $p_\mathrm{target}$, which is formally stated in \Cref{subsec:potentialfunction}.

% \yohe{talk about the relationship of three Bs.}

% \begin{table}[t]
%     \centering
%     \caption{Caption}
%     \label{tab:my_label}
%     \begin{tabular}{c|ccc}
%     \toprule
%         $\bar{\sigma}$ & 
%         $< B_{\mathrm{inf}}^{\mathrm{high}}(q^{\mathrm{high}})$ &
%         $\leq B_{\mathrm{sup}}^{\mathrm{low}}(q^{\mathrm{low}})$ &
%         $>B_{\mathrm{sup}}^{\mathrm{low}}(q^{\mathrm{low}})$ \\
%     \midrule 
%         $\Pr[f(m+\sigma\cdot z) \leq f(m)]$ &
%         $\geq q^\mathrm{high}$ &
%         $\geq Q$ &
%         $\leq q^\mathrm{low}$ \\
%     \bottomrule
%     \end{tabular}
%     \begin{tabular}{c|ccc}
%     \toprule
%         $\sigma$ & 
%         $< \frac{s \cdot \sqrt{ L f(m_t) }}{ \aup\cdot E_\mathcal{Q}}$ &
%          &
%         $>\frac{\ell \cdot \norm{\nabla f(m_t)} }{ \sqrt{2}\adown\cdot E_{z_t}[\mathcal{Q}_{z_t}(\theta_t)]}$ \\
%     \midrule 
%         $\Pr[f(m+\sigma\cdot z) \leq f(m)]$ &
%         $\geq q^\mathrm{high}$ &
%         $\geq Q$ &
%         $\leq q^\mathrm{low}$ \\
%     \bottomrule
%     \end{tabular}
% \end{table}


\section{Convergence Rate}

Our main results, an upper bound of the upper convergence rate and a lower bound of the lower convergence rate of the (1+1)-ES (\Cref{def:algorithm}) on objective function $h \in \mathcal{P}_{d, L, U}^{\gamma}$ (\Cref{def:problem}), are proved in this section. 
The proofs rely on the strong law of large numbers for martingale summarized in \Cref{prop:convergence_rate}.
The analysis of the convergence rate comes down to the analysis of a single-step behavior of the (1+1)-ES.
In light of \Cref{prop:f-norm}, it is sufficient to consider objective function $f$ satisfying \Cref{asm:L_U_2diff} and \Cref{asm:tr_cond} (i.e., $f \in \mathcal{S}_{d, L, U}$).

%    \mori{
%    Memo for convergence rate on smooth strongly convex function :
%    
%    We are bothered with the term about 
%    \begin{equation}
%       \left|\log\left(\frac{\E[\mathcal{Q}_z(\theta_{t})]}{\E[\mathcal{Q}_z(\theta_{t+1})]}\right)\right|
%        .
%    \end{equation}
%    In our convergence rate theorem used in GECCO'21, the following term emerges
%    \begin{equation}
%        \frac{1}{t}\sum_{i=0}^t
%        \left|\log\left(\frac{\E[\mathcal{Q}_z(\theta_{t})]}{\E[\mathcal{Q}_z(\theta_{t+1})]}\right)\right|
%        .
%    \end{equation}
%    We show this term converges to zero as $t\to\infty$ if $f$ is twice continuously differentiable.
%    Let $\bar m_\epsilon\in\R^d$ be a point in $\epsilon$-neighbor of $x^*$, that is, $\norm{\bar m_\epsilon - x^*}\leq \epsilon$ with $\epsilon>0$.
%    \todo{Prove : $(m_t, \sigma_t)\to(\bar m_\epsilon, 0)$ as $t\to\infty$.}
%    In light of \Cref{lemma:variance}, for any $\epsilon'>0$, there exists $T\in\mathbb{N}^+$ such that
%    \begin{equation}
%        |\E[\mathcal{Q}_z(\theta_{t})]-\Tr(\nabla f(\bar m_\epsilon)|\leq \epsilon'
%        \text{ for all } t>T
%        .
%    \end{equation}
%    Then, for $t>T$,
%    \begin{align}
%        &\frac{1}{t}\sum_{i=0}^t
%        \left|\log\left(\frac{\E[\mathcal{Q}_z(\theta_{t})]}{\E[\mathcal{Q}_z(\theta_{t+1})]}\right)\right|
%        \\
%        \leq&
%        \frac{1}{t}
%        \left(
%        \sum_{i=0}^T
%        \left|\log\left(\frac{\E[\mathcal{Q}_z(\theta_{t})]}{\E[\mathcal{Q}_z(\theta_{t+1})]}\right)\right|
%        +
%        (t+1-T)\cdot
%        \log\left(\frac{\Tr(\nabla f(\bar m_\epsilon)+\epsilon'}{\Tr(\nabla f(\bar m_\epsilon)-\epsilon'}\right)
%        \right)
%        \\
%        \to&
%        \log\left(\frac{\Tr(\nabla f(\bar m_\epsilon)+\epsilon'}{\Tr(\nabla f(\bar m_\epsilon)-\epsilon'}\right)
%        \text{ as } t\to\infty
%        .
%    \end{align}
%    Since $\epsilon'>0$ is arbitrary, by taking $\inf$, 
%    \begin{equation}
%        \frac{1}{t}\sum_{i=0}^t
%        \left|\log\left(\frac{\E[\mathcal{Q}_z(\theta_{t})]}{\E[\mathcal{Q}_z(\theta_{t+1})]}\right)\right|
%        \to 0
%        \text{ as } t\to\infty
%        .
%    \end{equation}
%    }

\subsection{Potential function}\label{subsec:potentialfunction}

The first step of the analysis of the upper convergence rate is to define a potential function of the state of the (1+1)-ES, $V:\Thetazero \to \R$, to measure the progress toward the optimum.
If we construct a potential function such that $\log(f(m))\leq V(\theta)$ holds for all $\theta \in \Thetazero$, then we have for any $\theta_0 \in \Thetazero$,
\begin{equation}
\underset{t\to\infty}{\limsup}\frac{1}{t}\log\left(\frac{f(m_t)}{f(m_0)}\right)
\leq
\underset{t\to\infty}{\limsup}\frac{1}{t}\log\left(\frac{V(\theta_t)}{V(\theta_0)}\right)
.
\end{equation}
Therefore, in light of \Cref{prop:f-norm}, we can obtain the upper convergence rate of the (1+1)-ES on $h \in \mathcal{T}^{\gamma}(f)$ by analyzing the convergence of $V(\theta_t)$. 
It results in analyzing the drift of $X_t = V(\theta_t)$ in light of \Cref{prop:convergence_rate}.

Before constructing a potential function, we explain the motivation of introducing a potential function instead of analyzing $\log(f(m))$ directly.
It has been known empirically and theoretically that the (1+1)-ES exhibits two different phases through the optimization on several functions \cite{akimoto2018drift,akimoto2020global,morinaga2019generalized}.
One is that the step-size $\sigma$ is well adapted and $\log(f(m))$ decreases steadily. 
The other is that the step-size $\sigma$ is on the way to a desirable range and $\log(f(m))$ has almost no progress subsequently.
In fact, when trying to evaluate the drift of $X_t = \log(f(m_t))$ using \Cref{lemma:qualitygain}, it is easy to see that the upper bound of the expected log-progress is non-negative if $\sigma / \norm{\nabla f(m)}$ is too large. 
Moreover, the expected log-progress diminishes if $\sigma / \norm{\nabla f(m)}$ approaches zero.
Hence, the drift of $X_t = \log(f(m_t))$ cannot be upper-bounded by a negative constant for all $\theta \in \Thetazero$. 
In related works \cite{akimoto2018drift,akimoto2020global,morinaga2019generalized}, this issue has been mitigated by constructing a potential function that penalizes the situation that the step-size is not well-adapted.

The potential function used in this work is defined below. The basic idea is borrowed from \cite{morinaga2019generalized,morinaga2021convergence}, but the constants are revised to obtain the desired convergence rate bound.


\begin{definition}[Potential function]\label{def:potential}
Consider the (1+1)-ES solving $f \in \mathcal{S}_{d,L,U}$ with $\aup$ and $\adown$ satisfying $p_\mathrm{target} \in I_q$, where $I_q$ is defined in \eqref{eq:def_Iq}.
Then, in light of \Cref{lemma:cases}, we can choose 
constants $q^\mathrm{low}\in I_q$ and $q^\mathrm{high}\in I_q^\mathrm{high}(q^\mathrm{low})$
so that  
\begin{align}
 q^\mathrm{low}<p_\mathrm{target}<q^\mathrm{high} .
 \label{eq:ptarget_condition}
\end{align}
For $\theta\in\Thetazero$, the potential function $V(\theta)$ for the (1+1)-ES solving $f$ is defined as
\begin{multline}
V(\theta)
= \log\left(f(m)\right) \\
+ v \cdot \log^+ \left(\frac{s\cdot\sqrt{Lf(m)}}{\sigma E_\mathcal{Q} }\right) 
+  v \cdot \log^+ \left(\frac{\sigma E_\mathcal{Q}}{\ell\cdot\sqrt{Lf(m)}}\right) ,
\label{eq:potential}
\end{multline}
where 
$\log^+(x) := \log(x) \cdot \ind{x\geq 1}$, 
$s = \sqrt{ 2 } \aup\cdot B_{\inf}^\mathrm{high}(q^\mathrm{high})$,
$\ell = \sqrt{2} \adown\cdot 
B_{\sup}^\mathrm{low}(q^\mathrm{low})$,
$E_\mathcal{Q} =
\sup_{\theta\in\Thetazero} \left\{\E[\mathcal{Q}_z(\theta)]\right\}$,
and
\begin{equation}
v = \min\left\{ \frac{w}{4 \log(\aup / \adown)} , 1 \right\}
\label{eq:v}
\end{equation}
with
% \begin{multline}
%     w = 
%     \frac{L}{E_\mathcal{Q}}
%     \cdot
%     \frac{
%     \underset{\theta\in\Theta\setminus\{\theta|m\neq x^*\}}{\inf}
%     \left\{
%     B_\theta^\mathrm{high}(q^\mathrm{high})
%     \right\}           
%     }{\kappa_{\inf}}
%     \\
%     \cdot
%     \left(
%     \sqrt{\frac{2}{\pi}}\cdot\kappa_{\inf}
%     -
%     \underset{\theta\in\Theta\setminus\{\theta|m\neq x^*\}}{\sup}
%     \left\{
%      B_\theta^\mathrm{low}(q^\mathrm{low})
%      \right\}
%     \right)
%     \cdot Q
%     >0
%   .
%   \label{eq:w}
% \end{multline}
\begin{equation}
w = 
\frac{L}{E_\mathcal{Q}}
\cdot
\frac{
B_{\inf}^\mathrm{high}(q^\mathrm{high})
}{\kappa_{\inf}}
\cdot
\left(
\sqrt{\frac{2}{\pi}}\cdot\kappa_{\inf}
-
B_{\sup}^\mathrm{low}(q^\mathrm{low})
\right)
\cdot Q .
\label{eq:w}
\end{equation}
%        \mori{$w_\theta$ is essentially an upper drift we seek for. When applying convergence rate theorem and drift theorem, we may want to have a constant upper drift which does not depend on $\theta$. One naive solution to eliminate dependency of $\theta$ is to take $\inf$ or $\sup$ of $B_\theta^\mathrm{low}, B_\theta^\mathrm{high}$ and $\E[\mathcal{Q}(\theta)]$ in $w_\theta$. I suppose that to take $\inf_\theta B_\theta^\mathrm{low}(q^\mathrm{low})$ and $\inf_\theta B_\theta^\mathrm{high}(q^\mathrm{high})$ probably does not affect the rate we expect, but we will lost some tightness of the rate if taking $\sup$ like $\sup_\theta \E[\mathcal{Q}_z(\theta)] = dU$. For now, I give up on the tightness of the rate $L/\Tr(\nabla f(\bar m))$ in the case of smooth storong convex function, and I will reflect after finishing the proof.}
%        \yohe{I see. Yes, tightness seems to be very difficult (as I guess the tightness for the first-order methods are also not established.) For our future work, I suggest to keep all lemmas as they are, and make a corollary of Lemma 4 where we take $\inf B^{high}$ and $\sup B^{low}$ and derive the corresponding results. }
\end{definition}

%It is easy to see that $0 < s < \ell$ for any $\theta\in\Thetazero$ from the definition of $q^\mathrm{high}$ in \eqref{eq:def_qhigh}.

% Although we have $E_\mathcal{Q} = dU$ assuming only \Cref{asm:L_U_2diff}, we also obtain $E_\mathcal{Q}$ is exactly the trace of the Hessian matrix in the case that $f$ is a convex quadratic function which is a subset of functions satisfying \Cref{asm:L_U_2diff}.

%In the following section, we show that the rate of convergence is essentially represented by $\Psi$ which is the infimum of the expected progress rate of $\log(f(m))$ over $\Theta_r$.
%Here $\Theta_r$ is supposed to be a set of $\theta$ in which the (1+1)-ES makes a good progress of optimization.

%\yohe{was there any problem with setting $s$ and $\ell$ accordingly and not touching the range of the reasonable step size?}
%    \mori{
%    By setting $s, \ell$ so, we surely can deal with the term $\log(\E[\mathcal{Q}_{t+1}]/\E[\mathcal{Q}_t])$ which emerges when analyzing the penalty terms at $t+1$ (which motivates us to change the current analysis), but changing $s, \ell$ will also affect the penalty terms at $t$, which is the terms we want to remain as it is.
%    Just changing the reasonable step-size range does not change the potential function itself, and is more convenient for us.
%    }
%    
%        \mori{
%    [memo] We may want to enlarge the reasonable step-size range as
%        \begin{multline}
%            \frac{s\cdot\sqrt{L f(m)}}{\aup\cdot\inf_\theta\{\E[\mathcal{Q}_z(\theta)]\}}
%            \leq \sigma \leq
%            \frac{\ell\cdot\norm{\nabla f(m)}}{\sqrt{2}\adown\cdot\sup_\theta\{\E[\mathcal{Q}_z(\theta)]\}}
%            ,
%        \end{multline}
%        while let the potential function remain as it is.
%        This leads to the same upper drift we obtain in GECCO'21 except for the upper bound of the expected log-progress.
%        By analyzing with the above range, we will obtain an upper drift with an upper bound of the expected log-progress
%        \begin{equation}
%            \frac{L\cdot \inf_\theta\{B_\theta^\mathrm{high}(q^\mathrm{high}\}}{2\sup_\theta\{\E[\mathcal{Q}_z(\theta)]\}}
%            \cdot\left(
%            \frac{\sup_\theta\{B_\theta^\mathrm{low}(q^\mathrm{low})\}}{\kappa_{\inf}}\cdot\frac{\E[\mathcal{Q}_z(\theta)\cdot\ind{z_e\leq 0}]}{\inf_\theta\{\E[\mathcal{Q}_z(\theta)\}} - \sqrt{\frac{2}{\pi}}
%            \right)\cdot Q
%            .
%        \end{equation}
%        This will lock up $p_\mathrm{target}$ into $1/2$ as $d\to\infty$.
%        }

The difference of the potential function values in consecutive steps is bounded in several ways depending on the current state. 
The proof is provided in \Cref{apdx:lemma:potential_decrease}.
%\yohe{This lemma is very nice. I didn't know we can have such a nice lemma. I think this is quite helpful to understand the proof idea.}
%\mori{And also note that this lemma does not assume any property of $f$ :) }
\begin{lemma}[Potential bound]\label{lemma:potential_decrease}
Consider the (1+1)-ES solving $f \in \mathcal{S}_{d,L,U}$ with $\aup$ and $\adown$ satisfying $p_\mathrm{target} \in I_q$. Let $\{\theta_t\}_{t\geq 0}$ be a sequence of the parameters of the (1+1)-ES.
Suppose $\theta_t\in\Thetazero$. 
%Suppose $m_t \neq x^*$. 
Then,
\begin{equation}
V(\theta_{t+1}) - V(\theta_t)\leq 
\left(1-\frac{v}{2}\right)
\cdot \log\left(\frac{f(m_{t+1})}{f(m_t)}\right) + v\cdot\log\left(\frac{\aup}{\adown}\right)
,
\label{eq:general_potential_decrease}
\end{equation}
and
\begin{equation}
V(\theta_{t+1}) - V(\theta_t)
>
\left(1+v\right)\cdot\log\left(\frac{f(m_{t+1})}{f(m_t)}\right)
-2v\cdot\log\left(\frac{\aup}{\adown}\right)
.
\label{eq:constant_lower_bound_of_V}
\end{equation}
In addition, the following statements holds.
%\begin{enumerate}
%\renewcommand{\labelenumi}{(\roman{enumi})}

(i) If $\sigma_t < \frac{s \cdot \sqrt{ L f(m_t) }}{ \aup\cdot E_\mathcal{Q}}$, letting $\inddown := \ind{f(m_t + \sigma_t\cdot z_t) > f(m_t)}$, we have
\begin{equation}
V(\theta_{t+1}) - V(\theta_t) \leq
v \cdot\left( \log\left(\frac{\aup}{\adown}\right)\cdot\inddown - \log(\aup)\right)
.
\label{eq:small_sigma_potential_decrease}
\end{equation}
%\todo{Is this definition correct? Check this definition and the first equality in the proof (86--92).}

(ii) If $\sigma_t > \frac{\ell \cdot \sqrt{L f(m_t)} }{ \adown\cdot E_\mathcal{Q}}$, letting $\indup := \ind{f(m_t + \sigma_t\cdot z_t) \leq f(m_t)}$, we have
\begin{equation}
V(\theta_{t+1}) - V(\theta_t)\leq 
v \cdot\left( \log\left(\frac{\aup}{\adown}\right)\cdot\indup - \log\left(\frac{1}{\adown}\right)\right)
.
\label{eq:large_sigma_potential_decrease}
\end{equation}
%\end{enumerate}
\end{lemma}
% Eq. \eqref{eq:small_sigma_potential_decrease} and \eqref{eq:large_sigma_potential_decrease} correpond to a situation that the 2nd term of $V(\theta)$ compensates for a stagnation of $\log(f(m))$, and a situation that the 3rd term do it.

Finally, we derive condition $C1$ of \Cref{prop:convergence_rate} for each of the three situations: ($i$) too small step-size situation, where the progress is made by increasing the step-size; ($ii$) too large step-size situation, where the progress is made by decreasing the step-size; and ($iii$) reasonable step-size situation, where the progress is made by moving toward the optimum.
Its proof is provided in \Cref{apdx:cor:expected_potential_decrease}.
\begin{corollary}[Expected potential decrease]\label{cor:expected_potential_decrease}
Consider the (1+1)-ES solving $f \in \mathcal{S}_{d,L,U}$ with $\aup$ and $\adown$ satisfying $p_\mathrm{target} \in I_q$. Let $\{\theta_t\}_{t\geq 0}$ be a sequence of the parameters of the (1+1)-ES.
Suppose $\theta_t \in\Thetazero$. 
%Suppose $m_t \neq x^*$. 
Then, the following statements hold.
% \begin{enumerate}
% \renewcommand{\labelenumi}{(\roman{enumi})}

(i) If $\sigma_t < \frac{s \cdot \sqrt{ L f(m_t) }}{ \aup\cdot E_\mathcal{Q}}$,
\begin{multline}
\E\left[V(\theta_{t+1}) - V(\theta_t) |\mathcal{F}_t\right]
\\
\leq
\min
\left\{
    \frac{w}{4}, \log\left(\frac{\aup}{\adown}\right)
\right\}
\cdot
    (p_\mathrm{target} - q^\mathrm{high})
<0
.
\label{eq:small_sigma_expected_potential_decrease}
\end{multline}

(ii) If $\sigma_t > \frac{\ell \cdot \norm{\nabla f(m_t)} }{ \sqrt{2}\adown\cdot \E_{z_t}[\mathcal{Q}_{z_t}(\theta_t)]}$,
\begin{multline}
\E\left[V(\theta_{t+1}) - V(\theta_t) |\mathcal{F}_t\right]
\\
\leq
\min
\left\{
    \frac{w}{4}, \log\left(\frac{\aup}{\adown}\right)
\right\}
\cdot
    (q^\mathrm{low} - p_\mathrm{target})
<0
.
\label{eq:large_sigma_expected_potential_decrease}
\end{multline}

(iii) If  $ \frac{s \cdot \sqrt{ L f(m_t) }}{ \aup\cdot E_\mathcal{Q}} \leq \sigma_t \leq \frac{\ell \cdot \norm{\nabla f(m_t)} }{ \sqrt{2}\adown\cdot \E_{z_t}[\mathcal{Q}_{z_t}(\theta_t)]}$,
\begin{equation}
\E\left[V(\theta_{t+1}) - V(\theta_t) |\mathcal{F}_t\right]
\leq -\frac{w}{4} 
<0
.
\label{eq:reasonable_sigma_expected_potential_decrease}
\end{equation}
%\end{enumerate}
\end{corollary}

% \mori{BTW, I'm not confident whether this word "expected exponential potential decrease" sounds natural.}\yohe{Not sure as well. Does it make sense to say "multiplicative drift?"}\mori{Sounds not too bad while it does not precisely describe the result of the following lemma. As I found many examples to use a word "multiplicative decrease", I state as multiplicative potential decrease for now (of course while this is not presice similarly).}
% In light of \Cref{lemma:potential_decrease}, a sufficient condition for \Cref{theo:additive_drift} is also offered.
% Its proof is provided in \Cref{apdx:cor:expected_multiplicative_potential_decrease}.
% \begin{corollary}[Expected multiplicative potential decrease]\label{cor:expected_multiplicative_potential_decrease}
%         In light of \Cref{lemma:cases} and \Cref{lemma:potential_decrease}, we have a constant negative upper bound of the expected multiplicative potential decrease as follows.
%         \begin{enumerate}
%         \renewcommand{\labelenumi}{(\roman{enumi})}
%         \item If $\sigma_t < \frac{s \cdot \sqrt{ L f(m_t) }}{ \aup\cdot E_\mathcal{Q}}$,
%         \begin{multline}
%             \E\left[\exp(\eta(V(\theta_{t+1}) - V(\theta_t) ))|\mathcal{F}_t\right]
%             \leq
%             \rho_s
%             ,
%             \\
%             \text{ where }
%             \rho_s
%             =
%             \left(
%             \exp\left(\eta
%                 \cdot\min\left\{\frac{w}{4}, \log\left(\frac{\aup}{\adown}\right)\right\}
%             \right)
%             \cdot(1-q^\mathrm{high})
%             + q^\mathrm{high}
%             \right)
%              \\
%             \cdot
%             \exp\left(-\eta
%                 \cdot (1 - p_\mathrm{target})
%                 \cdot\min\left\{\frac{w}{4}, \log\left(\frac{\aup}{\adown}\right)\right\}
%             \right)               
%             .
%             \label{eq:small_sigma_expected_multiplicative_potential_decrease}
%         \end{multline}
%         Moreover, there exists $\eta_s>0$ such that given any $v\in(0, 1)$, $\rho_s<1$ if $\eta<\eta_s$.

%         \item If $\sigma_t > \frac{\ell \cdot \norm{\nabla f(m_t)} }{ \sqrt{2}\adown\cdot E_{z_t}[\mathcal{Q}_{z_t}(\theta_t)]}$,
%         \begin{multline}
%             \E\left[\exp(\eta(V(\theta_{t+1}) - V(\theta_t) ))|\mathcal{F}_t\right]
%             \leq
%             \rho_\ell
%             ,
%             \\
%             \text{ where }
%             \rho_\ell
%             =
%             \left(
%             \exp\left(\eta
%                 \cdot\min\left\{
%                 \frac{w}{4},
%                 \log\left(\frac{\aup}{\adown}\right)
%                 \right\}
%             \right)
%             \cdot q^\mathrm{low}
%             + 1-q^\mathrm{low}
%             \right)
%             \\
%             \cdot
%             \exp\left(-\eta\cdot
%                 p_\mathrm{target}
%                 \cdot\min\left\{
%                 \frac{w}{4},
%                 \log\left(\frac{\aup}{\adown}\right)
%                 \right\}
%             \right)    
%             .
%             \label{eq:large_sigma_expected_multiplicative_potential_decrease}
%         \end{multline}
%         Moreover, there exists $\eta_\ell>0$ such that given any $v\in(0, 1)$, $\rho_\ell<1$ if $\eta<\eta_\ell$.

%         \item Otherwise, that is, if  $ \frac{s \cdot \sqrt{ L f(m_t) }}{ \aup\cdot E_\mathcal{Q}} \leq \sigma_t \leq \frac{\ell \cdot \norm{\nabla f(m_t)} }{ \sqrt{2}\adown\cdot E_{z_t}[\mathcal{Q}_{z_t}(\theta_t)]}$,
%         \begin{multline}
%             \E\left[\exp(\eta(V(\theta_{t+1}) - V(\theta_t) ))|\mathcal{F}_t\right]
%             \leq
%             \rho_r
%             \\
%             \text{ where }
%             \rho_r
%             =
%             \exp\left(-\eta\cdot p_\mathrm{target}\cdot\min\left\{\frac{w}{4}, \log\left(\frac{\aup}{\adown}\right)\right\}\right)
%             <1
%             .
%             \label{eq:reasonable_sigma_expected_multiplicative_potential_decrease}
%         \end{multline}
%       \end{enumerate}

% \end{corollary}


Condition $C3$ of \Cref{prop:convergence_rate} is obtained in the following lemma.
Its proof is provided in \Cref{apdx:lemma:v_variance_bound}.
%\todo{This lemma is not necessary to be an independent lemma. Incorporate the proof into that of Theorem 1.}\yohe{In my openion, it is not a bad idea to keep the following lemma as is. My reasoning is that by separating this kind of (not very interesting but not very short) proof, the main stream of the proof of the main theorem will be clearer. }\mori{I see. I'm now sure it is useful for readers to understand the proof idea.}
\begin{lemma}[Variance bound of the potential function]\label{lemma:v_variance_bound}
Consider the (1+1)-ES solving $f \in \mathcal{S}_{d,L,U}$ with $d > 3$ and $\aup$ and $\adown$ satisfying $p_\mathrm{target} \in I_q$. Let $\{\theta_t\in\Thetazero\}_{t\geq 0}$ be a sequence of the parameters of the (1+1)-ES.% and suppose that $m_0 \neq x^*$. 
Then, $\sum_{t=1}^{\infty}{\color{blue}\Var[V(\theta_{t})]} / t^2 < \infty$.
\end{lemma}



\subsection{Upper bound of the upper convergence rate}
%\yohe{
%\eqref{eq:order_B} can be refined for \Cref{asm:L_U_2diff} with twice continuous differentiability around $x^*$. The proof idea is as follows. 
%
%Our convergence rate theorem implies the probability one convergence of $(m^t, \sigma^t) \to (x^*, 0)$. That is, there is a null set and take the random variable $\omega$ outside the null set. Then, for each $\omega$ outside the null set, we have $\lim_t \norm{m^t(\omega) - x^*} = 0$ and $\lim_t \sigma^t(\omega) = 0$. It implies that for any $\epsilon > 0$, there exists $T_{\omega}$ such taht $\norm{m^t(\omega) - x^*} < \epsilon$ and $\sigma^t < \epsilon$ for all $t \geq T_{\omega}$. Obviously, the convergence rate of the original sequence under $\omega$ is the same as the convergence rate of the sequence trancated for $t \geq T_{\omega}$. Then, it means, $\sup_{\theta \in \Theta}$ in the definition of $E_\mathcal{Q}$ can be replaced with $\sup_{\theta \in \Theta \wedge \norm{m - x^*} < \epsilon \wedge \sigma < \epsilon}$. Because of Lemma 1 (4), we see that $E_\mathcal{Q}$ can be replaced with $\Tr(\nabla^2 f(x^*))$. 
%    }
%\mori{I'm wondering about how come $\lim_{t\to\infty} \sigma^t(\omega) = 0$ in light of our convergence rate theorem.}
%\yohe{You are right. It was a bit surprising to me that the potential function can diverge toward minus infinity even if sigma is constant. So, please forget about the above point. (In the potential function that I am using to analyze the weighted recombination ES, we can show the convergence of sigma from the convergence of potential. )}
%\mori{
%    Let me try to show the convergence of $\log(\sigma_t)$.
%    
%    Let $\{x_i^*\in\R^d\}_{i\in\mathbb{N}}$ be a set of stationary points of $f$, that is, $\norm{\nabla f(x_i^*)} = 0$ for all $i\in\mathbb{N}$.
%    Assume $\Pr_{z_t}[f(m_t+\sigma_t \cdot z_t)\leq f(m_t)]\leq 1/2$ on $f$ for any $\theta\in\Theta\setminus\{\theta\mid m\neq x_i^*, \forall i\in\mathbb{N}\}$.
%    Set $\aup$ and $\adown$ so that $\aup\adown<1$.
%    Given any $\sigma_t>0$, 
%    \begin{equation}
%        \log\left(\frac{\sigma_{t+1}}{\sigma_t}\right) = \log(\aup)\cdot \indup + \log(\adown)\cdot \inddown
%        .
%    \end{equation}
%    Then, 
%    \begin{align}
%        \E\left[\log\left(\frac{\sigma_{t+1}}{\sigma_t}\right) \mid\mathcal{F}_t\right]
%        =& \log\left(\frac{\aup}{\adown}\right)\cdot \Pr_{z_t}[f(m_t+\sigma_t \cdot z_t)\leq f(m_t)] + \log(\adown)
%        \\
%        \leq& \frac12 \cdot\log\left(\frac{\aup}{\adown}\right) + \log(\adown)
%        \\
%        =& \log\left(\sqrt{\aup \adown}\right)
%        \\
%        <& 0
%        .
%    \end{align}
%    Therefore, in light of \Cref{prop:convergence_rate}, there exists a linear function with a negative gradient which bounds $\log(\sigma_t)$ from above, and then $\log(\sigma_t) \to -\infty$ as $t\to\infty$.
%    
%    [remark] We can find $\aup>1$ and $\adown<1$ such that $\aup\adown<1 \Leftrightarrow \frac{\aup}{\adown}<\frac{1}{\adown^2}$, and it does not contradict the setting $\log(\aup/\adown)\in\Theta(1/d)$ at least as long as we set $\adown$ such that $\log(1/\adown^2)\in\Omega(1/d)$ (of course this is always satisfied as $d\to\infty$).
%    So we can just care whether $\log(\aup/\adown)\in\Theta(1/d)$ is satisfied.
%}
%\yohe{I believe $\aup \adown < 1$ is not acceptable. Shouldn't it indicate that the target probability $> 1/2$?}
%\mori{Ah, definitely. It shouldn't be so sweet deal.}

An upper bound of the upper convergence rate of the (1+1)-ES on objective function $h\in\mathcal{P}_{d, L, U}^\gamma$ (\Cref{def:convergence}) is derived by using \Cref{prop:convergence_rate} and \Cref{cor:expected_potential_decrease,lemma:v_variance_bound}.

\begin{theorem}[Upper convergence rate bound]\label{theorem:lower}
Let $\{(\tilde{\theta}_t, \mathcal{F}_t)\}_{t \geq 0} = \texttt{ES}(h, \tilde{\theta}_0, \{z_t\}_{t \geq 0})$ be the state sequence of the (1+1)-ES solving $h\in\mathcal{P}_{d, L, U}^\gamma$ with $d > 3$, where $\tilde{\theta}_0 \in \Theta\setminus\{\tilde{\theta}\mid \tilde{m}= x^\mathrm{opt}\}$ and $h(\tilde{m}_0) \leq \gamma$.
Suppose that $\aup$ and $\adown$ are set so that $p_\mathrm{target} \in I_q$.
%\yohe{what are the points of introducing two sequences? probably, it is a copy-paste error.}
%   Let $f\in\mathcal{P}_{d, L, U}$ be a function satisfying \Cref{asm:L_U_2diff}.
% Assume that function $f$ corresponding to $h$ in \Cref{def:problem} satisfies \eqref{eq:tr_cond}.
% Suppose that $\aup$ and $\adown$ in \Cref{algo} are set so that $p_\mathrm{target}\in I_q$, where $p_\mathrm{target}$ is defined in \eqref{eq:p_target-def} and $I_q$ is defined in \eqref{eq:def_Iq}. 
%
Then, the upper convergence rate on $h$ satisfies $\exp(-\mathrm{CR}^\mathrm{upper}) \leq \exp(-B_{d, L, U}^\mathrm{upper})$, where
\begin{multline}
B_{d, L, U}^\mathrm{upper}
= \frac12\cdot\sup_{q^\mathrm{low}, q^\mathrm{high}}
\min\left\{\frac{w}{4}, \log\left(\frac{\aup}{\adown}\right)\right\}
\\
\cdot \min\{p_\mathrm{target} -  q^\mathrm{low}, q^\mathrm{high} - p_\mathrm{target} \} > 0 ,
\label{eq:B}
\end{multline}
where $\sup_{q^\mathrm{low}, q^\mathrm{high}}$ is taken over $q^\mathrm{low}\in I_q$ and $q^\mathrm{high}\in I_q^\mathrm{high}(q^\mathrm{low})$ satisfying \eqref{eq:ptarget_condition} with 
$I_q^\mathrm{high}$ defined in \eqref{eq:def_qhigh}, and $w$ is defined in \eqref{eq:w}.
%\mori{Think again later : we can obtain a better, or larger, $B$ with a larger $q^\mathrm{high}$. However, I define $q^\mathrm{high}$ as \eqref{eq:def_qhigh} which takes $\inf$. Why I take $\inf$?}\yohe{I don't have the anser, but the conditionfor $q^\mathrm{high}$ can be satisfied with $q$ arbitrarily close to $1/2$ (this is how we prove the existance). I was also wondering why we want to take $\inf$ (and $\sup$ over $\Theta$). Let me try to reformulate the Lemma statement.}\yohe{What I see now is that the statement "we can obtain a better, or larger, $B$ with a larger $q^\mathrm{high}$" is not necessarily correct. A greater $q^\mathrm{high}$ results in a smaller $w$. }\mori{Yes, I see.}
% 
%   Moreover, \yohe{state the condition on $p_\text{target}$. I think it is one of the strong point of the approach that we can set the hyper-parameter $p_\text{target}$ without using structures of the objective function. }
%   %\mori{I state it in the above about $\aup, \adown$. Or, are you saying it should be emphasized when stating the following our main result?}
%   as long as the hyper-parameter $\aup$ and $\adown$ are set so that $p_\mathrm{target}\in\left(\Phi\left(-\sqrt{\frac{2}{\pi}}\right), \frac12 \right)$, the following statement holds.
%   \begin{equation}
%     \lim_{d \to \infty}
%     \left\{
%     \inf_{h \in \mathcal{P}_{d, L, U}}
%     \left\{
%     \frac{ \mathrm{CR}^\text{upper}}{ \min\left\{ \frac{L}{E_\mathcal{Q}},  \log\left(\frac{\aup}{\adown}\right)\right\} }
%     \right\}
%     \right\}
%     \in (0, \infty) .
%     \label{eq:order_B}
%   \end{equation}
\end{theorem}


\begin{proof}[Proof of \Cref{theorem:lower}]
%\mori{I correct the proof so that each reasoning is stated on $f$, not $h$.}
Let $h \in \mathcal{T}^{\gamma}(f)$ and $f \in \mathcal{S}_{d, L, U}$ as defined in \Cref{def:problem}.
Let $S_T : (x, \log(\sigma)) \mapsto (x - x^\mathrm{opt}, \log(\sigma))$ and $\{(\theta_t, \mathcal{F}_t)\}_{t \geq 0} = \texttt{ES}(f, S_T(\tilde{\theta}_0), \{z_t\}_{t \geq 0})$ be the state sequence of the (1+1)-ES solving $f$.
Because the singleton $\{x^*\}$ is of zero Lebesuge measure, given $m_0 \neq x^*$, we have $m_t \neq x^*$ for all $t$ with probability one.
In light of \Cref{prop:f-norm}, the upper convergence rate on $h$, $-\mathrm{CR}^\mathrm{upper}$ is equal to
$\limsup_{t \to \infty} \frac{1}{2t} \log\left(\frac{f(\theta_t)}{f(\theta_0)}\right)$
with probability one.

As shown in the proof of \Cref{cor:expected_potential_decrease}, $w$ in \eqref{eq:w} is positive since $q^\mathrm{low}\in I_q$ satisfies $\sqrt{\frac{2}{\pi}}\cdot\kappa_\mathrm{inf}>B_\mathrm{sup}^\mathrm{low}(q^\mathrm{low})$ in light of \Cref{lemma:cases}.
Therefore, it is clear that $v$ defined in \eqref{eq:v} and hence, $B_{d, L, U}^\mathrm{upper}$ defined in \eqref{eq:B} are positive.

Let $X_t = V(\theta_t)$ in \Cref{prop:convergence_rate} with $V$ defined in \Cref{def:potential}.
Condition C1 --- $\E[V(\theta_{t+1}) - V(\theta_t)\mid\mathcal{F}_t]\leq -B$ for all $t\geq 0$ --- is satisfied with $B = 2\cdot B_{d, L, U}^\mathrm{upper}$ in light of \Cref{cor:expected_potential_decrease}.
Condition C3 --- $\sum_{t=1}^{\infty}\Var[V(\theta_{t}) \mid \mathcal{F}_{t-1}] / t^2 < \infty$ --- is satisfied for $d>3$ in light of \Cref{lemma:v_variance_bound}.
Therefore, with probability one, we obtain
%   \begin{equation}
%   \limsup_{t \to \infty} \frac{1}{t} \frac{V(\theta_t)}{V(\theta_0)} \leq - B .
% \end{equation}
$\limsup_{t \to \infty} \frac{1}{t} \left( V(\theta_t) - V(\theta_0) \right) \leq - 2\cdot B_{d, L, U}^\mathrm{upper}$.
Because $\log(f(m)) \leq V(\theta)$ for any $\theta \in \Thetazero$, we have
$\limsup_{t \to \infty} \frac{1}{2 t} \log\left(\frac{f(\theta_t)}{f(\theta_0)}\right) \leq - B_{d, L, U}^\mathrm{upper}$.
Hence, we obtain $\exp(-\mathrm{CR}^\text{upper}) \leq \exp(- B_{d, L, U}^\mathrm{upper}) < 1$.
\end{proof}

If $d$ is sufficiently large and the hyper-parameters $\aup$ and $\adown$ are selected so that $\Phi(- \sqrt{2/\pi}) < p_\mathrm{target} < 1/2$, the bound of the upper convergence rate is $\exp(- B_{d,L,U}^\mathrm{upper}) \in \exp\left(-\Omega\left(\min\left\{ \frac{L}{E_\mathcal{Q}},  \log\left(\frac{\aup}{\adown}\right)\right\}\right)\right)$.
\begin{corollary}[Scaling of the Upper Convergence Rate]\label{cor:upper}
Let $q^\mathrm{low} \in \left(\Phi\left(-\sqrt{\frac{2}{\pi}}\right), \frac{1}{2}\right)$ and 
$q^\mathrm{high} \in \left(
%1 - \Phi\left( \frac{\adown}{\aup} \Phi^{-1}\left( 1 - q^\mathrm{low}\right)\right)
\Phi\left( \frac{\adown}{\aup} \Phi^{-1}\left( q^\mathrm{low}\right)\right)
, \frac{1}{2}\right)$, which are independent of $d$. 
Suppose that $\aup$ and $\adown$ are set so that $q^\mathrm{low} < p_\mathrm{target} < q^\mathrm{high}$ for all $d$ ($\aup$ and $\adown$ can depend on $d$). 
Then, there exists $D > 3$ such that for any $d \geq D$, (i) \Cref{asm:tr_cond} is satisfied under \Cref{asm:L_U_2diff}, (ii)
$q^\mathrm{low} \in I_q$ and $q^\mathrm{high} \in I_q^\mathrm{high}(q^\mathrm{low})$, and (iii) the following statement holds
\begin{equation}
\lim_{d \to \infty}
\left\{
\inf_{h \in \mathcal{P}_{d, L, U}^\gamma}
\left\{
\frac{ B_{d, L, U}^\mathrm{upper} }{ \min\left\{ \frac{L}{E_\mathcal{Q}},  \log\left(\frac{\aup}{\adown}\right)\right\} }
\right\}
\right\}
\in (0, \infty) ,
\label{eq:order_B}
\end{equation}    
where $B_{d, L, U}^\mathrm{upper}$ is as defined in \eqref{eq:B}.
Moreover, we have $\frac{L}{E_\mathcal{Q}} \geq \frac{L}{dU}$. If $h \in \mathcal{T}^\gamma(f)$ for $f(x) = \frac{1}{2} x^\mathrm{T} H x$, where $H$ is a symmetric matrix whose eigenvalues are bounded in $[L, U]$, we have $\frac{L}{E_\mathcal{Q}} = \frac{L}{\Tr(H)}$.
\end{corollary}
{
    We remark on the consequences.
    The hyper-parameters $\aup$ and $\adown$ are typically chosen such that $\log(\aup / \adown)\in\Omega_{d \to \infty}(1/d)$.
    In light of \Cref{cor:upper}, the upper convergence rate $\exp\left(-\mathrm{CR}^\mathrm{upper}\right)$ is in $ \exp\left(- \Omega_{d \to \infty}\left(\frac{L}{E_\mathcal{Q}} \right)\right)$ under a choice such that $\log(\aup / \adown) \in \Omega_{d \to \infty}(1/d)$.
    Otherwise (i.e. if $\log(\aup / \adown) \in o_{d \to \infty}(1/d)$), we have $\exp\left(-\mathrm{CR}^\mathrm{upper}\right) \in\exp\left(- \Omega_{d \to \infty}\left(\log\left(\frac{\aup}{\adown}\right) \right)\right) = \exp\left(-\Omega_{d \to \infty}\left(\frac{\log\left( 1/\adown\right)}{p_\mathrm{target}}\right)\right)$.
    This is rather intuitive for the following reason. $\norm{m_t - x^*}$ does not converge faster than $\sigma_t$ because $\sigma_t$ must be proportional to $\norm{m_t - x^*}$ to produce a sufficient decrease. The speed of the decrease in $\sigma_t$ is $\adown$. Therefore, the upper convergence rate should not be smaller than $\adown$.

The dependency of the convergence rate on the trace of the Hessian matrix is derived under the optimal step-size situation for a different ES variant \cite{akimoto2020tcs}. It has also been derived for a different ES variant in \cite{beyer2013dynamics} though it is not a rigorous convergence rate analysis.}
\begin{proof}
The first claim, (i), is proved in \Cref{cor:tr_cond}. In the following, we prove claims (ii) and (iii).

Continuing from the proof of \Cref{theorem:lower}, we consider $f$ corresponding to $h$. Let $\{\theta_t\}$ be defined in the proof of \Cref{theorem:lower}.

In light of \Cref{lemma:variance},
given any $0 < L \leq U$,
we have
\begin{equation}
\underset{d\to\infty}{\lim}
\left\{
\underset{f\in\mathcal{S}_{d, L, U}}{\sup}
\left\{
\mathcal{V}_\mathrm{std} 
\right\}
\right\}
\leq 
\underset{d\to\infty}{\lim}
\left\{
\frac{4U^2}{dL^2}
\right\}
= 0
.
\label{eq:tr_cond_asymptotic}
\end{equation}
Then, we have
$
\underset{d\to\infty}{\lim}
\left\{
B_{\inf}^\mathrm{high}(q)
\right\}
= 2\cdot\Phi^{-1}(1-q)
$
and
$
\underset{d\to\infty}{\lim}
\left\{
B_{\sup}^\mathrm{low}(q)
\right\}
= 2\cdot\Phi^{-1}(1-q)
$.
Moreover, as we prove in the proof of \Cref{cor:tr_cond}, we have
$\kappa_{\inf} \to 2$ as $d\to\infty$.
Therefore, the lower limit of $I_q$ defined as \eqref{eq:def_Iq} approaches $ \Phi\left(-\sqrt{\frac{2}{\pi}}\right)$ as $d\to\infty$ and the lower limit of $I_q^\mathrm{high}(q^\mathrm{low})$ approaches
$\Phi\left( \frac{\adown}{\aup} \Phi^{-1}\left( q^\mathrm{low}\right)\right)$. 
Hence, there exists an integer $D'$ such that claim (ii) holds for all $d \geq D'$.

Finally, we prove claim (iii). 
From \eqref{eq:def_Q}, it is clear that $Q \to q^\text{low}$ as $d\to\infty$.
Therefore, as $d \to \infty$, we obtain
\begin{multline}
\frac{\inf_{f \in \mathcal{S}_{d,L,U}} \{w\}}{L/E_\mathcal{Q}} 
\to 
\Phi^{-1}(1- q^\mathrm{high})
\\
\cdot\left(\sqrt{\frac{2}{\pi}} - \Phi^{-1}(1-q^\mathrm{low})\right) \cdot q^\mathrm{low}
\in (0, \infty).\qedhere
\end{multline}
\end{proof}


% \begin{theorem}[Upper bound of the expected first hitting time]\label{theorem:FHT}
%     Assume $d>3$.
%   Let $\{(\tilde{\theta}_t, \mathcal{F}_t)\}_{t \geq 0} = \texttt{ES}(h, \tilde{\theta}_0, \{z_t\}_{t \geq 0})$ be the state sequence of the (1+1)-ES solving $h\in\mathcal{P}_{d, L, U}$, where $\tilde{\theta}_0 \in \Theta\setminus\{\tilde{\theta}\mid \tilde{m}\neq x^\mathrm{opt}\}$ .
%   Let $f\in\mathcal{P}_{d, L, U}$ be a function satisfying \Cref{asm:L_U_2diff}.
%   Assume $\mathcal{Q}_z(\theta)$ defined as \eqref{eq:q-def} satisfies \eqref{eq:tr_cond}.
%   Set $\aup$ and $\adown$ in \Cref{algo} so that $p_\mathrm{target}\in I_q$ and $p_\mathrm{target} \in\left(\Phi\left(-\sqrt{\frac{2}{\pi}}\right), \frac12\right)$ where $p_\mathrm{target}$ is defined as \eqref{eq:p_target-def} .
%   Let
%   \begin{multline}
%     \rho = 
%     \sup_{q^\mathrm{low}, q^\mathrm{high}}
%     \left\{
%     \max\left\{
%         \exp\left(
%             \eta\cdot\min\left\{\frac{w}{4}, \log\left(\frac{\aup}{\adown}\right)\right\}
%         \right)
%         \cdot (1-q^\mathrm{high})
%         + q^\mathrm{high}
%         ,
%         \right.
%         \right.
%         \\
%         \left.
%         \exp\left(
%             \eta\cdot\min\left\{\frac{w}{4}, \log\left(\frac{\aup}{\adown}\right)\right\}
%         \right)
%         \cdot q^\mathrm{low}
%         + 1 - q^\mathrm{low}
%     \right\}
%     \\
%     \left.
%     \cdot
%     \exp\left(-\eta\cdot p_\mathrm{target}\cdot\min\left\{\frac{w}{4}, \log\left(\frac{\aup}{\adown}\right)\right\}\right)
%     \right\}
%     ,
%     \label{eq:rho}
%   \end{multline}
%   where $\sup_{q^\mathrm{low}, q^\mathrm{high}}$ is taken over $q^\mathrm{low}\in I_q$ and $q^\mathrm{high}\in I_q^\mathrm{high}(q^\mathrm{low})$ satisfying \eqref{eq:ptarget_condition}.
%   Let $X_t = \log\norm{\tilde m_t - x^\mathrm{opt}}$ and $T_\beta^X$ be a first hitting time of $\{X_t\}_{t\geq 0}$ for $\beta\leq\log\norm{\tilde m_0 - x^\mathrm{opt}}$.
%   Fix constants $\eta_s>0$ and $\eta_\ell>0$ in \Cref{cor:expected_multiplicative_potential_decrease} such that $\rho_s<1$ and $\rho_\ell<1$ respectively.
%   Let $\eta = \min\{\eta_s, \eta_\ell\}$.
%   Subsequently, for any $m_0\in\R^d$ and any target value $\beta\leq\log\norm{\tilde m_0 - x^\mathrm{opt}}$,
%   \begin{equation}
%       \E[T_\beta^X]
%       \leq
%       (\log(\rho^{-1/2}))^{-1}\cdot 
%       \left(
%           \eta\left(
%             \log\norm{\tilde m_0 - x^\mathrm{opt}}
%             - \beta
%             -\log\left(\frac{L}{2}\right)
%           \right)
%           +
%           \log(\rho^{-1/2}) + \log(2)
%       \right)
%       \ensapce.
%       \label{eq:FHT_bound}
%   \end{equation}
%   Moreover, 
%   as long as the hyper-parameter $\aup$ and $\adown$ are set so that $p_\mathrm{target}\in\left(\Phi\left(-\sqrt{\frac{2}{\pi}}\right), \frac12 \right)$, the following statement holds.
%   \begin{equation}
%     \lim_{d \to \infty}
%     \left\{
%     \sup_{h \in \mathcal{P}_{d, L, U}}
%     \left\{
%      \E[T_\beta^X]\cdot \min\left\{ \frac{L}{E_\mathcal{Q}},  \log\left(\frac{\aup}{\adown}\right)\right\}
%     \right\}
%     \right\}
%     \in (0, \infty)
%     .
%     \label{eq:order_FHT}
%   \end{equation}
% \end{theorem}

% \begin{proof}
%     Let $T^f$ be the FHT of $\log(f(m_t))$ for $\beta+\log\left(\frac{L}{2}\right)\leq\log(f(m_0))$.
%     Then, in light of \Cref{prop:FHT_mnorm}, $\E[T_\beta^X]\leq\E[T^f]$ for any $h, f\in\mathcal{P}_{d, L, U}$.
%     In light of \Cref{cor:expected_multiplicative_potential_decrease},
%     $\E[\exp(\eta(V(\theta_{t+1}) - V(\theta_t)))]\leq \rho$ .
%     Then, in light of \Cref{theo:additive_drift}, 
%     \begin{multline}
%       \E[T_\beta^X]
%       \leq
%       \E[T^f]
%       \leq
%       \\
%       (\log(\rho^{-1/2}))^{-1}\cdot 
%       \left(
%           \eta\left(
%             \log\norm{\tilde m_0 - x^\mathrm{opt}}
%             - \beta
%             -\log\left(\frac{L}{2}\right)
%           \right)
%           +
%           \log\left(
%           \frac{\rho^{-1/2} - \rho^{1/2}}{1-\rho^{1/2}}
%           \right)
%       \right)
%       \ensapce.
%     \end{multline}
%     Moreover,
%     \begin{equation}
%         \log\left(
%           \frac{\rho^{-1/2} - \rho^{1/2}}{1-\rho^{1/2}}
%         \right)
%         =
%         \log\left(\rho^{-1/2}\right)
%         +
%         \log\left(
%             \frac{1 - \rho}{1-\rho^{1/2}}
%         \right)
%         .
%     \end{equation}
%     Given $\rho\in(0, 1)$,
%     $
%         \log\left(
%             \frac{1 - \rho}{1-\rho^{1/2}}
%         \right)
%         =
%         \log\left(1 + \rho^{1/2}\right)
%         \leq
%         \log(2)
%     $ .
%     Then, 
%     \begin{equation}
%         \log\left(\rho^{-1/2}\right)
%         +
%         \log\left(
%             \frac{1 - \rho}{1-\rho^{1/2}}
%         \right)
%         \leq
%         \log(\rho^{-1/2}) + \log(2)
%         .
%     \end{equation}
%     Now \eqref{eq:FHT_bound} is proved.

%     Next, \eqref{eq:order_FHT} is to be proved by proving
%     \begin{equation}
%         \lim_{d \to \infty}
%         \left\{
%             \sup_{f \in \mathcal{P}_{d, L, U}}
%         \left\{
%         (\log(\rho^{-1/2}))^{-1}\cdot \min\left\{ \frac{L}{E_\mathcal{Q}},  \log\left(\frac{\aup}{\adown}\right)\right\}
%         \right\}
%         \right\}
%         \in(0, \infty) .
%     \end{equation}
%     %As $w>0$ is defined as \eqref{eq:w},
%     %\begin{equation}
%     %    w
%     %    \leq 
%     %    \frac{L}{dU}
%     %    \cdot
%     %    B_{\inf}^\mathrm{high}(q^\mathrm{high})
%     %    \cdot
%     %    \sqrt{\frac{2}{\pi}}
%     %    \cdot Q
%     %    \to 0
%     %     \text{ as } d\to\infty
%     %  ,
%     %\end{equation}
%     %for any sequence of $f$ in $\{f\in\mathcal{P}_{d, L, U}\}_{d\geq 3}$.
%     \todo{WIP below}
%     \begin{multline}
%         0>
%         \\
%         -\frac12
%         \log\left(
%         \max\left\{
%             \exp\left(
%                 \eta\cdot\min\left\{\frac{w}{4}, \log\left(\frac{\aup}{\adown}\right)\right\}
%             \right)
%             \cdot (1-q^\mathrm{high})
%             + q^\mathrm{high}
%             ,
%             \right.
%             \right.
%             \\
%             \left.
%             \left.
%             \exp\left(
%                 \eta\cdot\min\left\{\frac{w}{4}, \log\left(\frac{\aup}{\adown}\right)\right\}
%             \right)
%             \cdot q^\mathrm{low}
%             + 1 - q^\mathrm{low}
%             \right\}
%             \right)
%         \\
%         >
%          -\frac12
%         \log\left(
%             \exp\left(
%                 \eta\cdot\min\left\{\frac{w}{4}, \log\left(\frac{\aup}{\adown}\right)\right\}
%             \right)
%             + 1
%             \right) 
%         \\
%         \geq
%          -\frac12
%             \exp\left(
%                 \eta\cdot\min\left\{\frac{w}{4}, \log\left(\frac{\aup}{\adown}\right)\right\}
%             \right)
%         .
%     \end{multline}
%     for any sequence of $f$ in $\{f\in\mathcal{P}_{d, L, U}\}_{d\geq 3}$.
%     Moreover, 
%     \begin{equation}
%         \frac{
%             \log\left(
%             \left(\exp\left(-\eta\cdot p_\mathrm{target}\cdot\min\left\{\frac{w}{4}, \log\left(\frac{\aup}{\adown}\right)\right\}\right)
%             \right)^{-1/2}
%             \right)
%         }{
%             \min\left\{\frac{w}{4}, \log\left(\frac{\aup}{\adown}\right)\right\}
%         }
%         =
%         \frac{\eta\cdot p_\mathrm{target}}{2}
%         >0
%         \ensapce.
%     \end{equation}
%     Therefore,
%     \begin{equation}
%         \lim_{d \to \infty}
%         \left\{
%         \sup_{f \in \mathcal{P}_{d, L, U}}
%         \left\{
%         \frac{ \log(\rho^{-1/2})}{ \min\left\{ \frac{L}{E_\mathcal{Q}},  \log\left(\frac{\aup}{\adown}\right)\right\} }
%         \right\}
%         \right\}
%         =
%         \frac{\eta\cdot p_\mathrm{target}}{2}
%         \in (0, \infty)
%         .
%     \end{equation}

%     
% \end{proof}

%        \begin{lemma}    
%        Let   
%        \begin{align} % Q^\mathrm{high} &= \inf \left\{ Q : B_\theta^\mathrm{low}(Q) \leq \sqrt{\frac{L}{U}}B_\theta^\mathrm{high}(q^\mathrm{high}) \right\} % \\
%        Q &:= 
%        \underset{\theta\in\Theta_r\setminus\{\theta|m\neq x^*\}}{\inf}
%        \left\{
%        \sup \left\{ Q_\theta : B_\theta^\mathrm{high}(Q_\theta) \geq B_\theta^\mathrm{low}(q^\mathrm{low}) \right\} 
%        \right\} .
%              \label{eq:qh}
%      \end{align}
%      Then, 
%        \begin{multline}
%          \E \left[\log\left( \frac{ f(m + \sigma z) }{ f(m) } \right) \cdot\ind{f(m + \sigma z) \leq f(m)}\right]
%          \\
%          \leq
%          \frac{ L \cdot  B_\theta^\mathrm{high}(q^\mathrm{high}) }{ 2 \E[\mathcal{Q}_z(\theta)] }  \left( 
%          \frac{B_\theta^\mathrm{low}(q^\mathrm{low})}{\kappa_{\inf}} - \sqrt{\frac{2}{\pi}} \right) \cdot Q
%          < 0
%          .
%          \label{eq:lemma:reasonable:qualitygain}
%        \end{multline}
%        \end{lemma}
%
%Under the condition $B_\theta^\mathrm{high}(q^\mathrm{high}) \cdot \sqrt{ 2 L f(m) } / \E[\mathcal{Q}_z(\theta)] \leq \sigma$, we have
%  \begin{align*}
%    \frac{ \sigma \norm{\nabla f(m)} }{ f(m) } \geq B_\theta^\mathrm{high}(q^\mathrm{high}) \frac{ \sqrt{ 2 L } }{ \E[\mathcal{Q}_z(\theta)] } \frac{ \norm{\nabla f(m)} }{ \sqrt{ f(m) } }
%    .
%  \end{align*}
%  Using the relation $\sqrt{f(m)} \leq \norm{\nabla f(m)} / \sqrt{2 L}$ on $L$-storong convexity, we obtain
%  \begin{align*}
%    \frac{ \sigma \norm{\nabla f(m)} }{ f(m) } \geq \frac{ 2 L \cdot B_\theta^\mathrm{high}(q^\mathrm{high}) }{ \E[\mathcal{Q}_z(\theta)] }
%    .
%  \end{align*}
%  Under the condition $\sigma \leq B_\theta^\mathrm{low}(q^\mathrm{low}) \cdot \norm{\nabla f(m)} / \E[\mathcal{Q}_z(\theta)]$, we have
%  \begin{align*}
%    \frac{ \sigma \E[\mathcal{Q}_z(\theta)] }{ \norm{\nabla f(m)} } \leq B_\theta^\mathrm{low}(q^\mathrm{low}).
%  \end{align*}
%  From \Cref{lemma:qualitygain}, we obtain \eqref{eq:lemma:reasonable:qualitygain}.
%  The negativity follows $B_\theta^\mathrm{low}(q^\mathrm{low}) <\kappa_{\inf}\cdot\sqrt{\frac{2}{\pi}}$.


\subsection{Lower Convergence Rate Bound}

The lower bound of the lower convergence rate can also be derived.

\begin{theorem}[Lower Convergence Rate Bound]\label{thm:lowerconvergencerate}
Let $\{(\tilde{\theta}_t, \mathcal{F}_t)\}_{t \geq 0} = \texttt{ES}(h, \tilde{\theta}_0, \{z_t\}_{t \geq 0})$ be the state sequence of the (1+1)-ES solving $h \in \mathcal{P}_{d,L,U}^\gamma$ with $d > 3$. Suppose that $\tilde{\theta}_0 \in \Theta\setminus\{\tilde{\theta}\mid \tilde{m}= x^\mathrm{opt}\}$.
Then, the lower convergence rate on $h$ satisfies $\exp(-\mathrm{CR}^\mathrm{lower}) \geq \exp(-1/d)$.
\end{theorem}
\begin{proof}
We prove it by using \Cref{prop:f-norm,prop:convergence_rate}. 
Continuing from the proof of \Cref{theorem:lower}, we consider $f$ corresponding to $h$. Let $\{\theta_t\}$ be defined in the proof of \Cref{theorem:lower}.

We let $X_t = \log(\norm{m_t})$ in \Cref{prop:convergence_rate}. 
Lemma 4.6 of \cite{akimoto2020global} shows that
\begin{multline}
    \E[\log(\norm{m_{t+1}}/\norm{m_t}) \mid \mathcal{F}_t ] \\
    \geq 
    \E[\min\{ \log(\norm{m_{t+1}}/\norm{m_t}), 0 \} \mid \mathcal{F}_t ] \geq - 1/d
\end{multline}
for an arbitrary measurable objective function $h$ with $d \geq 2$. 
This shows condition C2 in \Cref{prop:convergence_rate} with $B^\mathrm{lower} = 1/d$. Hence, it suffices to show condition C3.

Under \Cref{asm:L_U_2diff}, we have $\log(2/U) \leq \log(\norm{x}^2) - \log(f(x)) \leq \log(2/L)$. Therefore,
\begin{multline}
    \frac{1}{2}\log\left(\frac{L}{U}\right)
    \leq
    \log\left(\frac{\norm{m_{t+1}}}{\norm{m_t}}\right)
    - \frac{1}{2} \log\left(\frac{f(m_{t+1})}{f(m_t)}\right)
    \\
    \leq
    \frac{1}{2}\log\left(\frac{U}{L}\right).
\end{multline}
\Cref{lemma:variancebound} along with {$x^2 \leq 2 (\exp\abs{x} - 1)$} and the above inequality implies
\begin{equation}
\begin{aligned}[t]
&{\color{blue}\Var\left[\log\left(\frac{\norm{m_{t+1}}}{\norm{m_t}}\right)\right]}
\leq
{\color{blue}\E\left[\left(\log\left(\frac{\norm{m_{t+1}}}{\norm{m_t}}\right)\right)^2\right] }
\\
&\leq \frac{1}{2}{\color{blue}\E\left[\left(\log\left(\frac{f(m_{t+1})}{f(m_t)}\right)\right)^2\right]} + \frac{1}{2}\left(\log\left(\frac{U}{L}\right)\right)^2
\\
&\leq
\left(\frac{U}{L}\cdot\left(1+\frac{1}{d-3}\right) -1\right) + \frac{1}{2}\left(\log\left(\frac{U}{L}\right)\right)^2
<\infty
.
\end{aligned}
\end{equation}
This shows condition C3 in \Cref{prop:convergence_rate}.
This completes the proof.
\end{proof}


\begin{figure}[t]
    \centering
    \begin{subfigure}{0.33\hsize}%
    \includegraphics[width=\hsize, trim=10 15 10 10, clip]{fig/estype0_func2.pdf}%
    \\
    \includegraphics[width=\hsize, trim=10 15 10 10, clip]{fig/estype1_func2.pdf}%
    \\
    \includegraphics[width=\hsize, trim=10 15 10 10, clip]{fig/estype2_func2.pdf}%
    \caption{$H_1$}\label{fig:cigar}%
    \end{subfigure}%
    \begin{subfigure}{0.33\hsize}%
    \includegraphics[width=\hsize, trim=10 15 10 10, clip]{fig/estype0_func1.pdf}%
    \\
    \includegraphics[width=\hsize, trim=10 15 10 10, clip]{fig/estype1_func1.pdf}%
    \\
    \includegraphics[width=\hsize, trim=10 15 10 10, clip]{fig/estype2_func1.pdf}%
    \caption{$H_2$}\label{fig:ellipsoid}%
    \end{subfigure}%
    \begin{subfigure}{0.33\hsize}%
    \includegraphics[width=\hsize, trim=10 15 10 10, clip]{fig/estype0_func3.pdf}%
    \\
    \includegraphics[width=\hsize, trim=10 15 10 10, clip]{fig/estype1_func3.pdf}%
    \\
    \includegraphics[width=\hsize, trim=10 15 10 10, clip]{fig/estype2_func3.pdf}%
    \caption{$H_3$}\label{fig:discus}%
    \end{subfigure}%
    \caption{Results on convex quadratic functions with Hessian matrix $H_1$, $H_2$, and $H_3$. Average and standard error of $\widehat{\mathrm{CR}} \times (\Tr(H) / L)$ with varying $d$ and varying $\kappa$: ($\star$) $0$, ($\times$) $1$, ($\circ$) $2$, ($\triangle$) $3$, ($\triangledown$) $4$, ($\lhd$) $5$, ($\rhd$) $6$. Top: $\aup = \exp(1)$. Middle: $\aup = \exp(1/\sqrt{d})$. Bottom: $\aup = \exp(1/d)$.}
    \label{fig:all}
\end{figure}

\begin{figure}[t]
    \centering
    \begin{subfigure}{0.33\hsize}%
    \includegraphics[width=\hsize, trim=10 15 10 10, clip]{fig/estype0_func2_woscale.pdf}%
    \\
    \includegraphics[width=\hsize, trim=10 15 10 10, clip]{fig/estype1_func2_woscale.pdf}%
    \\
    \includegraphics[width=\hsize, trim=10 15 10 10, clip]{fig/estype2_func2_woscale.pdf}%
    \caption{$H_1$}\label{fig:cigar_woscale}%
    \end{subfigure}%
    \begin{subfigure}{0.33\hsize}%
    \includegraphics[width=\hsize, trim=10 15 10 10, clip]{fig/estype0_func1_woscale.pdf}%
    \\
    \includegraphics[width=\hsize, trim=10 15 10 10, clip]{fig/estype1_func1_woscale.pdf}%
    \\
    \includegraphics[width=\hsize, trim=10 15 10 10, clip]{fig/estype2_func1_woscale.pdf}%
    \caption{$H_2$}\label{fig:ellipsoid_woscale}%
    \end{subfigure}%
    \begin{subfigure}{0.33\hsize}%
    \includegraphics[width=\hsize, trim=10 15 10 10, clip]{fig/estype0_func3_woscale.pdf}%
    \\
    \includegraphics[width=\hsize, trim=10 15 10 10, clip]{fig/estype1_func3_woscale.pdf}%
    \\
    \includegraphics[width=\hsize, trim=10 15 10 10, clip]{fig/estype2_func3_woscale.pdf}%
    \caption{$H_3$}\label{fig:discus_woscale}%
    \end{subfigure}%
    \caption{Results on convex quadratic functions with Hessian matrix $H_1$, $H_2$, and $H_3$. Average and standard error of $\widehat{\mathrm{CR}}$ with varying $d$ and varying $\kappa$: ($\star$) $0$, ($\times$) $1$, ($\circ$) $2$, ($\triangle$) $3$, ($\triangledown$) $4$, ($\lhd$) $5$, ($\rhd$) $6$. Top: $\aup = \exp(1)$. Middle: $\aup = \exp(1/\sqrt{d})$. Bottom: $\aup = \exp(1/d)$.}
    \label{fig:all_woscale}
\end{figure}

\subsection{Numerical Evaluation}\label{sec:exp}

We conducted numerical experiments to estimate the convergence rate of the (1+1)-ES with success-based step-size adaptation on different convex quadratic functions with different hyper parameter value $\aup$. We set $\adown = \aup^{-1/4}$ and $\aup = \exp(1)$, $\exp(1/\sqrt{d})$ and $\exp(1/d)$. 

We consider three convex quadratic functions $f(x) = \frac{1}{2} x^\mathrm{T} H x$ with Hessian matrices defined as 
\begin{align}
    H_1 &= \diag\left(1, 10^{\kappa}, \dots, 10^{\kappa}\right),
    \\
    H_2 &= \diag\left(10^{\kappa\cdot\frac{0}{d-1}}, 10^{\kappa\cdot\frac{1}{d-1}}, \dots, 10^{\kappa\cdot\frac{d-1}{d-1}}\right),
    \\
    H_3 &= \diag\left(1, \dots, 1, 10^{\kappa}\right).
\end{align}
In all cases, the greatest and smallest eigenvalues of the Hessian matrices are $U = 10^{\kappa}$ and $L = 1$, while the trace of the Hessian is $\Tr(H) = 1 + (d-1)10^{\kappa}$, $\sum_{i=1}^{d} 10^{\kappa \cdot \frac{i-1}{d-1}}$, and $d-1 + 10^{\kappa}$ for $H_1$, $H_2$ and $H_3$, respectively. Note that they all reduce to $H_0 = I$ when $\kappa = 0$.

The convergence rate is estimated by running the (1+1)-ES for $T = 10000 + 1000 \cdot d$ iterations. We estimate $\mathrm{CR}$ in \Cref{def:convergence} as $\widehat{\mathrm{CR}}$ by the least square linear regression of $\log(m_t)$ over $t \in \llbracket 0.9 T +1, T\rrbracket$. 
To avoid the numerical error, if $f(m_{\tilde{t}}) < 10^{-100}$ is reached at iteration $\tilde{t}$, we stop the run and set $T = \tilde{t}$ when estimating the convergence rate.
The initial solution $m_0$ is drawn from $\mathcal{N}(0, I)$ for each trial and the initial step-size is set to $\sigma_0 = \norm{\nabla f(m_0)} / \Tr(H)$. 
We run 10 independent trials.

\Cref{fig:all,fig:all_woscale} show the average and the standard deviation of $\widehat{\mathrm{CR}}$ on functions with $H_1$, $H_2$, and $H_3$, respectively, with condition number $10^{\kappa}$ for $\kappa = 0, \dots, 6$ on dimension $d = 1, 3, 10, 30, 100, 300, 1000, 3000, 10000$. 
Particularly, \Cref{cor:upper} implies $\mathrm{CR} \in \Omega(L / Tr(H))$ as $d \to \infty$ and \Cref{thm:lowerconvergencerate} implies $\mathrm{CR} \leq 1/d$. 

First, we observe that $\widehat{\mathrm{CR}}$ values observed for different $\aup$ values are very similar. 
Although it takes $d$ times more iterations to increase or decrease the step-size for the same amount when $\aup = \exp(1/d)$ than when $\aup = \exp(1)$, the difference in $\widehat{\mathrm{CR}}$ was less than the factor of $2$. However, as \Cref{theorem:lower} states, $\mathrm{CR}$ will be decreased as $d$ increases if we set $\log(\aup) \in o(1/d)$.  

Second, we observe the greatest $\widehat{\mathrm{CR}}$ when $\kappa = 0$ for all Hessian settings. We observe $\widehat{\mathrm{CR}} \approx 0.1 / d$. Therefore, the result of \Cref{thm:lowerconvergencerate} is tight up to a constant factor of approximately $0.1$. 

Third, we observe that $\widehat{\mathrm{CR}} \times (\Tr(H) / L) \geq 0.1$ for all cases.
This implies the implication of \Cref{cor:upper} (i.e., $\mathrm{CR} \in \Omega(L/\Tr(H))$) is valid for small $d$ although it is derived for the limit of $d$ to infinity.

Fourth, we observe that $\widehat{\mathrm{CR}} \times (\Tr(H) / L) \in [0.1, 2.0]$ on $H_1$ and $H_3$ for all $\kappa$ and for all $d$. This implies $\mathrm{CR} \in \Omega(L/\Tr(H))$ is a tight bound up to a constant factor on $H_1$ and $H_3$. 
However, we observe that $\widehat{\mathrm{CR}} \times (\Tr(H) / L) \in [0.1, 2.0]$ on $H_2$ for $\kappa \leq 3$ and for all $d$, but $\widehat{\mathrm{CR}} \times (\Tr(H) / L)$ becomes greater as $\kappa$ increases for $\kappa \geq 4$. We conjecture that it is because the (1+1)-ES did not reach the stationary regime within the given budget on $H_2$ with $\kappa \geq 4$ and $\mathrm{CR}$ was overestimated. Thus, we observed a significantly greater $\widehat{\mathrm{CR}}$ if we reduce $T$ by the factor of $10$.

Fifth, we observe that $\widehat{\mathrm{CR}}$ scales differently with respect to $d$ on $H_1$ and $H_3$. If we consider the naive bound derived in \Cref{cor:upper} (i.e., $\mathrm{CR} \in \Omega(L/(d\cdot U)))$), the bounds for these two problems are the same. However, if we consider the bound specialized for convex quadratic function derived in \Cref{cor:upper} (i.e., $\mathrm{CR} \in \Omega(L/\Tr(H))$), we find that the bounds are $1 / (1 + (d-1)10^{\kappa})$ for $H_1$ and $1 / (d-1 + 10^{\kappa})$ for $H_3$. The bound for $H_1$ is well approximated by the naive bound. The bound for $H_3$ is approximated by $1/10^{\kappa}$ if $d < 10^{\kappa}$ and $1/d$ if $d > 10^{\kappa}$. This clearly appears in \Cref{fig:discus_woscale}. Namely, the convergence rate of the (1+1)-ES does not always rely on the ratio $L/(d\cdot U)$ as in the naive bound. However, it can be as good as $1/d$ when the number of axes sensitive to the objective function (corresponding to the eigen vectors for large eigenvalues of the Hessian matrix) is limited.

% \begin{figure}
%     \centering
%     \begin{subfigure}{0.33\hsize}%
%     \includegraphics[width=\hsize, trim=10 15 10 10, clip]{fig/estype0_func2.pdf}%
%     \\
%     \includegraphics[width=\hsize, trim=10 15 10 10, clip]{fig/estype0_func2_woscale.pdf}%
%     \caption{$\aup=\exp(1)$}\label{fig:es0f2w}%
%     \end{subfigure}%
%     \begin{subfigure}{0.33\hsize}%
%     \includegraphics[width=\hsize, trim=10 15 10 10, clip]{fig/estype1_func2.pdf}%
%     \\
%     \includegraphics[width=\hsize, trim=10 15 10 10, clip]{fig/estype1_func2_woscale.pdf}%
%     \caption{$\aup=\exp(1/\sqrt{d})$}\label{fig:es1f2w}%
%     \end{subfigure}%
%     \begin{subfigure}{0.33\hsize}%
%     \includegraphics[width=\hsize, trim=10 15 10 10, clip]{fig/estype2_func2.pdf}%
%     \\
%     \includegraphics[width=\hsize, trim=10 15 10 10, clip]{fig/estype2_func2_woscale.pdf}%
%     \caption{$\aup=\exp(1/d)$}\label{fig:es2f2w}%
%     \end{subfigure}%
%     \caption{Results on convex quadratic function with Hessian matrix $H_1$. Average and standard error of $\widehat{\mathrm{CR}} \times (\Tr(H) / L)$ (top) and $\widehat{\mathrm{CR}}$ (bottom) with different $d$ and different $\kappa$: ($\star$) $0$, ($\times$) $1$, ($\circ$) $2$, ($\triangle$) $3$, ($\triangledown$) $4$, ($\lhd$) $5$, ($\rhd$) $6$.}
%     \label{fig:cigar}
% \end{figure}

% \begin{figure}
%     \centering
%     \begin{subfigure}{0.33\hsize}%
%     \includegraphics[width=\hsize, trim=10 15 10 10, clip]{fig/estype0_func1.pdf}%
%     \\
%     \includegraphics[width=\hsize, trim=10 15 10 10, clip]{fig/estype0_func1_woscale.pdf}%
%     \caption{$\aup=\exp(1)$}\label{fig:es0f1w}%
%     \end{subfigure}%
%     \begin{subfigure}{0.33\hsize}%
%     \includegraphics[width=\hsize, trim=10 15 10 10, clip]{fig/estype1_func1.pdf}%
%     \\
%     \includegraphics[width=\hsize, trim=10 15 10 10, clip]{fig/estype1_func1_woscale.pdf}%
%     \caption{$\aup=\exp(1/\sqrt{d})$}\label{fig:es1f1w}%
%     \end{subfigure}%
%     \begin{subfigure}{0.33\hsize}%
%     \includegraphics[width=\hsize, trim=10 15 10 10, clip]{fig/estype2_func1.pdf}%
%     \\
%     \includegraphics[width=\hsize, trim=10 15 10 10, clip]{fig/estype2_func1_woscale.pdf}%
%     \caption{$\aup=\exp(1/d)$}\label{fig:es2f1w}%
%     \end{subfigure}%
%     \caption{Results on convex quadratic function with Hessian matrix $H_2$. Average and standard error of $\widehat{\mathrm{CR}} \times (\Tr(H) / L)$ (top) and $\widehat{\mathrm{CR}}$ (bottom) with different $d$ and different $\kappa$: ($\star$) $0$, ($\times$) $1$, ($\circ$) $2$, ($\triangle$) $3$, ($\triangledown$) $4$, ($\lhd$) $5$, ($\rhd$) $6$.}
%     \label{fig:ellipsoid}
% \end{figure}

% \begin{figure}
%     \centering
%     \begin{subfigure}{0.33\hsize}%
%     \includegraphics[width=\hsize, trim=10 15 10 10, clip]{fig/estype0_func3.pdf}%
%     \\
%     \includegraphics[width=\hsize, trim=10 15 10 10, clip]{fig/estype0_func3_woscale.pdf}%
%     \caption{$\aup=\exp(1)$}\label{fig:es0f3w}%
%     \end{subfigure}%
%     \begin{subfigure}{0.33\hsize}%
%     \includegraphics[width=\hsize, trim=10 15 10 10, clip]{fig/estype1_func3.pdf}%
%     \\
%     \includegraphics[width=\hsize, trim=10 15 10 10, clip]{fig/estype1_func3_woscale.pdf}%
%     \caption{$\aup=\exp(1/\sqrt{d})$}\label{fig:es1f3w}%
%     \end{subfigure}%
%     \begin{subfigure}{0.33\hsize}%
%     \includegraphics[width=\hsize, trim=10 15 10 10, clip]{fig/estype2_func3.pdf}%
%     \\
%     \includegraphics[width=\hsize]{fig/estype2_func3_woscale.pdf}%
%     \caption{$\aup=\exp(1/d)$}\label{fig:es2f3w}%
%     \end{subfigure}%
%     \caption{Results on convex quadratic function with Hessian matrix $H_3$. Average and standard error of $\widehat{\mathrm{CR}} \times (\Tr(H) / L)$ (top) and $\widehat{\mathrm{CR}}$ (bottom) with different $d$ and different $\kappa$: ($\star$) $0$, ($\times$) $1$, ($\circ$) $2$, ($\triangle$) $3$, ($\triangledown$) $4$, ($\lhd$) $5$, ($\rhd$) $6$.}
%     \label{fig:discus}
% \end{figure}



\section{Discussion}

We derived the convergence rate bounds for the (1+1)-ES with success-based step-size adaptation on functions that are Lipschitz smooth and strongly convex around the global optimum and their strictly increasing transformation (i.e., $\mathcal{P}_{d, L, U}^{\gamma}$ defined in \Cref{def:problem}). To the best of our knowledge, this is the first study that shows the convergence rate dependency of the form $\frac{L}{d U}$ on dimension $d$, Lipschitz smoothness parameter $U$, and strong convexity parameter $L$ explicitly on $\mathcal{P}_{d, L, U}^{\gamma}$. The tighter convergence rate bounds of form $\frac{L}{\Tr(H)}$ on convex quadratic functions $\frac12 x^\T H x$ derived in this study provide further insight into the algorithmic behavior. In particular, the fact that the convergence rate can be independent of the condition number $\Cond(H) {\leq} U / L$ on functions where a few eigenvalues of $H$ are large and the others are all relatively small is a promising characteristic of the (1+1)-ES when applied to high dimensional problems with a few important variables.

Although our results are stated as asymptotic ones because our focus is to derive the convergence rate, we can easily derive a non-asymptotic bound if we limit our attention to Lipschitz smooth and strongly convex functions (i.e., $\mathcal{S}_{d,L,U}$). Noting that our potential function $V(\theta)$ satisfies $\log(f(m)) \leq V(\theta)$, we have
\begin{equation}
\begin{split}
    \MoveEqLeft[0]\E\left[\log\left( f(m_t) \right)\right] - \log\left(f(m_0)\right)\\
    &\leq \E\left[\log\left( V(\theta_t) \right)\right] - \log\left( V(\theta_0) \right) + \log\left( V(\theta_0) \right) - \log\left( f(m_0) \right)
    \\
    &\leq - 2\cdot B_{d, L, U}^\mathrm{upper} \cdot t +  \log\left( V(\theta_0) \right) - \log\left( f(m_0) \right) .
\end{split}
\end{equation}

Three strong points of the (1+1)-ES with success-based step-size adaptation over other derivative-free approaches analyzed on $\mathcal{S}_{d,L,U}$ are as follows. 
First, any hyper-parameters need not be set depending on some characteristic parameters such as $L$ and $U$.
Without the use of these parameters inside the algorithm, the convergence rate factor of $\frac{L}{d U}$ is achieved. 
As estimating $L$ and $U$ in simulation-based optimization settings is not trivial, it is a strong advantage of the (1+1)-ES.
Second, the (1+1)-ES is invariant to any strictly increasing transformation of the objective function; hence, our analysis is not limited to $\mathcal{S}_{d,L,U}$, but covers a broader class of functions, namely $\mathcal{P}_{d, L, U}^{\gamma}$.%$\mathcal{S}_{d,L,U}$.
Third, our results guarantee linear convergence of the (1+1)-ES, rather than locating an $\epsilon$-optimal solution such that $f(x) - f(x^*) \leq \epsilon$ for a finite $\epsilon$. 
Previous studies on derivative-free approaches \cite{2019gradientlesgolovins, nesterov2017random} set some hyper-parameters of the algorithms depending on this pre-defined threshold $\epsilon$ to derive these results. 


We close our analysis with discussion on two limitations of this study. 
First, we have derived the convergence rate bound in \Cref{theorem:lower}.
However, the dependency of $B_{d, L, U}^\mathrm{upper}$ on $d$, $L$, and $U$ are not explicitly derived for finite $d$. 
Second, to evaluate $B_{d, L, U}^\mathrm{upper}$ in the limit $d \to \infty$, we have assumed that $\aup$ and $\adown$ are set such that $\Phi\left(-\sqrt{\frac{2}{\pi}}\right) < p^\mathrm{target} < \frac{1}{2}$, where $p^\mathrm{target} = \frac{\log(1/\adown)}{\log(\aup / \adown)}$. 
However, as $\Phi\left(-\sqrt{\frac{2}{\pi}}\right)\approx 0.212$, the most common setting of $\aup$ and $\adown$ is such that $p^\mathrm{target} = 1/5$ \cite{rechenberg1973evolution} and is not included in the above assumption.
These limitations will be addressed in future works.


\section*{Acknowledgments}
%This should be a simple paragraph before the References to thank those individuals and institutions who have supported your work on this article.
This paper is partially supported by JSPS KAKENHI Grant Number 19H04179.



\appendix

\section{Proof}

\subsection{Proof of \Cref{prop:invariance}}\label{apdx:prop:invariance}

For the first claim, it is sufficient to show that $\mathcal{G}(\theta, z; f) = \mathcal{G}(\theta, z; g \circ f)$ for all $\theta \in \Theta$ and $z \in \R^d$. This is trivial because $f(m + \sigma \cdot z) \leq f(m) \Leftrightarrow g(f(m + \sigma \cdot z)) \leq g(f(m))$.

For the second claim, we first show that $\mathcal{G}((m, \sigma), z; f) = \mathcal{G}((m+x^*, \sigma), z; f \circ T)$ for all $\theta \in \Theta$ and $z \in \R^d$. Because $f(T(m + x^* + \sigma \cdot z)) \leq f(T(m + x^*)) \Leftrightarrow f(m + \sigma \cdot z) \leq f(m)$, it is clear that $\mathcal{G}((m, \sigma), z; f) = \mathcal{G}((m+x^*, \sigma), z; f \circ T)$. Assume that $\tilde{m}_{t} = m_t + x^*$ and $\tilde{\sigma}_t = \sigma_t$. Then, we have $\mathcal{G}(\theta_t, z; f) = \mathcal{G}(\tilde{\theta}_t, z; f \circ T)$, and hence $\tilde{m}_{t+1} = m_{t+1} + x^*$ and $\tilde{\sigma}_{t+1} = \sigma_{t+1}$. Because the assumption holds for $t = 0$, by mathematical induction, we obtain the second claim.

For the third claim, it is sufficient to show that $\mathcal{G}(\theta, z; f) = \mathcal{G}(\theta, z; \tilde{f})$ for all $\theta \in \Theta$ and $z \in \R^d$. Because the (1+1)-ES guarantees the monotone improvement of $f(m_t)$, we have $m_t \in S(m_0)$. Therefore, $\tilde{f}(m_t) = f(m_t)$ for all $t \geq 0$. We have $f(m_t + \sigma_t \cdot z_t) \leq f(m_t) \Rightarrow \tilde{f}(m_t + \sigma_t \cdot z_t) \leq \tilde{f}(m_t)$. Moreover, because $\tilde{f}(x) > \tilde{f}(m_0)$ for all $x \in \R^d\setminus S(m_0)$, we have $\tilde{f}(m_t + \sigma_t \cdot z_t) \leq \tilde{f}(m_t) \Rightarrow m_t + \sigma_t \cdot z_t \in S(m_0)$, implying $f(m_t + \sigma_t \cdot z_t) \leq f(m_t)$. Hence, $f(m_t + \sigma_t \cdot z_t) \leq f(m_t) \Leftrightarrow \tilde{f}(m_t + \sigma_t \cdot z_t) \leq \tilde{f}(m_t)$ and we obtain $\mathcal{G}(\theta, z; f) = \mathcal{G}(\theta, z; \tilde{f})$.

\subsection{Proof of \Cref{prop:convergence_rate}}\label{apdx:prop:convergence_rate}
Let $Z_{t+1} = X_{t+1} - \E[X_{t+1} \mid \mathcal{F}_t]$.
Then, $\{Z_{t}\}_{t \geq 1}$ is a martingale difference sequence adapted to $\{\mathcal{F}_{t}\}$, and
\begin{multline}
\frac1t (X_t - X_0)
= \frac1t \sum_{i=1}^{t}(X_i - X_{i-1})
= \frac1t \sum_{i=1}^{t}(Z_i + \E[X_i - X_{i-1} \mid \mathcal{F}_{i-1}])
\\
= \frac1t \sum_{i=1}^{t}Z_i + \frac1t \sum_{i=1}^{t} \E[X_i - X_{i-1} \mid \mathcal{F}_{i-1}]
.
\end{multline}
Suppose that C3 holds.
{\color{blue}Beucause $\Var[X_t] = \Var[Z_t] + \Var[\E[X_t\mid\mathcal{F}_{t-1}]] \geq \Var[Z_t]$, }
C3 immediately implies that $\sum_{t=1}^{\infty}{\color{blue}\E[Z_{t}^2]} / t^2 < \infty$. Consequently, from the strong law of large numbers of martingale \cite{chow1967strong}, we obtain $\lim_{t\to\infty} \frac1t \sum_{i=1}^{t} Z_t = 0$ almost surely. 
We obtain statement (I) by taking $\limsup$ of the equation above and using C1 and $\limsup_{t\to\infty} \frac1t \sum_{i=1}^{t} Z_t = 0$.
Similarly, we obtain statement (II) by taking $\liminf$ and using C2 and $\liminf_{t\to\infty} \frac1t \sum_{i=1}^{t} Z_t = 0$. This completes the proof.

\subsection{Proof of \Cref{prop:f-norm}}\label{apdx:prop:f-norm}
Let $T: x \mapsto x - x^\text{opt}$. Then, $T^{-1}(x) = x + x^\text{opt}$ and $S_T$ in (2) of \Cref{prop:invariance} is equal to $S$ in \Cref{prop:f-norm}.
Hence, in light of (1) and (2) of \Cref{prop:invariance}, we have $\tilde{\theta}_t = S(\theta_t)$ for all $t \geq 0$.
This implies $\norm{\tilde{m}_t - x^\mathrm{opt}} = \norm{m_t}$ for all $t \geq 0$. Hence, we obtain \eqref{eq:liminf}. 
Moreover, under \Cref{asm:L_U_2diff}, for any $x \in \R^d$, we have
$\frac{2 f(x)}{ U } \leq \norm{x}^2 \leq \frac{2 f(x)}{ L }$.
Then, we obtain
$\frac{1}{2}\log \left(\frac{L}{U}\right)
\leq \log\left(\frac{ \norm{m_t}} { \norm{m_0} }\right) - \frac{1}{2}\log\left(\frac{f(m_t)}{f(m_0)}\right)
\leq \frac{1}{2}\log \left(\frac{U}{L}\right)$.
% \begin{equation}
% \frac{1}{2}\log \left(\frac{L}{U}\right)
% \leq \log\left(\frac{ \norm{m_t}} { \norm{m_0} }\right) - \frac{1}{2}\log\left(\frac{f(m_t)}{f(m_0)}\right)
% \leq \frac{1}{2}\log \left(\frac{U}{L}\right)  .
% \end{equation}
Taking $\limsup$ after multiplying all terms by $\frac{1}{t}$, we obtain \eqref{eq:limsup}.

\subsection{Proof of \Cref{lemma:variance}}\label{apdx:lemma:variance}

%The positivity of $\mathcal{Q}_z(\theta)$ is straight-forward from $L$-strong convexity of $f$.

We prove the statements one by one.

1. The $L$-strong convexity and $U$-Lipschitz smoothness of $f$ implies 
% \begin{equation}
%     L \norm{z}^2 \leq \mathcal{Q}_z(\theta) \leq U \norm{z}^2.
% \end{equation}
$L \norm{z}^2 \leq \mathcal{Q}_z(\theta) \leq U \norm{z}^2$.
This inequality and $\E[\norm{z}^2] = d$ complete the proof for the first statement.

2. \cite[Therem2]{chen1982inequality}
shows that
% \begin{align}
% \Var[\mathcal{Q}_z(\theta)] \leq \E[ \norm{\nabla_z \mathcal{Q}_z(\theta)}^2 ] .
% \end{align}
$\Var[\mathcal{Q}_z(\theta)] \leq \E[ \norm{\nabla_z \mathcal{Q}_z(\theta)}^2 ]$.
Using this fact and the $U$-Lipschitz continuity of $\nabla f$, we have
\begin{equation}
\begin{split}
\Var[\mathcal{Q}_z(\theta)] 
&\leq \E[ \norm{\nabla_z \mathcal{Q}_z(\theta)}^2 ] \\
&= 
%\frac{4}{\sigma^2} 
(4 / \sigma^2) 
\E[ \norm{ \nabla f(m + \sigma z) - \nabla f(m) }^2 ] 
\\
&\leq 
%\frac{4}{\sigma^2} 
(4 / \sigma^2) 
\E[ U^2 \sigma^2 \norm{ z }^2 ] 
= 4 d U^2 .
\end{split}
\end{equation}

3. We use the Cauchy-Schwarz inequality and obtain
\begin{equation}
\begin{split}
\MoveEqLeft[1]
\E[\mathcal{Q}_z(\theta) \cdot \ind{z_e \leq 0}] \\
&= \E[\E[\mathcal{Q}_z(\theta)] \cdot \ind{z_e \leq 0}] + \E[(\mathcal{Q}_z(\theta) - \E[\mathcal{Q}_z(\theta)]) \cdot \ind{z_e \leq 0}]
\\
&= \E[\mathcal{Q}_z(\theta)] / 2 + \E[(\mathcal{Q}_z(\theta) - \E[\mathcal{Q}_z(\theta)]) \cdot \ind{z_e \leq 0}]
\\
&\leq \E[\mathcal{Q}_z(\theta)] / 2 + (\E[(\mathcal{Q}_z(\theta) - \E[\mathcal{Q}_z(\theta)])^2] \cdot \E[\ind{z_e \leq 0}])^{1/2}
\\
&= \E[\mathcal{Q}_z(\theta)] / 2 + (\Var[\mathcal{Q}_z(\theta)] / 2)^{1/2} ,
\end{split}
\end{equation}
and similarly,
\begin{equation}
\begin{split}
\MoveEqLeft[1]
\E[\mathcal{Q}_z(\theta) \cdot \ind{z_e \leq 0}] \\
&= \E[\mathcal{Q}_z(\theta)] / 2 + \E[(\mathcal{Q}_z(\theta) - \E[\mathcal{Q}_z(\theta)]) \cdot \ind{z_e \leq 0}]
\\
&= \E[\mathcal{Q}_z(\theta)] / 2 - \E[(\mathcal{Q}_z(\theta) - \E[\mathcal{Q}_z(\theta)]) \cdot \ind{z_e > 0}]
\\
&\geq \E[\mathcal{Q}_z(\theta)] / 2 - (\E[(\mathcal{Q}_z(\theta) - \E[\mathcal{Q}_z(\theta)])^2] \cdot \E[\ind{z_e > 0}])^{1/2}
\\
&= \E[\mathcal{Q}_z(\theta)] / 2 - (\Var[\mathcal{Q}_z(\theta)] / 2)^{1/2} .
\end{split}
\end{equation}
Note that we have $\E[\mathcal{Q}_z(\theta)] \geq d L$. 
We have,
\begin{equation}
    \begin{split}
(\Var[\mathcal{Q}_z(\theta)]/ 2)^{1/2} 
&\leq (2 d)^{1/2} U 
\\
&= \E[\mathcal{Q}_z(\theta)] ((2 d)^{1/2} U / \E[\mathcal{Q}_z(\theta)] )
\\
&\leq \E[\mathcal{Q}_z(\theta)] ((2 / d)^{1/2} (U / L) ) .
\end{split}
\end{equation}
% \begin{multlined}[t]
% (\Var[\mathcal{Q}_z(\theta)] / 2)^{1/2} 
% \leq (2 d)^{1/2} U \\
% = \E[\mathcal{Q}_z(\theta)] ((2 d)^{1/2} U / \E[\mathcal{Q}_z(\theta)] ) 
% \leq \E[\mathcal{Q}_z(\theta)] ((2 / d)^{1/2} (U / L) ) .
% \end{multlined}
This ends the proof.

% 4. Since $f$ is twice continuously differentiable, we have $\mathcal{Q}_z(\theta) = z^\T H_z(\theta) z$, where
% \begin{equation}
%     H_z(\theta) = 2 \int_{0}^{1}\int_{0}^{t} \nabla^2 f(m + s\cdot \sigma\cdot z) \mathrm{d}s\mathrm{d}t .
% \end{equation}
% For any $\epsilon > 0$ there exist $\delta > 0$ such that the eigenvalues of $\nabla^2 f(\bar{m} + \Delta) - \nabla^2 f(\bar{m})$ are all in $(-\epsilon, \epsilon)$ for any $\Delta$ satisfying $\norm{\Delta} < \delta$. Let $\epsilon$ and $\delta$ be fixed. We have
% \begin{align}
%     &\abs{\E[\mathcal{Q}_z(\theta)] - \Tr(\nabla^2 f(\bar{m}))}\\
%     &= \abs{\E[z^\T H_z(\theta) z] - \E[z^\T \nabla^2 f(\bar{m}) z]} 
%     = \abs{\E[z^\T (H_z(\theta) - \nabla^2 f(\bar{m}) z ]} 
%     \leq \E[(z^\T (H_z(\theta) - \nabla^2 f(\bar{m}) z)^2 ]^{1/2} \\
%     &= \E[(z^\T (H_z(\theta) - \nabla^2 f(\bar{m})) z)^2 \ind{\norm{z} < \delta / (2\sigma)}]^{1/2}
%     + \E[(z^\T (H_z(\theta) - \nabla^2 f(\bar{m})) z)^2 \ind{\norm{z} \geq \delta / (2\sigma)}]^{1/2}
%     .
% \end{align}
% By using $L$-strong convexity and $U$-Lipschitz smoothness of $f$, we have that all the eigenvalues of $H_z(\theta) - \nabla^2 f(\bar{m})$ is bounded in $[L - U, U - L]$. Therefore, the second term on the right-most side (RMS) is bounded as
% \begin{align}
%     \E[(z^\T (H_z(\theta) - \nabla^2 f(\bar{m})) z)^2 \ind{\norm{z} \geq \delta / (2\sigma)}]^{1/2}
%     \leq (U - L) \E[\norm{z}^4 \ind{\norm{z} \geq \delta / (2\sigma)}]^{1/2} .
% \end{align}
% If $\norm{m - \bar{m}} < \delta / 2$, we have $\norm{m + s \cdot \sigma \cdot z - \bar{m}} \leq \norm{m - \bar{m}} + s \cdot \sigma \cdot \norm{z} < \delta$ for any $s \in [0, 1]$ under the condition $\norm{z} < \delta / (2\sigma)$. The, the first term on the RMS is bounded as 
% \begin{align}
%     \E[(z^\T (H_z(\theta) - \nabla^2 f(\bar{m})) z)^2 \ind{\norm{z} < \delta / (2\sigma)}]^{1/2}
%     \leq \epsilon \E[ \norm{z}^4 \ind{\norm{z} < \delta / (2\sigma)}]^{1/2}
%     \leq \epsilon \E[ \norm{z}^4 ]^{1/2} .
% \end{align}
% Note that $C = \E[ \norm{z}^4 ]$ is finite and $\lim_{\sigma \to 0} \E[\norm{z}^4 \ind{\norm{z} \geq \delta / (2\sigma)}] = 0$. Therefore, we have
% \begin{align}
% \limsup_{(m,\sigma) \to (\bar{m}, 0)} \abs{\E[\mathcal{Q}_z(\theta)] - \Tr(\nabla^2 f(\bar{m}))}
% \leq \limsup_{(m,\sigma) \to (\bar{m}, 0)} \E[(z^\T (H_z(\theta) - \nabla^2 f(\bar{m}) z)^2 ]^{1/2}
% \leq C \epsilon .
% \end{align}
% Since $\epsilon > 0$ is arbtrary, we obtain $\lim_{(m,\sigma) \to (\bar{m}, 0)} \E[\mathcal{Q}_z(\theta)] = \Tr(\nabla^2 f(\bar{m}))$.

% 5. Observe
% \begin{align}
%     &\abs{\Var[\mathcal{Q}_z(\theta)^2] - 2\Tr((\nabla^2 f(\bar{m}))^2)}\\
%     &= \abs{\E[\mathcal{Q}_z(\theta)^2] - \E[\mathcal{Q}_z(\theta)]^2 - 2\Tr((\nabla^2 f(\bar{m}))^2)}\\
%     &= \abs{\E[\mathcal{Q}_z(\theta)^2] - \Tr((\nabla^2 f(\bar{m})))^2 - 2\Tr((\nabla^2 f(\bar{m}))^2)}\\
%     &= \abs{\E[(z^\T H_z(\theta) z)^2] - \E[(z^\T \nabla^2 f(\bar{m}) z)^2]} \\
%     &= \abs{\E[ (z^\T H_z(\theta) z - z^\T \nabla^2 f(\bar{m}) z) \cdot (z^\T H_z(\theta) z + z^\T \nabla^2 f(\bar{m}) z)]} \\
%     &\leq \E[ (z^\T H_z(\theta) z - z^\T \nabla^2 f(\bar{m}) z)^2]^{1/2} \cdot \E[(z^\T H_z(\theta) z + z^\T \nabla^2 f(\bar{m}) z)^2]^{1/2} \\
%     &= \E[ (z^\T (H_z(\theta) - \nabla^2 f(\bar{m})) z)^2]^{1/2} \cdot \E[(z^\T (H_z(\theta) + \nabla^2 f(\bar{m})) z)^2]^{1/2} .
% \end{align}
% The second expectation on the RMS is upper bounded by $2U \E[\norm{z}^4]^{1/2}$ because of the $U$-Lipschitz smoothness of $f$. The first expectation on the RMS is shown to converge to $0$ in the proof of 4th statement. Hence we have
% \begin{align}
% \limsup_{(m,\sigma) \to (\bar{m}, 0)} \Var[\mathcal{Q}_z(\theta)] = 2\Tr((\nabla^2 f(\bar{m}))^2).
% \end{align}

% 6. Because of the symmetry of $z^\T \nabla^2 f(\bar{m}) z$ around the origin $z = 0$, it is easy to see that $\E[z^\T \nabla^2 f(\bar{m}) z \cdot \ind{z_e \leq 0}] = \E[z^\T \nabla^2 f(\bar{m}) z] / 2 = \Tr(\nabla^2 f(\bar{m})) / 2$ for any $\bar{m} \in \R^d$ and $m \in \R^d\setminus\{0\}$. Using this fact, we obtain
% \begin{align}
%     &\abs{\E[\mathcal{Q}_z(\theta) \cdot \ind{z_e \leq 0}] - \Tr(\nabla^2 f(\bar{m})) / 2}\\
%     &= \abs{\E[z^\T H_z(\theta) z\cdot \ind{z_e \leq 0}] - \E[z^\T \nabla^2 f(\bar{m}) z \cdot \ind{z_e \leq 0}]} \\
%     &= \abs{\E[( z^\T (H_z(\theta) - \nabla^2 f(\bar{m}) z ) \cdot \ind{z_e \leq 0} ]} \\
%     &\leq \E[(z^\T (H_z(\theta) - \nabla^2 f(\bar{m}) z)^2 ]^{1/2} \E[\ind{z_e \leq 0}]^{1/2} \\
%     &\leq \E[(z^\T (H_z(\theta) - \nabla^2 f(\bar{m}) z)^2 ]^{1/2} (1/2)^{1/2}
%     .
% \end{align}
% The convergence of the RMS of the above inequality to zero is proved in the proof of the 4th statement. This completes the proof.
%\end{proof}
%\hrule

\subsection{Proof of \Cref{cor:tr_cond}}\label{apdx:cor:tr_cond}
In light of \Cref{lemma:variance}, we have for any $\theta \in \Thetazero$
% \begin{equation}
% \abs*{
% \frac{\E[\mathcal{Q}_z(\theta)]}{\E[\mathcal{Q}_z(\theta)\cdot\ind{z_e\leq 0}]} - 2
% }
% \leq \frac{ 1 }{ (2 / d)^{1/2} (U/L) } \xrightarrow{d \to \infty} 0 .
% \end{equation}
\begin{equation}
\abs*{
\frac{\E[\mathcal{Q}_z(\theta)\cdot\ind{z_e\leq 0}]}{\E[\mathcal{Q}_z(\theta)]} - \frac{1}{2}
}
\leq (2 / d)^{1/2} (U/L) \xrightarrow{d \to \infty} 0 .
\end{equation}
It implies that for any $\epsilon > 0$ there exists $D' \in \N$ such that $\kappa_{\inf} \in (2 - \epsilon, 2+\epsilon)$ for all $d \geq D'$. 
Then, the RHS of \eqref{eq:tr_cond} is lower bounded by a positive constant, say $c > 0$, for all $d \geq D'$. 
Again in light of \Cref{lemma:variance}, we have
\begin{equation}
    \mathcal{V}_\mathrm{std}
    \leq 4 d U^2/ (d L)^2 = 4 U^2/ (d L^2 ) \xrightarrow{d \to \infty} 0.
\end{equation}
%\begin{equation}
%    \mathcal{V}_\mathrm{std}
%    \leq \frac{4 d U^2}{ (d L)^2 } = \frac{4 U^2}{ d L^2 } \xrightarrow{d \to \infty} 0.
%\end{equation}
It implies that there exists $D'' \in \N$ such that the LHS of \eqref{eq:tr_cond} is strictly smaller than the above defined $c > 0$ for all $d \geq D''$. 
Therefore, letting $D = \max\{D', D''\}$, we obtain \eqref{eq:tr_cond} for all $d \geq D$. This completes the proof.

\subsection{Proof of \Cref{lemma:qualitygain}}\label{apdx:lemma:qualitygain}

%\begin{proof}%[Proof of \Cref{lemma:qualitygain}]
Let $G_a(z) := (f(m + \sigma z) - f(m) ) \ind{ z_e \leq 0}$ and $G_b(z) := \ind{f(m + \sigma z) \leq f(m)}$. 
First, we prove that the LHS of \eqref{eq:qualitygain} is upper-bounded by $\E[G_a(z)]\E[G_b(z)] / f(m)$.
Because $f: \R^d \to \R$ is a differentiable convex function, we have the following property: 
(A) $\ind{f(m + \sigma z) \leq f(m)} \leq \ind{ z_e \leq 0}$, implying that $\ind{f(m + \sigma z) \leq f(m)} = \ind{ z_e \leq 0}\ind{f(m + \sigma z) \leq f(m)}$;
(B) $G_a(z)$ and $G_b(z)$ are negatively correlated, that is, $(G_a(z_1) - G_a(z_2))(G_b(z_1) - G_b(z_2)) \leq 0$ for all $z_1, z_2 \in \R^d$.
% \begin{enumerate}
% \item[(A)] $\ind{f(m + \sigma z) \leq f(m)} \leq \ind{ z_e \leq 0}$, implying that $\ind{f(m + \sigma z) \leq f(m)} = \ind{ z_e \leq 0}\ind{f(m + \sigma z) \leq f(m)}$;
% \item[(B)] $G_a(z)$ and $G_b(z)$ are negatively correlated, i.e., $(G_a(z_1) - G_a(z_2))(G_b(z_1) - G_b(z_2)) \leq 0$ for all $z_1, z_2 \in \R^d$.
% \end{enumerate}
Because $((f(m + \sigma z) - f(m)) / f(m)) \ind{f(m + \sigma z) \leq f(m)} = G_a(z) G_b(z) / f(m)$, where we used (A), the LHS of \eqref{eq:qualitygain} is  $\E[G_a(z) G_b(z)] / f(m)$.
Using (B) along with the chebyshev sum inequality \cite[Theorem 236]{hardy1952inequalities}, we have $\E[G_a(z)G_b(z)] \leq \E[G_a(z)]\E[G_b(z)]$. Hence, the LHS of \eqref{eq:qualitygain} is upper-bounded by $\E[G_a(z)]\E[G_b(z)] / f(m)$.

Next, we prove that the RHS of \eqref{eq:qualitygain} is $\E[G_a(z)]\E[G_b(z)] / f(m)$. 
First, we have $\E[G_b(z)] = \E[\ind{f(m + \sigma z) \leq f(m)}] = \Pr[f(m + \sigma z) \leq f(m)]$.
Second,
\begin{equation}
\begin{aligned}[t]
&\E[G_a(z)]
%\\
%=&\E[(f(m + \sigma z) - f(m)) \ind{z_e\leq 0}]
\\
=&
  \E\left[\left(\sigma \| \nabla f(m) \| z_e + \frac{\sigma^2 }{2} \mathcal{Q}_z(\theta) \right) \ind{z_e \leq 0}\right]
\\
=& - \frac{\sigma \| \nabla f(m) \|}{\sqrt{2\pi}} + \frac{\sigma^2}{2}\cdot\E\left[\mathcal{Q}_z(\theta) \cdot \ind{z_e \leq 0}\right]
\\
=& \sigma \norm{\nabla f(m)}\cdot\left(
- \frac{1}{\sqrt{2\pi}}
+ \frac{\sigma}{2\norm{\nabla f(m)}}\cdot\E\left[\mathcal{Q}_z(\theta) \cdot \ind{z_e \leq 0}\right]
\right)
%\\
%&\leq \sigma \norm{\nabla f(m)}\cdot\left(
%- \frac{1}{\sqrt{2\pi}}
%+ \frac{\sigma}{2\norm{\nabla f(m)}}\cdot\e\left[z^\mathrm{t} h_z(\theta) z\right]
%\right)
  .
\end{aligned}
\end{equation}
These equalities demonstrate that the RHS of \eqref{eq:qualitygain} is $\E[G_a(z)]\E[G_b(z)] / f(m)$.

%  \mori{the last inequality restricts a scope of our analysis in the view of $p_\mathrm{target}$. is there some way to compute this more accurately?}
%  \yohe{i see the point. so the problem is that we can not use the symmetory that we exploited in quadratic objective cases. does the holder's inequality or more simply schwartz inequality help? for more accurate argument, we might add an additional assumption something like: $\nabla^2 f(x) - \nabla^2 f(m) \preccurlyeq \epsilon$ for all $x$ and $m$. then, $\e\left[z^\mathrm{t} h_z(\theta) z \ind{z_e \leq 0} \right] \leq \tr(\nabla^2 f(m))/2 + d \cdot \epsilon / 2$.}
%  
%  \yohe{we can do a little bit better. consider the taylor expansion of $\log(1 - x) = - \sum_{n=1}^{\infty} \frac{x^n}{n}$. this implies $\log^+(1 - x) \leq - \max(x, 0) - \max(x, 0)^2/2$.
%  \begin{align*}
%  &\e\left[\log^+\left( \frac{ f(m + \sigma z) }{ f(m) } \right) \right]
%  \\
%  &=
%  \e\left[\log^+\left(1 -  \frac{f(m) - f(m + \sigma z) }{ f(m) } \right) \ind{f(m + \sigma z) \leq f(m)}\right]
%  \\
%  &\leq
%  \e\left[\left(-\left(\frac{f(m) - f(m + \sigma z) }{ f(m) } \right) - \frac{1}{2}\left(\frac{f(m) - f(m + \sigma z) }{ f(m) } \right)^2\right)  \ind{f(m + \sigma z) \leq f(m)} \right]
%  \\
%  &=
%  \e\left[\left(-\left(\frac{f(m) - f(m + \sigma z) }{ f(m) } \right) - \frac{1}{2}\left(\frac{f(m) - f(m + \sigma z) }{ f(m) } \right)^2\right)  \ind{\inner{\nabla f(m)}{z} \leq 0} \ind{f(m + \sigma z) \leq f(m)} \right]
%\end{align*}
%then, use lipschtz smoothness and we will obtian a better bound.          }
%\mori{
%let $h_2(z) = (f(m) -f(m+\sigma z))^2\cdot\ind{\inner{\nabla f(m)}{z} \leq 0}$ and $g(z) = \ind{f(m + \sigma z) \leq f(m)}$.
%these two are not negatively correlated, and then we cannot exploit the inequality we do in the case of $h(z)$ : $\E[h(z)g(z)] \leq \E[h(z)]\E[g(z)]$.
%}
%\yohe{
%the first term is what we have analyzed above. for the second term, use $\E[x^2 \cdot \ind{...}] = \E[(x \cdot \ind{...})^2] \geq \E[x \cdot \ind{...}]^2$, the right-most side of which is the square of the first term. therefore, we can obtain $(24) - (24)^2/2$ (here, i realized it does not help to enlarge the possible $p_{target}$.)
%}
%
%\yohe{some other 
%}
%\end{proof}
%\hrule


\subsection{Proof of \Cref{lemma:variancebound}}\label{apdx:lemma:variancebound}

Under \Cref{asm:L_U_2diff}, for any $x \in \R^d$, we have
$\frac{2 f(x)}{ U } \leq \norm{x}^2 \leq \frac{2 f(x)}{ L }$.
Then, we have
\begin{equation}
\begin{aligned}
\MoveEqLeft[3]
\log\left(\frac{f(m+\sigma z)}{f(m)}\right)\cdot\ind{f(m+\sigma z)\leq f(m)}
\\
=&
\min\left\{\log\left(\frac{f(m+\sigma z)}{f(m)}\right), 0 \right\}
\\
\geq&
\min\left\{
\log\left(\frac{\norm{m+\sigma z}^2}{\norm{m}^2}\right) + \log\left(\frac{L}{U}\right)
, 0 \right\}
\\
=&
\min\Bigg\{
\log\left(
1 + \frac{\norm{z}^2}{\norm{m}^2}\left(
\left(\sigma + \frac{\langle m, z \rangle}{\norm{z}^2}\right)^2
-\left(\frac{\langle m, z \rangle}{\norm{z}^2}\right)^2
\right)
\right) 
\\
&+ \log\left(\frac{L}{U}\right)
, 0 \Bigg\}
\\
\geq&
\log\left(
1 
-\frac{\langle m, z \rangle^2}{\norm{m}^2\cdot\norm{z}^2}
\right)
+ \log\left(\frac{L}{U}\right)
.
\end{aligned}
\end{equation}
Therefore,
\begin{equation}
\begin{aligned}[t]
&\exp\left(
\abs*{
\log\left(\frac{f(m+\sigma z)}{f(m)}\right)\cdot\ind{f(m+\sigma z)\leq f(m)}
}
\right)
\\
&\leq
\exp\left(
-
\log\left(
\frac{L}{U}\cdot\left(1 
-\frac{\langle m, z \rangle^2}{\norm{m}^2\cdot\norm{z}^2}
\right)
\right)
\right)
% =
% \frac{U}{L}\cdot
% \frac{\norm{z}^2}{\norm{z}^2-\left\langle \frac{m}{\norm{m}}, z \right\rangle^2}
\\
&=
\frac{U}{L}\cdot \left( 1 + \frac{1}{d-1}\cdot \frac{(d-1) \inner*{\frac{m}{\norm{m}}}{z}^2}{\norm{z}^2-\inner*{\frac{m}{\norm{m}}}{z}^2}\right)
.
\end{aligned}
\end{equation}
% Let $\{\mathcal{N}_i\}_{i=1}^d$ be a set of i.i.d random variables which follow the standard normal distribution.
% The RMS is stochastically dominated by
% % \begin{equation}
% % \frac{U}{L}\cdot
% % \left(
% % 1+ \frac{\mathcal{N}_1^2}{\sum_{i=2}^d\mathcal{N}_i^2}
% % \right)
% % .
% % \end{equation}
% $\frac{U}{L}\cdot\left(1+ \frac{\mathcal{N}_1^2}{\sum_{i=2}^d\mathcal{N}_i^2}\right)$.
% A fact is $\frac{(d-1)\cdot\mathcal{N}_1^2}{\sum_{i=2}^d\mathcal{N}_i^2}$ follows F-distribution and its expectation is $(d-1)/(d-3)$.
Here, $\frac{(d-1) \inner*{\frac{m}{\norm{m}}}{z}^2}{\norm{z}^2-\inner*{\frac{m}{\norm{m}}}{z}^2}$ is $F$-distributed and its expectation is $(d-1) / (d-3)$. 
Finally, we obtain \eqref{eq:lemma:variancebound}.
%\begin{multline}
%\E\left[\exp\left(
%\abs*{
%\log\left(\frac{f(m+\sigma z)}{f(m)}\right)\cdot\ind{f(m+\sigma z)\leq f(m)}
%}
%\right)
%\right]
%\\
%\leq
%\frac{U}{L}\cdot
%\left(
%1+\frac{1}{d-3}
%\right)
%.
%\end{multline}
This completes the proof.

%          We have
%          \begin{align}
%            &\log\left( \frac{f(m + \sigma \cdot z) }{ f(m)} \right) \cdot \ind{f(m + \sigma \cdot z) \leq f(m)}
%            \\
%            =& \log\left( 1 +  \frac{f(m+\sigma z)-f(m)}{f(m)}\right)
%            \cdot \ind{f(m + \sigma \cdot z) \leq f(m)}
%            \\
%            \geq&
%            \log\left(1+\frac{1}{f(m)}\cdot\left(\sigma\norm{\nabla f(m)} z_e+ \frac{\sigma^2 L \norm{z}^2}{2}\right)\right)
%            \cdot \ind{f(m + \sigma \cdot z) \leq f(m)}
%            \\
%            =&
%            \log\left(1+\frac{L\norm{z}^2}{2 f(m)}\cdot
%            \left(
%            \left(\sigma + \frac{\norm{\nabla f(m)} z_e}{L\norm{z}^2}\right)^2
%             - \left(\frac{\norm{\nabla f(m)} z_e}{L\norm{z}^2}\right)^2
%            \right)\right)
%            \cdot \ind{f(m + \sigma \cdot z) \leq f(m)}
%            \label{eq:progress_rate_with_LU}
%            .
%            \end{align}
%            The RMS is minimized on $\sigma>0$ when
%            \begin{equation}
%                \sigma = \sigma^* :=
%                \begin{cases}
%                    { 0  \text{ if } z_e >0,}
%                    \\
%                    {-\frac{\norm{\nabla f(m)} z_e}{L\norm{z}^2}
%                    \text{ if } z_e\leq 0 .}
%                \end{cases}
%            \end{equation}          
%            Moreover,
%            \begin{equation}
%            \left(\sigma_L^* + \frac{\norm{\nabla f(m)} z_e}{L\norm{z}^2}\right)
%             - \left(\frac{\norm{\nabla f(m)} z_e}{L\norm{z}^2}\right)^2
%             = -\min\left\{\frac{\norm{\nabla f(m)} z_e}{L\norm{z}^2}, 0\right\}^2
%             .
%            \end{equation}          
%            Therefore, the RMS of \eqref{eq:progress_rate_with_LU} is bounded from below by
%            \begin{align}
%                &
%                \log\left(1-\min\left\{\frac{\norm{\nabla f(m)}z_e}{2L\norm{z}^2 f(m)}, 0\right\}^2\right)
%                ,
%                \\
%                \geq&
%                \log\left(1-\frac{(\norm{\nabla f(m)}z_e)^2}{2L\norm{z}^2 f(m)}\right)
%            \end{align}
%            Further, as $\norm{\nabla f(m)}\leq \sqrt{U f(m)}$ with $U$-Lipschitz smoothness, the RHS of the above, and thus $\log(f(m+\sigma z)/f(m))\cdot\ind{f(m+\sigma z)\leq f(m)}$, is bounded from below as
%            \begin{equation}
%                \log\left(\frac{f(m+\sigma z)}{f(m)}\right)\cdot\ind{f(m+\sigma z)\leq f(m)}
%                \geq 
%                \log\left(1-\frac{U z_e^2}{2L\norm{z}^2}\right)
%                .
%            \end{equation}
%            \todo{WIP}
%            Finally, taking $\exp$ and the expectation on both sides of the above,
%            \begin{equation}
%                \E\left[\exp\left(\log\left(\frac{f(m+\sigma z)}{f(m)}\right)\cdot\ind{f(m+\sigma z)\leq f(m)}\right)\right]
%                \geq 
%                1-\frac{U}{2L}\cdot\E\left[\frac{z_e^2}{\norm{z}^2}\right]
%                .
%            \end{equation}
%            In light of \cite[Theorem 6]{magnus1986exact},
%            \begin{equation}
%                \E\left[\frac{z_e^2}{\norm{z}^2}\right]
%                = \int_0^\infty (1+2t)^{-\frac{d}{2}-1} dt
%                = \frac{1}{d}
%                .
%            \end{equation}
%            Therefore, we obtain
%            \begin{equation}
%                \E\left[\exp\left(\log\left(\frac{f(m+\sigma z)}{f(m)}\right)\cdot\ind{f(m+\sigma z)\leq f(m)}\right)\right]
%                \geq 
%                1-\frac{U}{2dL}
%                .
%            \end{equation}


\subsection{Proof of \Cref{lemma:successprobability}}\label{apdx:lemma:successprobability}

Let 
$A = \bar\sigma/2$
%\frac{\sigma \E[\mathcal{Q}_z(\theta)]}{2 \| \nabla f(m) \| }$
and $B = \mathcal{Q}_z(\theta)/\E[\mathcal{Q}_z(\theta)]$.
First, we derive a few inequalities before deriving the upper bound and the lower bound.
By using Chebyshev's inequality, we obtain
\begin{align}
\Pr\left[\left| B - 1 \right|\geq \epsilon\right]
%=
%\Pr\left[\left| \frac{\mathcal{Q}_z(\theta)}{\E[\mathcal{Q}_z(\theta)]} - 1 \right|\geq \epsilon\right]
\leq
\frac{1}{\epsilon^2} \cdot \mathcal{V}_\mathrm{std}
.
\end{align}
%
Because of \eqref{eq:q-def}, we can write the success probability as
\begin{equation}
\begin{aligned}[t]
\Pr[ f(m + \sigma z) \leq f(m)]
%&=
%\Pr\left[ \sigma \| \nabla f(m) \| z_e \leq - \frac{\sigma^2 }{2} \mathcal{Q}_z(\theta) \right] 
%\\
%&= \Pr\left[ z_e \leq -\frac{\bar\sigma}{2} \frac{\mathcal{Q}_z(\theta) }{\E[\mathcal{Q}_z(\theta)]} \right]
= \Pr\left[ z_e \leq - A \cdot B\right]
.
\end{aligned}
\end{equation}
%   %\yohe{The last equality is not correct. It is (25) $\leq$ (26). The reason is because $\ind{z_e \leq 0} = 0$ leads to equality holding inside the probability in (26).}
%   Since $- \frac{\sigma^2 }{2}\mathcal{Q}_z(\theta) \leq - \frac{\sigma^2 }{2}\mathcal{Q}_z(\theta) \cdot\ind{z_e \leq 0} $, we have the following upper bound:
%   \begin{align}
%       \Pr\left[ \sigma \| \nabla f(m) \| z_e \leq - \frac{\sigma^2 }{2}\mathcal{Q}_z(\theta) \right]
%     \leq&  \Pr\left[ \sigma \| \nabla f(m) \| z_e \leq - \frac{\sigma^2 }{2}\mathcal{Q}_z(\theta) \cdot\ind{z_e \leq 0} \right]
%     \\
%     =&  \Pr\left[ z_e \leq - \frac{\sigma \E[\mathcal{Q}_z(\theta)]}{2 \| \nabla f(m) \| } \frac{\mathcal{Q}_z(\theta) }{\E[\mathcal{Q}_z(\theta)]} \right].
%   \end{align}
% %For the other side, use
% %Since $z$ is a continuous random variable which follows the $d$-dimensional normal distribution, the RHS of the above is 
% In addition, we also have a similar lower bound:
%   \begin{align}
%       \Pr\left[ \sigma \| \nabla f(m) \| z_e \leq - \frac{\sigma^2 }{2}\mathcal{Q}_z(\theta) \right]
%     \geq& \Pr\left[ \sigma \| \nabla f(m) \| z_e < - \frac{\sigma^2 }{2}\mathcal{Q}_z(\theta) \right]
%     \\
%     =& \Pr\left[ \sigma \| \nabla f(m) \| z_e < - \frac{\sigma^2 }{2}\mathcal{Q}_z(\theta) \cdot \ind{z_e \leq 0}\right]
%     \\
%     =&  \Pr\left[ z_e < - \frac{\sigma \E[\mathcal{Q}_z(\theta)]}{2 \| \nabla f(m) \| } \frac{\mathcal{Q}_z(\theta) }{\E[\mathcal{Q}_z(\theta)]} \right]
%     %\\
%     %=&  \Pr\left[ z_e \leq - \frac{\sigma \E[\mathcal{Q}_z(\theta)]}{2 \| \nabla f(m) \| } \frac{\mathcal{Q}_z(\theta) }{\E[\mathcal{Q}_z(\theta)]} \right]
%     .
%   \end{align}
%   %\mori{To think twice : Is the above exchange of $\leq$ for $<$ really trivial, or even correct?}
%   %\yohe{Right, it is not crystal clear. But the result must be correct. 
%   %All we need to prove is 
%   %\begin{align*}
%   %\Pr\left[ z_e = - \frac{\sigma \E[\mathcal{Q}_z(\theta)]}{2 \| \nabla f(m) \| } \frac{\mathcal{Q}_z(\theta) }{\E[\mathcal{Q}_z(\theta)]} \right] = 0 .
%   %\end{align*}
%   %It is enough to show that
%   %\begin{align*}
%   %\Pr\left[ z_e = - \frac{\sigma \E[\mathcal{Q}_z(\theta)]}{2 \| \nabla f(m) \| } \frac{\mathcal{Q}_z(\theta) }{\E[\mathcal{Q}_z(\theta)]} \mid z_e \right] = 0 .
%   %\end{align*}
%   %(If the above conditional probability is zero, then its expectation over $z_e$ is zero as well.)
%   %If $z_e > 0$, the equality never holds. 
%   %Given $z_e \leq 0$, $\phi$ is a strongly convex function of the other components of $z$ orthogonal to $z_e$. 
%   %Since $z$ itself is absolutely continuous w.r.t.\ the Lebesgue measure, so is $\phi$. 
%   %It implies that the above probability is always zero. 
%   %}
%   %In the end, whether $\leq$ or $=$ should not matter as $z_e$ is normally distributed. These inequalities should suffice.

Now, we derive the upper bound of the success probability as
\begin{equation}
\begin{aligned}[t]
%\MoveEqLeft[3]
\Pr\left[f(m + \sigma z) \leq f(m) \right]
%=&
%\Pr\left[ z_e \leq - A \cdot B \right]
%\\
=&
\Pr\left[ \left(z_e \leq - A \cdot B\right)\wedge \left( B > 1-\epsilon\right)\right]
\\
&+\Pr\left[ \left(z_e \leq - A \cdot B\right)\wedge\left(B \leq 1-\epsilon\right)\right]
%\\
%\leq&
%\Pr\left[ z_e \leq - A \cdot(1-\epsilon)\right]
%+\Pr\left[ B \leq 1-\epsilon\right]
\\
\leq&
\Pr\left[ z_e \leq - A \cdot(1-\epsilon)\right]
+\Pr\left[ \left| B -1 \right| \geq \epsilon\right]
\\
\leq&
\Phi\left( - \frac12 \bar\sigma\cdot(1-\epsilon)\right)
+ \frac{1}{\epsilon^2} \cdot \mathcal{V}_\mathrm{std}
.
\end{aligned}
\end{equation}
% \begin{equation}
% \begin{aligned}[t]
% \MoveEqLeft[3]\Pr\left[f(m + \sigma z) \leq f(m) \right]
% \\
% =&
% \Pr\left[ z_e \leq - \frac12 \bar\sigma \frac{\mathcal{Q}_z(\theta)}{\E[\mathcal{Q}_z(\theta)]} \right]
% \\
% =&
% \Pr\left[ \left(z_e \leq - \frac12 \bar\sigma\cdot\frac{\mathcal{Q}_z(\theta)}{\E[\mathcal{Q}_z(\theta)]}\right)\wedge \left(\frac{\mathcal{Q}_z(\theta)}{\E[\mathcal{Q}_z(\theta)]}> 1-\epsilon\right)\right]
% \\
% &+\Pr\left[ \left(z_e \leq - \frac12 \bar\sigma\cdot\frac{\mathcal{Q}_z(\theta)}{\E[\mathcal{Q}_z(\theta)]}\right)\wedge\left(\frac{\mathcal{Q}_z(\theta)}{\E[\mathcal{Q}_z(\theta)]}\leq 1-\epsilon\right)\right]
% \\
% \leq&
% \Pr\left[ z_e \leq - \frac12 \bar\sigma\cdot(1-\epsilon)\right]
% +\Pr\left[ \frac{\mathcal{Q}_z(\theta)}{\E[\mathcal{Q}_z(\theta)]}\leq 1-\epsilon\right]
% \\
% \leq&
% \Pr\left[ z_e \leq - \frac12 \bar\sigma\cdot(1-\epsilon)\right]
% +\Pr\left[ \left| \frac{\mathcal{Q}_z(\theta)}{\E[\mathcal{Q}_z(\theta)]} -1 \right| \geq \epsilon\right]
% \\
% \leq&
% \Phi\left( - \frac12 \bar\sigma\cdot(1-\epsilon)\right)
% + \frac{1}{\epsilon^2} \cdot \mathcal{V}_\mathrm{std}
% .
% \end{aligned}
% \end{equation}

Analogously, we derive the lower bound:
\begin{equation}
\begin{aligned}[t]
%\MoveEqLeft[2]
\Pr\left[f(m + \sigma z) > f(m) \right]
%\\
%=&
%\Pr\left[ z_e > - A \cdot B \right]
%\\
%=&
%\Pr\left[ \left(z_e > - A \cdot B\right)\wedge\left(B < 1+\epsilon\right)\right]
%\\
%&+\Pr\left[ \left(z_e > - A \cdot B\right)\wedge\left(B \geq 1+\epsilon\right)\right]
%%\\
%%\leq&
%%\Pr\left[ z_e > - A \cdot(1+\epsilon)\right]
%%+\Pr\left[ B \geq 1+\epsilon\right]
%\\
\leq&
\Pr\left[ z_e > - A \cdot(1+\epsilon)\right]
+\Pr\left[ \left| B -1 \right| \geq \epsilon\right]
\\
\leq&
%\Pr\left[ z_e > - \frac12 \bar\sigma\cdot(1+\epsilon)\right]
%+ \frac{1}{\epsilon^2} \cdot \mathcal{V}_\mathrm{std}
%\\
%=&
1 - \Phi\left( - \frac12 \bar\sigma\cdot(1+\epsilon)\right)
+ \frac{1}{\epsilon^2} \cdot \mathcal{V}_\mathrm{std}
.
\end{aligned}
\end{equation}
% \begin{equation}
% \begin{aligned}[t]
% \MoveEqLeft[3]\Pr\left[f(m + \sigma z) > f(m) \right]
% \\
% =&
% \Pr\left[ z_e > - \frac12 \bar\sigma \frac{\mathcal{Q}_z(\theta)}{\E[\mathcal{Q}_z(\theta)]} \right]
% \\
% =&
% \Pr\left[ \left(z_e > - \frac12 \bar\sigma\cdot\frac{\mathcal{Q}_z(\theta)}{\E[\mathcal{Q}_z(\theta)]}\right)\wedge\left(\frac{\mathcal{Q}_z(\theta)}{\E[\mathcal{Q}_z(\theta)]}< 1+\epsilon\right)\right]
% \\
% &+\Pr\left[ \left(z_e > - \frac12 \bar\sigma\cdot\frac{\mathcal{Q}_z(\theta)}{\E[\mathcal{Q}_z(\theta)]}\right)\wedge\left(\frac{\mathcal{Q}_z(\theta)}{\E[\mathcal{Q}_z(\theta)]}\geq 1+\epsilon\right)\right]
% \\
% \leq&
% \Pr\left[ z_e > - \frac12 \bar\sigma\cdot(1+\epsilon)\right]
% +\Pr\left[ \frac{\mathcal{Q}_z(\theta)}{\E[\mathcal{Q}_z(\theta)]}\geq 1+\epsilon\right]
% \\
% \leq&
% \Pr\left[ z_e > - \frac12 \bar\sigma\cdot(1+\epsilon)\right]
% +\Pr\left[ \left| \frac{\mathcal{Q}_z(\theta)}{\E[\mathcal{Q}_z(\theta)]} -1 \right| \geq \epsilon\right]
% \\
% \leq&
% \Pr\left[ z_e > - \frac12 \bar\sigma\cdot(1+\epsilon)\right]
% + \frac{1}{\epsilon^2} \cdot \mathcal{V}_\mathrm{std}
% \\
% =&
% 1 - \Phi\left( - \frac12 \bar\sigma\cdot(1+\epsilon)\right)
% + \frac{1}{\epsilon^2} \cdot \mathcal{V}_\mathrm{std}
% .
% \end{aligned}
% \end{equation}
%Hence, we have
%\begin{equation}
%\Pr\left[f(m + \sigma z) \leq f(m) \right] 
%\geq \Phi\left( - \frac12 \bar\sigma\cdot(1+\epsilon)\right)
%- \frac{1}{\epsilon^2} \cdot \mathcal{V}_\mathrm{std}
%.
%\end{equation}
This completes the proof.

\subsection{Proof of \Cref{cor:successprobability}}\label{apdx:cor:successprobability}

%%\begin{proof}%[Proof of \Cref{cor:successprobability}]

First, we prove \eqref{eq:cor:successprobability:large}.
Define 
\begin{equation}
B_\theta^\mathrm{high}(q; \epsilon) := 2\cdot\Phi^{-1} \left(1 - \left(q + \frac{1}{\epsilon^2} \cdot \mathcal{V}_\mathrm{std}\right) \right)/ (1 + \epsilon)
\label{eq:b_high_q_e}
\end{equation}
for $q\in(0, 1/2)$ and $\epsilon > \sqrt{\frac{2}{1-2q}\cdot\mathcal{V}_\mathrm{std}}$.
Apparently, $B_\theta^\mathrm{high}(q; \epsilon)$ is continuous with respect to both arguments and $B_\theta^\mathrm{high}(q; \epsilon) > 0$. 
Since $B_\theta^\mathrm{high}(q; \epsilon) \to 0$ as $\epsilon \downarrow \sqrt{\frac{2}{1-2q}\cdot\mathcal{V}_\mathrm{std}}$ and $B_\theta^\mathrm{high}(q; \epsilon) \to 0$ as $\epsilon \uparrow + \infty$, there must exist at least one $\epsilon$ for each $q$ such that $B_\theta^\mathrm{high}(q; \epsilon)$ is maximized. Let $\epsilon^\text{high}(q)$ be the smallest one of such maximizers. Then, $B_\theta^\text{high}(q) = B_\theta^\text{high}(q; \epsilon^\text{high}(q))$.

In light of \Cref{lemma:successprobability}, we have
\begin{equation}
\begin{aligned}[t]
%\MoveEqLeft[2]
\bar\sigma \leq  B_\theta^\mathrm{high}(q; \epsilon)
&\Leftrightarrow 
q\leq \Phi\left(-\frac12 \bar\sigma\cdot (1+\epsilon)\right) - \frac{1}{\epsilon^2}\cdot\mathcal{V}_\mathrm{std}
\\
&\Rightarrow
\Pr\left[f(m+\sigma z)\leq f(m)\right] > q
%\label{eq:7}
\end{aligned}
\end{equation}
for any $q\in(0, 1/2)$ and $\epsilon > \sqrt{\frac{2}{1-2q}\cdot\mathcal{V}_\mathrm{std}}$.
By letting $\epsilon = \epsilon^\text{high}(q)$, the above relation implies \eqref{eq:cor:successprobability:large}. This completes the proof of \eqref{eq:cor:successprobability:large}.
It is trivial to see that $B_\theta^\mathrm{high}(q) \leq 2 \Phi^{-1}(1-q)$ for all $q \in \left(0, \frac12\right)$.
% Note that $\bar\sigma>0$ under the convexity of $f$ (\Cref{asm:L_U_2diff}).
% Moreover, 
% $B_\theta^\mathrm{high}(q;\epsilon)$ is positive for any $\theta\in\Theta\setminus\{\theta|m\neq x^*\}$ if $q\in(0, 1/2)$ and $\epsilon > \sqrt{\frac{2}{1-2q}\cdot\mathcal{V}_\mathrm{std}}$.
% By taking $\sup$ in $\epsilon > \sqrt{\frac{2}{1-2q}\cdot\mathcal{V}_\mathrm{std}}$,
% the proof of \eqref{eq:cor:successprobability:large} is completed.
% \yohe{I do not think that it is problematic, but the LHS of \eqref{eq:cor:successprobability:large} should be $... < B_\theta^\text{high}(q)$ (without equality) as we are taking $\sup$ and the existance of $\max$ is not guaranteed. But I guess it is easy to show that $\sup = \max$ by using the continuity of B in $\epsilon$. }

%      First, we prove \eqref{eq:cor:successprobability:large}.
%      In light of \Cref{lemma:successprobability}, it is enough to show that for each $q \in \left(0, \frac12\right)$, there exists an $\epsilon > 0$ such that
%        \begin{equation}
%        \bar\sigma\leq B_\theta^\mathrm{high}(q)
%        \Rightarrow 
%            q\leq \Phi\left(-\frac12 \bar\sigma\cdot (1+\epsilon)\right) - \frac{1}{\epsilon^2}\cdot\mathcal{V}_\mathrm{std}
%            .
%            \label{eq:4}
%        \end{equation}
%    %   is sufficient to prove
%    % \begin{equation}
%    %     \bar\sigma\leq B_\theta^\mathrm{high}(q)
%    %     \Rightarrow 
%    %     \mathrm{Pr}[f(m+\sigma z)\leq f(m)]>q
%    %     \label{eq:5}
%    %     .
%    % \end{equation}
%      Now suppose that $\theta$ satisfies $\bar\sigma\leq B_\theta^\mathrm{high}(q)$.
%    Given any $q\in(0, 1/2)$, the RHS of \eqref{eq:cor:successprobability:large} holds if the following condition is satisfied :
%    \begin{multline}
%        q\leq \Phi\left(-\frac12 \bar\sigma\cdot (1+\epsilon)\right) - \frac{1}{\epsilon^2}\cdot\mathcal{V}_\mathrm{std}
%        \\
%        \Leftrightarrow
%        \bar\sigma \leq \frac{ 2 \Phi^{-1} \left(1 - \left(q + \frac{1}{\epsilon^2} \cdot \mathcal{V}_\mathrm{std}\right) \right)}{ 1 + \epsilon}
%        =: B_\theta^\mathrm{high}(q; \epsilon)
%        \label{eq:7}
%        .
%    \end{multline}
%    Note that $\bar\sigma$ can never be negative with the definitions of $\sigma, \mathcal{Q}_z(\theta)$ under \Cref{asm:L_U_2diff}.
%    For any $q\in(0, 1/2)$, there exists a positive constant $\epsilon$ which satisfies $\epsilon>\sqrt{\frac{2}{1-2q}\cdot\mathcal{V}_\mathrm{std}}$,
%    and if $\epsilon>\sqrt{\frac{2}{1-2q}\cdot\mathcal{V}_\mathrm{std}}$, it holds
%    $B_\theta^\mathrm{high}(q; \epsilon)>0$.
%    Therefore, now we see that there exists $\epsilon^\text{high}(q) > \sqrt{ \frac{ 2 }{1 - 2q} \cdot \mathcal{V}_\mathrm{std} }$ such that $B_\theta^\text{high}(q) = B_\theta^\text{high}(q; \epsilon^\text{high}(q))$.
%    Finally, as \eqref{eq:4} holds for any $\epsilon>0$ given any $q\in(0, 1/2)$, the LHS of \eqref{eq:cor:successprobability:large} also holds if
%    \begin{equation}
%        \bar\sigma \leq \sup_{\epsilon>\sqrt{\frac{2}{1-2q}\cdot\mathcal{V}_\mathrm{std}}}\frac{ 2 \Phi^{-1} \left(1 - \left(q + \frac{1}{\epsilon^2} \cdot \mathcal{V}_\mathrm{std}\right) \right)}{ 1 + \epsilon}
%        =: B_\theta^\mathrm{high}(q)
%        .
%    \end{equation}
%    This completes the proof of \eqref{eq:cor:successprobability:large}.

%Next, we prove several properties of $B_\theta^\text{high}$.
%First, note that $B_\theta^\text{high}(q; \epsilon)$ is continuous and strictly decreasing with respect to $q$ for each $\epsilon$.
%We prove the right-continuity and strict-decrease of $B_\theta^\text{high}$ by contradiction.
%If $B_\theta^\text{high}$ does not strictly decrease at a point $q$ in the domain, then there must exist $q' > q$ such that $B_\theta^\text{high}(q) \leq B_\theta^\text{high}(q')$.
%However, $B_\theta^\text{high}(q') = B_\theta^\text{high}(q'; \epsilon^\text{high}(q')) < B_\theta^\text{high}(q; \epsilon^\text{high}(q')) \leq B_\theta^\text{high}(q; \epsilon^\text{high}(q)’) = B_\theta^\text{high}(q’)$, which contradicts $B_\theta^\text{high}(q) \leq B_\theta^\text{high}(q')$.
%Hence, $B_\theta^\text{high}$ is strictly decreasing.
%If $B_\theta^\text{high}$ is not right-continuous at a point $q$, there must exist $\delta > 0$ such that $B_\theta^\text{high}(q) - \delta \geq B_\theta^\text{high}(q')$ for all $q' > q$.
%Because $B_\theta^\text{high}(q; \epsilon^\text{high}(q)) = B_\theta^\text{high}(q’) > B_\theta^\text{high}(q') = B_\theta^\text{high}(q'; \epsilon^\text{high}(q')) \geq B_\theta^\text{high}(q'; \epsilon^\text{high}(q))$, we have $\abs{B_\theta^\text{high}(q) - B_\theta^\text{high}(q')} \leq \abs{B_\theta^\text{high}(q; \epsilon^\text{high}(q)) - B_\theta^\text{high}(q'; \epsilon^\text{high}(q))}$. The LHS must be no smaller than $\delta$ for any $q' > q$, whereas from the continuity of $B_\theta^\text{high}(\cdot; \epsilon^\text{high}(q))$, the RHS satisfies $\lim_{q' \downarrow q} \abs{B_\theta^\text{high}(q; \epsilon^\text{high}(q)) - B_\theta^\text{high}(q'; \epsilon^\text{high}(q))} = 0$, which is a contradiction.
%Hence, $B_\theta^\text{high}$ is right-continuous.

Similarly, we prove \eqref{eq:cor:successprobability:small}.
Define
\begin{equation}
B_\theta^\mathrm{low}(q; \epsilon)
:= 
2 \cdot\Phi^{-1} \left(1 - \left(q - \frac{1}{\epsilon^2} \cdot \mathcal{V}_\mathrm{std}\right) \right)/ (1 - \epsilon)\label{eq:b_low_q_e}
\end{equation}
for $q\in\left( \mathcal{V}_\mathrm{std} , \frac{1}{2}\right)$ and $\epsilon\in \left(\sqrt{\frac{1}{q}\cdot\mathcal{V}_\mathrm{std}}, 1\right)$.
It is easy to see that $B_\theta^\mathrm{low}(q; \epsilon)$ is continuous with respect to both arguments and $B_\theta^\mathrm{low}(q; \epsilon) > 0$.
Moreover, since $B_\theta^\mathrm{low}(q; \epsilon) \to +\infty$ as $\epsilon \uparrow 1$ and $B_\theta^\mathrm{low}(q; \epsilon) \to +\infty$ as $\epsilon \downarrow \sqrt{\frac{1}{q}\cdot\mathcal{V}_\mathrm{std}}$, there must exist at least one $\epsilon$ for each $q$ such that $B_\theta^\mathrm{low}(q; \epsilon)$ is minimized. Let $\epsilon^\text{low}(q)$ be the smallest one of such minimizers. Then, $B_\theta^\text{low}(q) = B_\theta^\text{low}(q; \epsilon^\text{low}(q))$.

In light of \Cref{lemma:successprobability}, we have
\begin{equation}
\begin{aligned}[t]
    \label{eq:8}
\MoveEqLeft[2]
\bar\sigma \leq 
B_\theta^\mathrm{low}(q; \epsilon)
\Leftrightarrow
q\geq \Phi\left(-\frac12 \bar\sigma\cdot (1-\epsilon)\right) + \frac{1}{\epsilon^2}\cdot\mathcal{V}_\mathrm{std}
\\
&\Rightarrow
\Pr\left[f(m+\sigma z)\leq f(m)\right] < q
,
\end{aligned}
\end{equation}
for any $q\in\left( \mathcal{V}_\mathrm{std} , \frac{1}{2}\right)$ and $\epsilon\in \left(\sqrt{\frac{1}{q}\cdot\mathcal{V}_\mathrm{std}}, 1\right)$.
By letting $\epsilon = \epsilon^\text{low}(q)$, the above relation implies \eqref{eq:cor:successprobability:small}. This completes the proof of \eqref{eq:cor:successprobability:small}.
It is trivial to see that $B_\theta^\mathrm{high}(q) \geq 2 \Phi^{-1}(1-q)$ for all $q\in\left( \mathcal{V}_\mathrm{std} , \frac{1}{2}\right)$.

% Note that $B_\theta^\mathrm{low}(q;\epsilon)$ is positive 
% for any $\theta\in\Theta\setminus\{\theta|m\neq x^*\}$
% if $q\in\left( \mathcal{V}_\mathrm{std} , \frac{1}{2}\right)$ and $\epsilon\in \left(\sqrt{\frac{1}{q}\cdot\mathcal{V}_\mathrm{std}}, 1\right)$.
% The condition $\epsilon < 1$ is also necessary to have the RHS of \eqref{eq:lemma:successprobability} smaller than $1/2$.
% By taking $\inf$ in $\epsilon\in \left(\sqrt{\frac{1}{q}\cdot\mathcal{V}_\mathrm{std}}, 1\right)$,
% the proof of \eqref{eq:cor:successprobability:small} is completed.

%      Similarly, we prove \eqref{eq:cor:successprobability:small}.
%      In light of \Cref{lemma:successprobability}, for each $q \in \left( \mathcal{V}_\mathrm{std} , \frac{1}{2}\right)$, it is sufficient to show that there exists an $\epsilon \in (0, 1)$ such that the RHS of \eqref{eq:lemma:successprobability} is no greater than $q$.
%      The condition $\epsilon < 1$ is necessary to have the RHS of \eqref{eq:lemma:successprobability} smaller than $1/2$.
%      By solving this inequality for $\bar\sigma$, we obtain
%      \begin{equation}
%        \bar\sigma \geq \frac{ 2 \Phi^{-1} \left(1 - \left(q - \frac{1}{\epsilon^2} \cdot \mathcal{V}_\mathrm{std}\right) \right)}{ 1 - \epsilon}  =: B_\theta^\text{low}(q; \epsilon).
%        \label{eq:cor:sp:small:eps}
%      \end{equation}
%      That is, if there exists an $\epsilon \in (0, 1)$ such that the condition above holds, then we have $\Pr\left[f(m + \sigma z) \leq f(m) \right] < q$.
%      In contrast, for each $q \in \left( \mathcal{V}_\mathrm{std} , \frac{1}{2}\right)$, we can easily see that there exists $\epsilon_H^\text{low}(q) \in \left(\sqrt{ \frac{ 1 }{q} \cdot \mathcal{V}_\mathrm{std} }, 1\right)$ such that $B_\theta^\text{low}(q) = B_\theta^\text{low}(q; \epsilon^\text{low}(q))$.
%      Hence, we obtain \eqref{eq:cor:successprobability:large}.

%Finally, we prove several properties of $B_\theta^\text{low}$.
%Note that $B_\theta^\text{high}(q; \epsilon)$ is continuous and strictly decreasing with respect to $q$ for each $\epsilon$.
%We prove the left-continuity and strict decrease of $B_\theta^\text{low}$ by contradiction.
%If $B_\theta^\text{low}$ is not strictly decreasing at a point $q$ in the domain, then there must exist $q' < q$ such that $B_\theta^\text{low}(q) \geq B_\theta^\text{low}(q')$.
%However, $B_\theta^\text{low}(q') = B_\theta^\text{low}(q'; \epsilon^\text{low}(q')) > B_\theta^\text{low}(q; \epsilon^\text{low}(q')) \geq B_\theta^\text{low}(q; \epsilon^\text{low}(q)’) = B_\theta^\text{low}(q’)$, which contradicts $B_\theta^\text{low}(q) \geq B_\theta^\text{low}(q')$.
%Hence, $B_\theta^\text{high}$ is strictly decreasing.
%If $B_\theta^\text{low}$ is not left-continuous at a point $q$, there must exist $\delta > 0$ such that $B_\theta^\text{low}(q) + \delta \leq B_\theta^\text{low}(q')$ for all $q' < q$.
%Because $B_\theta^\text{low}(q; \epsilon^\text{low}(q)) = B_\theta^\text{low}(q’) < B_\theta^\text{low}(q') = B_\theta^\text{low}(q'; \epsilon^\text{low}(q')) \leq B_\theta^\text{low}(q'; \epsilon^\text{low}(q))$, we have $\abs{B_\theta^\text{low}(q) - B_\theta^\text{low}(q')} \leq \abs{B_\theta^\text{low}(q; \epsilon^\text{low}(q)) - B_\theta^\text{low}(q'; \epsilon^\text{low}(q))}$. The LHS must be no smaller than $\delta$ for any $q' < q$, whereas from the continuity of $B_\theta^\text{low}(\cdot; \epsilon^\text{low}(q))$, the RHS satisfies $\lim_{q' \uparrow q} \abs{B_\theta^\text{low}(q; \epsilon^\text{low}(q)) - B_\theta^\text{low}(q'; \epsilon^\text{low}(q))} = 0$, which is a contradiction.
%Hence, $B_\theta^\text{low}$ is left-continuous.

%%\end{proof}




\subsection{Proof of \Cref{lemma:cases}}\label{apdx:lemma:cases}

%%\begin{proof}%[Proof of \Cref{lemma:cases}]
Proof of (i). 
The 1st term in $\max$ of the lower limit of $I_q$ is smaller than $1/8$. This is because the RHS of \eqref{eq:tr_cond} cannot be greater than or equal to $1/8$; hence, \eqref{eq:tr_cond} implies that 
$\underset{\theta\in\Thetazero}{\sup} \left\{\mathcal{V}_\mathrm{std}\right\} < 1/8$. 
%$\underset{\theta\in\Thetazero}{\sup} \left\{\Var[\mathcal{Q}_z(\theta)] / (\E[\mathcal{Q}_z(\theta)])^2\right\} < 1/8$. 
Therefore, it suffices to show that the 2nd term in $\max$ of the lower limit of $I_q$ is smaller than $1/2$.
Namely, we want to show that there exists $\bar{q} < 1/2$ such that $\kappa_{\inf} \cdot \sqrt{\frac{2}{\pi}} > B_{\sup}^\mathrm{low}(q)$ for all $q \in (\bar{q}, 1/2)$.
% we want to show $         
%   \inf\left\{
%         q:
%         \kappa_{\inf}
%         \cdot
%         \sqrt{\frac{2}{\pi}} > 
%         \underset{\theta\in\Theta\setminus\{\theta|m\neq x^*\}}{\sup}
%         B_\theta^\mathrm{low}(q) 
%         \right\}
%         < \frac12$.
Let $B_\theta^\text{low}(q;\epsilon)$ as defined in \eqref{eq:b_low_q_e}. 
Because 
$\mathcal{V}_\mathrm{std}< 1/8$,
%$\Var[\mathcal{Q}_z(\theta)] / (\E[\mathcal{Q}_z(\theta)])^2< 1/8$,
$B_\theta^\text{low}(q; \epsilon)$ is defined at least for $q \in [1/8, 1/2)$ and $\epsilon \in [1/\sqrt{8 q}, 1)$.
% and we have $B_\theta^\text{low}(q; 1/2) \geq B_\theta^\text{low}(q)$ for $q \in [1/8, 1/2)$ for any $\theta\in\Theta\setminus\{\theta\mid m\neq x^*\}$ as $1/2 \in [1/\sqrt{8 q}, 1)$ for any $q \in [1/8, 1/2)$. 
If we let 
\begin{equation}
\delta_1 = \Phi\left(\frac{\kappa_{\inf}}{2}\cdot\frac{1}{\sqrt{2\pi}}\right) - \frac12 - 4\cdot\underset{\theta\in\Thetazero}{\sup}
\mathcal{V}_\mathrm{std} ,
\end{equation}
it is strictly positive because of \eqref{eq:tr_cond}.    
For any $q \in [1/8, 1/2)$, we have
\begin{equation}
\begin{split}
\label{eq:sup_b_low_upper}
&B_{\sup}^\text{low}(q)
\leq
\sup_{\theta \in \Theta_0} B_\theta^\mathrm{low}(q; \epsilon=1/\sqrt{8q})
\\
&=
\sup_{\theta \in \Theta_0}
    \frac{2}{1 - 1/\sqrt{8q}} \Phi^{-1} \left(1 - q + 8 q \cdot \mathcal{V}_\mathrm{std}\right)
\\
&=
\frac{2}{1 - 1/\sqrt{8q}} \Phi^{-1} \left( 1 + 2q \left(\Phi\left(\frac{\kappa_{\inf}}{2}\cdot\frac{1}{\sqrt{2\pi}}\right) - 1 - \delta_1 \right)\right) .
\end{split}
\end{equation}    
Because the RMS of \eqref{eq:sup_b_low_upper} is continuous with respect to $q$ and it approaches $4 \Phi^{-1} \left( \Phi\left(\frac{\kappa_{\inf}}{2}\cdot\frac{1}{\sqrt{2\pi}}\right) - \delta_1 \right) <
4 \Phi^{-1} \left( \Phi\left(\frac{\kappa_{\inf}}{2}\cdot\frac{1}{\sqrt{2\pi}}\right) \right) = \kappa_{\inf}\cdot \sqrt{2 / \pi}$ as $q \to 1/2$, there exists a $\bar{q} \in [1/8, 1/2)$ such that 
$B_{\sup}^\text{low}(q) < \kappa_{\inf}\cdot\sqrt{\frac{2}{\pi}}$ holds for all $q \in (\bar{q}, 1/2)$. 
This completes the proof for (i).
%\yohe{The above derivation solved the issue of continuity. I think the right- or left- continuity of $B$ functions are not anymore needed. Non-increasing property may be necessary for the following argument, but it is easy to verify.}
%{\color{red}TBD: 
%Then, the following inequalities hold.
%\begin{multline}
%    \lim_{q \to 1/2} \left\{\sup_{\theta \in \Theta\setminus\{\theta \mid m = x^*\}} B_\theta^\text{low}(q) \right\}
%    \leq 
%    \lim_{q \to 1/2} \left\{\sup_{\theta \in \Theta\setminus\{\theta \mid m = x^*\}} B_\theta^\mathrm{low}(q; \epsilon=1/2)\right\}
%    \\
%        = 
%        \lim_{q \to 1/2} \left\{
%        \sup_{\theta \in \Theta\setminus\{\theta \mid m = x^*\}}
%        4 \Phi^{-1} \left(1- q + 4 \cdot \mathcal{V}_\mathrm{std}\right)
%        \right\}
%    <
%    4 \Phi^{-1} \left( \Phi\left(\frac{\kappa_{\inf}}{2}\cdot\sqrt{\frac{1}{2\pi}}\right)\right)
%    = \kappa_{\inf}\cdot \sqrt{\frac{2}{\pi}}
%\end{multline}
%Because $B_\theta^\text{low}$ is left-continuous and strictly decreasing, 
%there exists $q_\theta^\mathrm{low}\in(1/8, 1/2)$ such that
%   $\kappa_{\inf}\cdot\sqrt{\frac{2}{\pi}} > B_\theta^\mathrm{low}(q_\theta^\mathrm{low})$ for any $\theta\in\Theta\setminus\{\theta\mid m\neq x^*\}$ .}
%   Then, it holds
%         \begin{equation}
%         \inf\left\{
%         q\in\left(0, \frac12\right):
%         \kappa_{\inf}
%         \cdot
%         \sqrt{\frac{2}{\pi}} > 
%         \underset{\theta \in \Theta\setminus\{\theta \mid m = x^*\}}{\sup}
%         \left\{B_\theta^\mathrm{low}(q) \right\}
%         \right\}
%         % \\
%         % \leq
%         % \inf\left\{
%         % q\in\left(\frac18, \frac12\right):
%         % \kappa_{\inf}
%         % \cdot
%         % \sqrt{\frac{2}{\pi}} > 
%         % \underset{\theta \in \Theta\setminus\{\theta \mid m = x^*\}}{\sup}
%         % \left\{B_\theta^\mathrm{low}(q) \right\}
%         % \right\}
%         < \frac12
%         .
%       \end{equation}
%   Therefore, we see the lower limit of $I_q$ is strictly smaller than $1/2$, and then $I_q$ is non-empty.

Proof of (ii).  
Let $B_\theta^\mathrm{high}(\cdot;\epsilon)$ as defined in \eqref{eq:b_high_q_e}.
To prove that $I_q^\mathrm{high}(q^\mathrm{low})$ is nonempty, it suffices to show there exists $\bar{q} < 1/2$ such that 
$B_\theta^\mathrm{low}(q^\mathrm{low}) > \frac{\aup}{\adown}\cdot B_\theta^\mathrm{high}(q)$ holds for all $\theta \in \Theta_0$ and for all $q \in (\bar{q}, 1/2)$.
Let $\bar{q} = \Phi\left( - \left(\adown / \aup\right)\cdot \Phi^{-1}\left(1 - q^\mathrm{low} \right) \right)$.
Then $\bar{q} \in (0, 1/2)$. 
Because $B_\theta^\mathrm{low}(q) \geq 2\Phi^{-1}(1-q) \geq B_\theta^\mathrm{high}(q)$ for all $q$ for which $B_\theta^\mathrm{low}(q)$ and $B_\theta^\mathrm{high}(q)$ are defined, we have for all $q \in [\bar{q}, 1/2)$ and for all $\theta \in \Theta_0$
\begin{multline}
\frac{\adown}{\aup} \cdot B_\theta^\mathrm{low}(q^\mathrm{low}) 
\geq 2 \frac{\adown}{\aup} \Phi^{-1}(1-q^\mathrm{low}) 
\\
= 2\Phi^{-1}(1-\bar{q}) 
\geq 2\Phi^{-1}(1-q) 
\geq B_\theta^\mathrm{high}(q) .
\end{multline}
This completes the proof of (ii).

%   Since $2\Phi^{-1}(1-q)$ is strictly decreasing with respect to $q$ and $2\Phi^{-1}(1-q) \downarrow 0$ as $q \uparrow 1/2$, there exists a constant $q' \in \left(0, \frac12 \right)$ such that $2\Phi^{-1}(1-q) < (\adown / \aup) \cdot B_\theta^\mathrm{low}(q^\mathrm{low})$ for all $q \in \left(q', \frac12\right)$. Because $B_\theta^\mathrm{high}(q) \leq 2\Phi^{-1}(1-q)$ for any $q \in \left(q', \frac12\right)$, we have $B_\theta^\mathrm{high}(q) < (\adown / \aup) \cdot B_\theta^\mathrm{low}(q^\mathrm{low})$ for all $q \in \left(q', \frac12\right)$.
% Then, we have
% \begin{equation}
% (\max\{q', q^\mathrm{low}\}, 1/2) 
% \subseteq 
% \left\{
%         q \in \left(q^\mathrm{low}, \frac12 \right) :
%         B_\theta^\mathrm{low}(q^\mathrm{low}) 
%         > \frac{\aup}{\adown}\cdot B_\theta^\mathrm{high}(q)
% \right\}
% \end{equation}
% and hence $q^\mathrm{high} \leq \max\{q', q^\mathrm{low}\}$ exists.

Proof of (iii).
For any $q^\mathrm{low} \in I_q$, we have 
$\kappa_{\inf}\cdot\sqrt{\frac{2}{\pi}} > B_{\sup}^\mathrm{low}(q^\mathrm{low})$. 
Therefore, if there exists $\bar{q} > 0$ such that $B_{\inf}^\text{high}(q) \geq\kappa_{\inf}\cdot\sqrt{\frac{2}{\pi}}$ for all $q \in (0, \bar{q}]$, we have $Q \geq \bar{q} > 0$.
Because $\mathcal{V}_\mathrm{std} < \frac18$ under condition \eqref{eq:tr_cond}, $B_\theta^\mathrm{high}(q; \epsilon)$ is defined for any $q \in (0, 1/2)$ and at least for $\epsilon \in \left[\frac{1}{2\sqrt{1 - 2q}}, 1 \right)$. 
Let 
\begin{equation}
\delta_2 = 
 1 - \Phi\left( \frac{\kappa_{\inf}}{2}\cdot\frac{3}{\sqrt{2\pi}}\right)
 - 4\cdot\underset{\theta\in\Thetazero}{\sup}
\left\{\mathcal{V}_\mathrm{std}  \right\}
 .
\end{equation}
Then, it is positive because of \eqref{eq:tr_cond}. 
For any $q \in (0, 1/2)$, we have
\begin{equation}
\begin{aligned}[t]
\label{eq:inf_b_high_lower}
&B_{\inf}^\text{high}(q)
\geq
\underset{\theta \in \Theta_0}{\inf}\left\{
B_\theta^\text{high}\left(q; \epsilon=\frac{1}{2\sqrt{1 - 2q}}\right) \right\}
\\
&=
\underset{\theta \in \Theta_0}{\inf}
\Bigg\{
\frac{4 \sqrt{1 - 2q}}{2 \sqrt{1 - 2q} + 1} 
\cdot \Phi^{-1} \left(1 - q - 4(1 - 2q) \cdot \mathcal{V}_\mathrm{std}\right)  
\Bigg\}
\\
&=
\frac{4 \sqrt{1 - 2q}}{2 \sqrt{1 - 2q} + 1}
\\
&\quad\cdot \Phi^{-1} \left(1 - q - (1 - 2q) \cdot \left(1 - \Phi\left( \frac{\kappa_{\inf}}{2}\cdot\frac{3}{\sqrt{2\pi}}\right) - \delta_2\right) \right) 
.
\end{aligned}
\end{equation}
Because the RMS of \eqref{eq:inf_b_high_lower} is continuous with respect to $q$ and it approaches $\frac{4}{3} \Phi^{-1} \left( \Phi\left( \frac{\kappa_{\inf}}{2}\cdot\frac{3}{\sqrt{2\pi}}\right) + \delta_2\right) >\kappa_{\inf}\cdot\sqrt{\frac{2}{\pi}}$  as $q \to 0$, there exists a $\bar{q} \in (0, 1/2)$ such that 
$B_{\inf}^\text{high}(q) \geq \kappa_{\inf}\cdot\sqrt{\frac{2}{\pi}}$ holds for all $q \in (0, \bar{q}]$. 
This completes the proof of (iii).

Proof of (iv). 
Condition $\bar\sigma < B_{\inf}^\mathrm{high}(q^\mathrm{high})$ implies $\bar\sigma < B_\theta^\mathrm{high}(q^\mathrm{high})$ for all $\theta\in\Thetazero$.
%Condition $\sigma < B_{\inf}^\mathrm{high}(q^\mathrm{high}) \cdot \norm{\nabla f(m)} / \E[\mathcal{Q}_z(\theta)]$ implies $\sigma < B_\theta^\mathrm{high}(q^\mathrm{high}) \cdot \norm{\nabla f(m)} / \E[\mathcal{Q}_z(\theta)]$ for all $\theta\in\Thetazero$.
Then, from \Cref{cor:successprobability}, we immediately find \eqref{eq:lemma:toosmall}.

Proof of (v). 
Condition $\bar\sigma > B_{\sup}^\mathrm{low}(q^\mathrm{low})$ implies
$\bar\sigma > B_\theta^\mathrm{low}(q^\mathrm{low})$ for all $\theta\in\Thetazero$.
%Condition $\sigma > B_{\sup}^\mathrm{low}(q^\mathrm{low}) \cdot \norm{\nabla f(m)} / \E[\mathcal{Q}_z(\theta)]$ implies
%$\sigma > B_\theta^\mathrm{low}(q^\mathrm{low}) \cdot \norm{\nabla f(m)} / \E[\mathcal{Q}_z(\theta)]$ for all $\theta\in\Thetazero$.
Then, from \Cref{cor:successprobability}, we immediately find \eqref{eq:lemma:toolarge}.

Proof of (vi). 
Definition of $Q$ implies that $B_{\sup}^\mathrm{low}(q^\mathrm{low}) \leq B_{\inf}^\mathrm{high}(Q - \xi)$ for any $\xi > 0$. 
Therefore, condition $\bar\sigma \leq B_{\sup}^\mathrm{low}(q^\mathrm{low}) $ implies $\bar\sigma \leq B_{\theta}^\mathrm{high}(Q - \xi)$ for any $\theta\in\Thetazero$.
%Therefore, condition $\sigma \leq B_{\sup}^\mathrm{low}(q^\mathrm{low}) \cdot \norm{\nabla f(m)} / \E[\mathcal{Q}_z(\theta)]$ implies $\sigma \leq B_{\theta}^\mathrm{high}(Q - \xi) \cdot \norm{\nabla f(m)} / \E[\mathcal{Q}_z(\theta)]$ for any $\theta\in\Thetazero$.
From \Cref{cor:successprobability}, we find $\Pr\left[f(m + \sigma z) \leq f(m) \right] \geq Q - \xi$. Because $\xi > 0$ is arbitrary, we obtain \eqref{eq:lemma:reasonable}. 


%           \begin{multline}
%             \E\left[\log\left( \frac{ f(m + \sigma z) }{ f(m) } \right) \ind{f(m + \sigma z) \leq f(m)}\right]
%             \leq \\
%             \frac{\sigma \norm{\nabla f(m)}}{f(m)}\cdot\left(
%             - \frac{1}{\sqrt{2\pi}}
%             + \frac{\sigma \cdot \E\left[\mathcal{Q}_z(\theta)\cdot \ind{z_e\leq 0} \right] }{2\norm{\nabla f(m)}}
%             \right)
%             \cdot \Pr[f(m + \sigma z) \leq f(m) ]
%             .
%           \end{multline}
%   \todo{WIP on the rest}

%   If $\sigma > B_\theta^\mathrm{low}(q^\mathrm{low}) \cdot \norm{\nabla f(m)} / \E[\mathcal{Q}_z(\theta)]$,
%   we find \eqref{eq:lemma:toolarge} immediately from \Cref{cor:successprobability}.

% If $\sigma \geq B_\theta^\mathrm{high}(q^\mathrm{high}) \cdot \sqrt{ 2 L f(m) } / \E[\mathcal{Q}_z(\theta)]$, using the relation $\norm{\nabla f(m)} / \sqrt{2 U} \leq \sqrt{f(m)}$ we have for any $\xi^\text{high} > 0$
% \begin{align*}
%   \frac{ \sigma \E[\mathcal{Q}_z(\theta)] }{ \norm{\nabla f(m)} } \geq \sqrt{\frac{L}{U}}B_\theta^\mathrm{high}(q^\mathrm{high}) \geq B_\theta^\mathrm{low}(Q^\mathrm{high} + \xi^\text{high}).
% \end{align*}
% From \Cref{cor:successprobability}, we find
% \begin{equation}
%   \Pr\left[f(m + \sigma z) \leq f(m) \right] \leq Q^\mathrm{high} + \xi^\text{high}
% \end{equation}
% for any $\xi^\text{high} > 0$. Taking $\inf$ over $\xi^\text{high}$, we obtain the right-hand side of \eqref{eq:lemma:reasonable}.

%   If $\sigma \leq B_\theta^\mathrm{low}(q^\mathrm{low}) \cdot \norm{\nabla f(m)} / \E[\mathcal{Q}_z(\theta)]$,
%   we have
%   \begin{align*}
%     \frac{ \sigma \E[\mathcal{Q}_z(\theta)] }{ \norm{\nabla f(m)} } \leq B_\theta^\mathrm{low}(q^\mathrm{low}) \leq B_\theta^\mathrm{high}(Q - \xi^\text{low})
%   \end{align*}
%   for any $\xi^\text{low} > 0$.
%   From \Cref{cor:successprobability}, we find
%   \begin{equation}
%     \Pr\left[f(m + \sigma z) \leq f(m) \right] \geq Q - \xi^\text{low} .
%   \end{equation}
%   Taking $\inf$ over $\xi^\text{low}$, we obtain the LHS of \eqref{eq:lemma:reasonable}.

%   Under the condition $B_\theta^\mathrm{high}(q^\mathrm{high}) \cdot \sqrt{ 2 L f(m) } / \E[\mathcal{Q}_z(\theta)] \leq \sigma$, we have
%   \begin{align*}
%     \frac{ \sigma \norm{\nabla f(m)} }{ f(m) } \geq B_\theta^\mathrm{high}(q^\mathrm{high}) \frac{ \sqrt{ 2 L } }{ \E[\mathcal{Q}_z(\theta)] } \frac{ \norm{\nabla f(m)} }{ \sqrt{ f(m) } }
%     .
%   \end{align*}
%   Using the relation $\sqrt{f(m)} \leq \norm{\nabla f(m)} / \sqrt{2 L}$ on $L$-storong convexity, we obtain
%   \begin{align*}
%     \frac{ \sigma \norm{\nabla f(m)} }{ f(m) } \geq \frac{ 2 L \cdot B_\theta^\mathrm{high}(q^\mathrm{high}) }{ \E[\mathcal{Q}_z(\theta)] }
%     .
%   \end{align*}
%   Under the condition $\sigma \leq B_\theta^\mathrm{low}(q^\mathrm{low}) \cdot \norm{\nabla f(m)} / \E[\mathcal{Q}_z(\theta)]$, we have
%   \begin{align*}
%     \frac{ \sigma \E[\mathcal{Q}_z(\theta)] }{ \norm{\nabla f(m)} } \leq B_\theta^\mathrm{low}(q^\mathrm{low}).
%   \end{align*}
%   From \Cref{lemma:qualitygain}, we obtain \eqref{eq:lemma:reasonable:qualitygain}.
%   The negativity follows $B_\theta^\mathrm{low}(q^\mathrm{low}) <\kappa_{\inf}\cdot\sqrt{\frac{2}{\pi}}$.

%\hrule

\subsection{Proof of \Cref{lemma:potential_decrease}}\label{apdx:lemma:potential_decrease}

In the following proofs, we use abbreviations of the indicator functions as follows:
% \begin{itemize}
% \item $\inds = \ind{s\sqrt{ Lf(m_{t+1})} \geq  E_\mathcal{Q}\cdot \sigma_{t+1}}$
% \item $\indl = \ind{\ell \sqrt{ L f(m_{t+1})} \leq E_\mathcal{Q} \cdot \sigma_{t+1}}$
% \end{itemize}
$\inds = \ind{s\sqrt{ Lf(m_{t+1})} \geq  E_\mathcal{Q}\cdot \sigma_{t+1}}$ and $\indl = \ind{\ell \sqrt{ L f(m_{t+1})} \leq E_\mathcal{Q} \cdot \sigma_{t+1}}$.
Note that $\inds$ and $\indl$ are exclusive to each other, as $s<\ell$.

Let 
\begin{equation}
\begin{aligned}[t]
    \Delta_s &= \log\left(\frac{\aup}{\adown}\right)\inds\inddown
    -\log\left(\aup\right)\cdot\inds
    \\
    &\quad+ \log\left(\frac{s\sqrt{ Lf(m_{t})}}{\sigma_t\cdot E_\mathcal{Q} }\right)\inds
    - \log^+ \left(\frac{s\sqrt{Lf(m_t)}}{\sigma_t \cdot E_\mathcal{Q}}\right),
    \label{eq:delta_s}
\end{aligned}
\end{equation}
\begin{equation}
\begin{aligned}[t]
    \Delta_{\ell} &=
    \log\left(\frac{\aup}{\adown}\right)\indl\indup
    -\log\left(\frac{1}{\adown}\right)\indl
    \\
    &\quad+ \log\left(\frac{\sigma_t\cdot E_\mathcal{Q}}{\ell\sqrt{ Lf(m_{t})}}\right)\indl
    - \log^+ \left(\frac{\sigma_t \cdot E_\mathcal{Q}}{\ell\sqrt{Lf(m_t)}}\right).
    \label{eq:delta_l}
\end{aligned}
\end{equation}
Then, we can write
\begin{multline}
V(\theta_{t+1}) - V(\theta_t)
\\
=
\left(1+(\inds-\indl)\frac{v}{2}\right)\cdot\log\left(\frac{f(m_{t+1})}{f(m_t)}\right)
+ v \cdot \Delta_s + v \cdot \Delta_\ell.
\label{eq:v_decompose}
\end{multline}

\paragraph{Proof of \eqref{eq:general_potential_decrease} for any $\theta\in\Theta\setminus\{\theta\mid m= x^*\}$}
Clearly, the sums of the last two terms of \eqref{eq:delta_s} and \eqref{eq:delta_l} are both non-positive and the second terms of \eqref{eq:delta_s} and \eqref{eq:delta_l} are both non-positive. 
The sum of the first terms of \eqref{eq:delta_s} and \eqref{eq:delta_l} is upper-bounded by $v\cdot\log(\aup / \adown)$ because both $\log(\aup)$ and $\log(1/\adown)$ are positive and $\inds\inddown + \indl\indup \leq 1$.
Because the first term of \eqref{eq:v_decompose} is non-positive, we have
\begin{align}
V(\theta_{t+1}) - V(\theta_t)
\leq& \left(1-\frac{v}{2}\right)\cdot\log\left(\frac{f(m_{t+1})}{f(m_t)}\right)
+v\cdot\log\left(\frac{\aup}{\adown}\right)
\label{eq:hold_above_in_general_case}
.
\end{align}
This completes the proof of \eqref{eq:general_potential_decrease}.

% \paragraph{Proof of \eqref{eq:general_potential_decrease} for any $\theta\in\Theta\setminus\{\theta\mid m= x^*\}$}

% Extracting $\log(f(m_{t+1})/f(m_t))$, $\log(\aup)$ and $\log(\adown)$ from $V(\theta_{t+1})$,
% %\yohe{I suggest to write the second term on (102) as $-v \cdot \log\left(\frac{1}{\adown}\right)$ so that the sign is clearer.}
% \begin{subequations}
% \begin{align}
% V(\theta_{t+1}) - V(\theta_t)
% =&
% \label{eq:main_terms_of_V_diff_start}
% \left(1+(\inds-\indl)\frac{v}{2}\right)\cdot\log\left(\frac{f(m_{t+1})}{f(m_t)}\right) \\
% \label{eq:main_terms_of_V_diff_2}
% &+v\cdot\log\left(\frac{\aup}{\adown}\right)\cdot\inds\inddown
% -v\cdot\log\left(\aup\right)\cdot\inds \\
% \label{eq:main_terms_of_V_diff_end}
% &+v\cdot\log\left(\frac{\aup}{\adown}\right)\cdot\indl\indup
% -v\cdot\log\left(\frac{1}{\adown}\right)\cdot\indl
% \\
% &+ v\cdot\log\left(\frac{s\sqrt{ Lf(m_{t})}}{\sigma_t\cdot E_\mathcal{Q} }\right)\cdot\inds
% \label{eq:penalty_start}
% \\
% &+ v\cdot\log\left(\frac{\sigma_t\cdot E_\mathcal{Q}}{\ell\sqrt{ Lf(m_{t})}}\right)\cdot\indl
% \label{eq:penalty_start+1}
% \\
% &- v \cdot \log^+ \left(\frac{s\sqrt{Lf(m_t)}}{\sigma_t \cdot E_\mathcal{Q}}\right)
% \label{eq:penalty_end-1}
% \\
% &-  v \cdot \log^+ \left(\frac{\sigma_t \cdot E_\mathcal{Q}}{\ell\sqrt{Lf(m_t)}}\right)
% \label{eq:penalty_end}
% .
% \end{align}
% \end{subequations}
% The sum of \eqref{eq:penalty_start} and \eqref{eq:penalty_end-1} is non-positive as 
% $\text{\eqref{eq:penalty_end-1}} < 0 \Rightarrow  \text{\eqref{eq:penalty_start}} \leq \abs{\text{\eqref{eq:penalty_end-1}}}$ 
% and $\text{\eqref{eq:penalty_end-1}} = 0 \Rightarrow \text{\eqref{eq:penalty_start}} \leq 0$.
% The sum of \eqref{eq:penalty_start+1} and \eqref{eq:penalty_end} is non-positive as 
% $\text{\eqref{eq:penalty_end}} < 0 \Rightarrow \text{\eqref{eq:penalty_start+1}} \leq \abs{\text{\eqref{eq:penalty_end}}}$ 
% and $\text{\eqref{eq:penalty_end}} = 0 \Rightarrow \text{\eqref{eq:penalty_start+1}} \leq 0$.
% %    \yohe{Am I correct that $(103) + (105) \leq 0$ is trivial as $(105) < 0 \Rightarrow (103) \leq \abs{(105)}$ and $(105) = 0 \Rightarrow (103) = 0$?}
% %    We evaluate the sum of \eqref{eq:penalty_start} - \eqref{eq:penalty_end}.
% %  Trivial upper bounds of the terms \eqref{eq:penalty_end-1} and \eqref{eq:penalty_end} are
% %   \begin{align}
% %    - v \cdot \log^+ \left(\frac{s\sqrt{Lf(m_t)}}{\sigma_t \cdot E_\mathcal{Q}}\right)
% %    &\leq
% %    - v \cdot \log^+ \left(\frac{s\sqrt{Lf(m_t)}}{\sigma_t \cdot E_\mathcal{Q}}\right)
% %    \cdot (\inds+\indl)
% %    ,
% %    \\
% %    -  v \cdot \log^+ \left(\frac{\sigma_t \cdot E_\mathcal{Q}}{\ell\sqrt{Lf(m_t)}}\right)
% %    &\leq
% %    -  v \cdot \log^+ \left(\frac{\sigma_t \cdot E_\mathcal{Q}}{\ell\sqrt{Lf(m_t)}}\right)
% %    \cdot (\inds+\indl)
% %      .
% %  \end{align}
% %  %
% %  Next, we evaluate the products of $\inds, \indl, \indup$, and $\inddown$ to inspect the sum of the penalty terms as follows.
% %  %
% %  \begin{align}
% %    \inds\indup
% %    &=\ind{ 1\leq \frac{ s\sqrt{Lf(m_{t+1})} }{E_\mathcal{Q}\aup\sigma_t} }\cdot\indup,
% %    \\
% %    \inds\inddown
% %    &=\ind{ 1\leq \frac{ s\sqrt{Lf(m_{t})} }{E_\mathcal{Q}\adown\sigma_t} }\cdot\inddown \\
% %    &=\left(\ind{ 1\leq \frac{ s\sqrt{Lf(m_{t})} }{\sigma_t \cdot E_\mathcal{Q}} }
% %      +\ind{ \adown \leq \frac{ s\sqrt{Lf(m_{t})} }{\sigma_t \cdot E_\mathcal{Q}} < 1}\right)\cdot\inddown.
% %  \end{align}
% %  %
% %  %
% %  \begin{align}
% %    \indl\indup
% %    &= \ind{1\leq \frac{E_\mathcal{Q}\aup\sigma_t}{\ell \sqrt{Lf(m_{t+1})} } }\cdot\indup\\
% %    &= \ind{1\leq \frac{\sigma_t \cdot E_\mathcal{Q}}{\ell \sqrt{Lf(m_{t+1})} } }\cdot\indup
% %      + \ind{\frac{1}{\aup}\leq \frac{\sigma_t \cdot E_\mathcal{Q}}{\ell \sqrt{Lf(m_{t+1})} } <1 }\cdot\indup\\
% %    &= \ind{1\leq \frac{\sigma_t \cdot E_\mathcal{Q}}{\ell \sqrt{Lf(m_{t})} } }\cdot\indup\\
% %    &+ \ind{\frac{\sigma_t \cdot E_\mathcal{Q}}{\ell \sqrt{Lf(m_{t})} } <1 \leq \frac{\sigma_t \cdot E_\mathcal{Q}}{\ell \sqrt{Lf(m_{t+1})} } }\cdot\indup, \\
% %    &+ \ind{\frac{1}{\aup}\leq \frac{\sigma_t \cdot E_\mathcal{Q}}{\ell \sqrt{Lf(m_{t+1})} } <1 }\cdot\indup
% %    \\
% %    \indl\inddown
% %    &= \ind{1\leq \frac{E_\mathcal{Q}\adown\sigma_t}{\ell \sqrt{Lf(m_{t})} } }\cdot\inddown
% %  \end{align}
% %  %
% %  Here we utilize the following deductions :
% %  $\inds\indup=1 \Rightarrow 1\leq \frac{s \sqrt{Lf(m_t)}}{\sigma_t \cdot E_\mathcal{Q}}$
% %  and
% %  $\indl\inddown=1 \Rightarrow 1\leq \frac{\sigma_t \cdot E_\mathcal{Q}}{\ell \sqrt{Lf(m_t)}}$.
% %  Exploiting these compositions of the indicator functions,  the sum of \eqref{eq:penalty_start}--\eqref{eq:penalty_end} is proved to be non-positive given any $\theta\in\Theta\setminus\{\theta|m\neq x^*\}$ as follows:
% %  %
% %  \begin{equation}
% %    \left(v\cdot\log\left(\frac{s\sqrt{ Lf(m_{t})}}{E_\mathcal{Q} \sigma_{t}}\right)
% %      - v \cdot \log^+ \left(\frac{s\sqrt{Lf(m_t)}}{\sigma_t \cdot E_\mathcal{Q}}\right)\right)\inds\indup
% %        =0
% %       ,
% %    \label{eq:penalty_diff_inds_indup}
% %    \end{equation}
% %  \begin{multline}
% %    \label{eq:penalty_diff_inds_inddown}
% %    \left(v\cdot\log\left(\frac{s\sqrt{ Lf(m_{t})}}{E_\mathcal{Q} \sigma_{t}}\right)
% %      - v \cdot \log^+ \left(\frac{s\sqrt{Lf(m_t)}}{\sigma_t \cdot E_\mathcal{Q}}\right)\right)\inds\inddown\\
% %    =
% %       v\cdot\log\left(\frac{s\sqrt{ Lf(m_{t})}}{E_\mathcal{Q} \sigma_{t}}\right)
% %       \cdot\ind{ \adown \leq \frac{ s\sqrt{Lf(m_{t})} }{\sigma_t \cdot E_\mathcal{Q}} < 1}
% %      \cdot\inddown
% %    \leq 0
% %          ,
% %    \end{multline}
% %    \begin{multline}
% %    \label{eq:penalty_diff_indl_indup}
% %    \left( v\cdot\log\left(\frac{\sigma_t\cdot E_\mathcal{Q}}{\ell\sqrt{ Lf(m_{t})}}\right)
% %      -  v \cdot \log^+ \left(\frac{\sigma_t \cdot E_\mathcal{Q}}{\ell\sqrt{Lf(m_t)}}\right)\right)\indl\indup\\
% %    =
% %        \left( v\cdot\log\left(\frac{\sigma_t\cdot E_\mathcal{Q}}{\ell\sqrt{ Lf(m_{t})}}\right)
% %            - v \cdot \log^+ \left(\frac{\sigma_t \cdot E_\mathcal{Q}}{\ell\sqrt{Lf(m_t)}}\right)
% %            \right)
% %          \\
% %          \quad \cdot
% %          \Biggl(
% %          \ind{\frac{\sigma_t \cdot E_\mathcal{Q}}{\ell \sqrt{Lf(m_{t})} } <1 \leq \frac{\sigma_t \cdot E_\mathcal{Q}}{\ell \sqrt{Lf(m_{t+1})} } }
% %          \\
% %          \quad + \ind{\frac{1}{\aup}\leq \frac{\sigma_t \cdot E_\mathcal{Q}}{\ell \sqrt{Lf(m_{t+1})} } <1 }
% %          \Biggr)
% %            \cdot\indup
% %    \leq 0
% %          ,
% %  \end{multline}
% %    %
% %    \begin{equation}
% %    \left( v\cdot\log\left(\frac{\sigma_t\cdot E_\mathcal{Q}}{\ell\sqrt{ Lf(m_{t})}}\right)
% %      -  v \cdot \log^+ \left(\frac{\sigma_t \cdot E_\mathcal{Q}}{\ell\sqrt{Lf(m_t)}}\right)\right)\indl\inddown
% %    =0
% %       .
% %    \label{eq:penalty_diff_indl_inddown}
% %  \end{equation}
% %  %
% % 
% The sum of \eqref{eq:main_terms_of_V_diff_end} and \eqref{eq:main_terms_of_V_diff_2} is upper-bounded by $v\cdot\log(\aup / \adown)$ because both $\log(\aup)$ and $\log(1/\adown)$ are positive and $\inds\inddown + \indl\indup \leq 1$.
% Subsequently, 
% %\yohe{I feel (107)--(121) are not necessary to derive the following. See the comment above. I suggest to place these relations where they are really used.}
% %\yohe{Here again, I suggest to replace $v \cdot \log(\adown)$ with $v \cdot \log\left(\frac{1}{\adown}\right)$. }
% %
% \begin{align}
% V(\theta_{t+1}) - V(\theta_t)
% % \leq& \left(1+(\inds-\indl)\frac{v}{2}\right)\cdot\log\left(\frac{f(m_{t+1})}{f(m_t)}\right) \\
% % &+v\cdot\log\left(\frac{\aup}{\adown}\right)\cdot\inds\inddown
% %   -v\cdot\log\left(\aup\right)\cdot\inds, \\
% % &+v\cdot\log\left(\frac{\aup}{\adown}\right)\cdot\indl\indup
% %   -v\cdot\log\left(\frac{1}{\adown}\right)\cdot\indl\\
% \leq& \left(1-\frac{v}{2}\right)\cdot\log\left(\frac{f(m_{t+1})}{f(m_t)}\right)
% +v\cdot\log\left(\frac{\aup}{\adown}\right)
% \label{eq:hold_above_in_general_case}
% .
% \end{align}
% This completes the proof of \eqref{eq:general_potential_decrease}.

\paragraph{Proof of \eqref{eq:constant_lower_bound_of_V} for any $\theta\in\Theta\setminus\{\theta\mid m= x^*\}$}
%\yohe{I feel that I can simplify this part significantly. Let me try it below. The result becomes even tighter.}
Note that $\log^+(a \cdot b) \leq \log^+(a) + \log^+(b)$. This implies that $\log^+(b) - \log^+(a \cdot b) \geq - \log^+(a) \geq -\abs{\log(a)}$. Now, let $a = \frac{\sigma_{t+1} \sqrt{ f(m_{t})}}{ \sigma_{t} \sqrt{f(m_{t+1})} }$ and $b = \frac{s\sqrt{ Lf(m_{t+1})}}{\sigma_{t+1}\cdot E_\mathcal{Q} }$. Then, the above inequality implies 
\begin{equation}
\begin{split}
\MoveEqLeft[2]\log^+\left( \frac{s\sqrt{ Lf(m_{t+1})}}{\sigma_{t+1}\cdot E_\mathcal{Q} } \right) - \log^+\left( \frac{s\sqrt{ Lf(m_{t})}}{\sigma_{t}\cdot E_\mathcal{Q} } \right) 
\\
&\geq - \abs*{\log\left( \frac{\sigma_{t+1} \sqrt{ f(m_{t})}}{ \sigma_{t} \sqrt{f(m_{t+1})} } \right)}
\\
&\geq - \abs*{\log\left(\frac{\sigma_{t+1}}{\sigma_t}\right)} - \frac12 \abs*{\log\left( \frac{ f(m_{t+1}) } { f(m_{t}) }\right)}
\\
&> - \log\left(\frac{\aup}{\adown}\right) + \frac12 \log\left( \frac{ f(m_{t+1}) } { f(m_{t}) }\right) .
\end{split}
\end{equation}
Analogously, we have
\begin{multline}
\log^+\left( \frac{\sigma_{t+1} \cdot E_\mathcal{Q}}{\ell\sqrt{Lf(m_{t+1})}} \right) - \log^+\left( \frac{\sigma_t \cdot E_\mathcal{Q}}{\ell\sqrt{Lf(m_t)}} \right) 
\\
> - \log\left(\frac{\aup}{\adown}\right) + \frac12 \log\left( \frac{ f(m_{t+1}) } { f(m_{t}) }\right) .
\end{multline}
Therefore, we obtain
\begin{equation}
V(\theta_{t+1}) - V(\theta_t) > (1 + v) \cdot \log\left( \frac{ f(m_{t+1}) } { f(m_{t}) }\right) - 2 v \cdot \log\left(\frac{\aup}{\adown}\right) .
\end{equation}


\paragraph{Proof of \eqref{eq:small_sigma_potential_decrease} under $\sigma_t < \frac{s\sqrt{L f(m_t)}}{\aup E_\mathcal{Q}}$}

Under this condition, 
we have $\frac{s\sqrt{Lf(m_t)}}{\sigma_t \cdot E_\mathcal{Q}}>\aup>1$
and $\frac{\sigma_t \cdot E_\mathcal{Q}}{\ell\sqrt{Lf(m_t)}} < \frac{\aup \cdot \sigma_t \cdot E_\mathcal{Q}}{\ell\sqrt{Lf(m_t)}}<\frac{s}{\ell}<1$.
The latter implies that the last term of \eqref{eq:delta_l} is $0$.
Note that the sum of the first three terms of \eqref{eq:delta_l} is $v\cdot\log\left(\frac{\sigma_{t+1}\cdot E_\mathcal{Q}}{\ell\sqrt{ Lf(m_{t})}}\right)\cdot\indl$ and $\sigma_{t+1} \leq \aup \cdot \sigma_t$, implying that the sum of the first three terms of \eqref{eq:delta_l} is non-positive. That is, $\Delta_{\ell} \leq 0$. 
Because the first term of \eqref{eq:v_decompose} is always non-positive, we have $V(\theta_{t+1}) - V(\theta_t) \leq v \Delta_s$ and
\begin{equation}
\begin{aligned}[t]
&v \cdot\Delta_s
%\\
%=&
% v\cdot\log\left(\frac{s\sqrt{ Lf(m_{t})}}{E_\mathcal{Q} \aup\sigma_{t}}\right)\cdot\inds\indup
% + v\cdot\log\left(\frac{s\sqrt{ Lf(m_{t})}}{E_\mathcal{Q} \adown\sigma_{t}}\right)\cdot\inds\inddown
% - v \cdot \log \left(\frac{s\sqrt{Lf(m_t)}}{\sigma_t \cdot E_\mathcal{Q}}\right)
% \\
%v\cdot\log\left(\frac{\aup}{\adown}\right)\cdot\inds\inddown
%- v\cdot\log\left(\aup\right)\cdot\inds
%- v \cdot \log \left(\frac{s\sqrt{Lf(m_t)}}{\sigma_t \cdot E_\mathcal{Q}}\right) \cdot (1 - \inds)
\leq
v\cdot\log\left(\frac{\aup}{\adown}\right)\cdot\inds\inddown
- v\cdot\log\left(\aup\right)\cdot\inds
- v \cdot \log \left(\aup\right) \cdot (1 - \inds)
\\
&=
v\cdot\log\left(\frac{\aup}{\adown}\right)\cdot\inds\inddown
- v \cdot \log \left(\aup\right).
%- v \cdot \log \left(\aup\right) \cdot (1 - \inds)
%\\
%&=
%v\cdot\log\left(\frac{\aup}{\adown}\right)\cdot\inds\inddown
%- v \cdot \log \left(\aup\right).
\end{aligned}
\end{equation}
%where we used $\aup < \frac{s\sqrt{L f(m_t)}}{\sigma_t E_\mathcal{Q}}$.  
By using $\inds \leq 1$, we obtain \eqref{eq:small_sigma_potential_decrease}.
%   %
%   The first term of the RHS of the expression above is non-positive under the current condition $\sigma_t < \frac{\aup s\sqrt{L f(m_t)}}{E_\mathcal{Q}}$\yohe{shouldn't it be $\sigma_t < \frac{s\sqrt{L f(m_t)}}{\aup E_\mathcal{Q}}$?}.
%   Moreover, replacing $\inds$ with 1 in the second term provides the following upper bound:
%   %
%   \begin{align}
%     V(\theta_{t+1}) - V(\theta_t)
%     \leq
%     v\cdot\log\left(\frac{\aup}{\adown}\right)\cdot\inddown
%     -v\log(\aup)
%     .
%   \end{align}
%   This completes the proof of \eqref{eq:small_sigma_potential_decrease}.
%   %
%   \yohe{I have shorten the derivation as follows}.


\paragraph*{Proof of \eqref{eq:large_sigma_potential_decrease} under $\sigma_t > \frac{\ell \sqrt{Lf(m_t)}}{\adown  E_\mathcal{Q}}$}

Under this condition, the last term of \eqref{eq:delta_s} is $0$ as $s<\ell$ and $\adown < 1$.
Moreover, we have $\indl = \ind{\ell \sqrt{ L f(m_{t+1})} \leq E_\mathcal{Q} \cdot \sigma_{t+1}} = 1$ and hence $\inds = 0$ because $\sigma_{t+1} \geq \adown\sigma_{t}$ and $\sqrt{ f(m_{t+1})} \leq \sqrt{ L f(m_{t})}$. 
Therefore, the first three terms of \eqref{eq:delta_s} and the sum of the last two terms of \eqref{eq:delta_l} are all $0$. 
Inserting $\indl = 1$ and $\inds = 0$ into \eqref{eq:v_decompose} and \eqref{eq:delta_l}, we obtain
\begin{multline}
V(\theta_{t+1}) - V(\theta_t)
= \left(1-\frac{v}{2}\right)\cdot\log\left(\frac{f(m_{t+1})}{f(m_t)}\right) 
\\
+v\cdot\log\left(\frac{\aup}{\adown}\right)\cdot\indup
-v\cdot\log\left(\frac{1}{\adown}\right)
  ,
\end{multline}
where the first term is always non-positive.
This completes the proof of \eqref{eq:large_sigma_potential_decrease}.
%     Under this condition, the potential function at $t$ is
%   %
%   \begin{align}
%     V(\theta_t)
%     =& \log f(m_t)
%       + v \cdot \log \left(\frac{\sigma_t \cdot E_\mathcal{Q}}{\ell\sqrt{Lf(m_t)}}\right) ,
%   \end{align}
%   as $s<\ell$ and $\adown < 1 < \aup$.
%   %
%   In addition, if $\sigma_t > \frac{\ell \sqrt{L f(m_t)}}{\adown  E_\mathcal{Q}}$, $\indl = 1$ and naturally $\inds = 0$ because
%   %
%   \begin{align}
%     \frac{\ell \sqrt{ L f(m_{t+1})}}{E_\mathcal{Q} \sigma_{t+1}}
%     \leq
%     \frac{\ell \sqrt{ L f(m_{t})}}{E_\mathcal{Q} \adown\sigma_{t}}
%     <1
%     .
%   \end{align}
%   %
%   Consequently, a one-step difference of the potential function under $\sigma_t > \frac{\ell \sqrt{L f(m_t)}}{\adown  E_\mathcal{Q}}$ is
%   %
%   \begin{align}
%     V(\theta_{t+1}) - V(\theta_t)
%     =& \left(1-\frac{v}{2}\right)\cdot\log\left(\frac{f(m_{t+1})}{f(m_t)}\right) \\
%     &+ v\cdot\log\left(\frac{E_\mathcal{Q}\aup\sigma_{t}}{\ell\sqrt{ Lf(m_{t})}}\right)\cdot\indup\\
%     &+ v\cdot\log\left(\frac{E_\mathcal{Q}\adown\sigma_{t}}{\ell\sqrt{ Lf(m_{t})}}\right)\cdot\inddown\\
%     &-  v \cdot \log \left(\frac{\sigma_t \cdot E_\mathcal{Q}}{\ell\sqrt{Lf(m_t)}}\right)
%     \\
%     =& \left(1-\frac{v}{2}\right)\cdot\log\left(\frac{f(m_{t+1})}{f(m_t)}\right) \\
%     &+v\cdot\log\left(\frac{\aup}{\adown}\right)\cdot\indup
%       -v\cdot\log\left(\frac{1}{\adown}\right)
%     \\
%     \leq&
%           v\cdot\log\left(\frac{\aup}{\adown}\right)\cdot\indup
%           -v\cdot\log\left(\frac{1}{\adown}\right)
%           .
%   \end{align}
%   This completes the proof of \eqref{eq:large_sigma_potential_decrease}.
%   %\yohe{I prefer to writing $-\log(1/\adown)$ instead of $\log(\adown)$.}


%%\end{proof}
%%\hrule

\subsection{Proof of \Cref{cor:expected_potential_decrease}}\label{apdx:cor:expected_potential_decrease}
We divide $\Theta\setminus\{\theta\mid m= x^*\}$ into the following three cases:
 (i) $\sigma_t < \frac{s\sqrt{L f(m_t)}}{\aup E_\mathcal{Q}}$; 
(ii) $\sigma_t > \frac{\ell \norm{\nabla f(m_t)}}{\sqrt{2} \adown  E_z[\mathcal{Q}_z(\theta_t)]}$;
(iii) $\frac{s\sqrt{L f(m_t)}}{\aup E_\mathcal{Q}} \leq \sigma_t \leq \frac{\ell \norm{\nabla f(m_t)}}{\sqrt{2} \adown  E_z[\mathcal{Q}_z(\theta_t)]}$.

In the following, we use the fact $w > 0$. 
From \eqref{eq:w}, it is easy to see that $w > 0$ if $\sqrt{\frac{2}{\pi}}\cdot\kappa_{\inf} > B_{\sup}^\mathrm{low}(q^\mathrm{low})$,  
which follows from $q^\mathrm{low} \in I_q$ and the definition of $I_q$, \eqref{eq:def_Iq}.

\paragraph*{(i) $\sigma_t < \frac{s\sqrt{L f(m_t)}}{\aup E_\mathcal{Q}}$}
In light of \Cref{lemma:potential_decrease}, we have \eqref{eq:small_sigma_potential_decrease}.
Moreover, we have
\begin{multline}
\hspace{-0.8em}
\sigma_t < \frac{s\sqrt{L f(m_t)}}{\aup E_\mathcal{Q}}
= \frac{B_{\inf}^\mathrm{high}(q^\mathrm{high}) \cdot \sqrt{2 L f(m_t)}}{E_\mathcal{Q}}
\\
\leq \frac{B_{\inf}^\mathrm{high}(q^\mathrm{high}) \cdot \norm{\nabla f(m_t)}}{E_\mathcal{Q}}
\leq \frac{B_{\inf}^\mathrm{high}(q^\mathrm{high}) \cdot \norm{\nabla f(m_t)}}{\E[\mathcal{Q}_z(\theta)]}.
\end{multline}
Therefore, the condition of (v) in \Cref{lemma:cases} holds and we have
$\E[\inddown] = 1 - \mathrm{Pr}\left[f(m_t + \sigma_t z_t)\leq f(m_t) \mid \mathcal{F}_t\right] \leq 1 - q^\mathrm{high}$.
By taking the expectation on both sides of \eqref{eq:small_sigma_potential_decrease}, we obtain
\begin{equation}
\begin{split}
\MoveEqLeft[1]\E[ V(\theta_{t+1}) - V(\theta_t) \mid \mathcal{F}_t ]
\\
&\leq
v\cdot \log\left(\frac{\aup}{\adown}\right)\cdot \left( 1 - q^\text{high} - \frac{\log(\aup)}{\log(\aup / \adown)} \right)
\\
&=
v\cdot \log\left(\frac{\aup}{\adown}\right)\cdot \left( p_\mathrm{target} - q^\text{high}\right)
\\
&=
\min\left\{ \frac{w}{4} , \log\left(\frac{\aup}{\adown}\right)\right\}
\cdot \left( p_\mathrm{target} - q^\text{high}\right)
<0
\label{eq:small_expected_potential_decrease}
.
\end{split}
\end{equation}

\paragraph*{(ii) $\sigma_t > \frac{\ell \norm{\nabla f(m_t)}}{\sqrt{2}\adown  E_z[\mathcal{Q}_z(\theta_t)]}$}
This condition implies the condition of (ii) in \Cref{lemma:potential_decrease} as we have $\norm{\nabla f(m_t)} \geq \sqrt{2 L f(m_t)}$. 
Therefore, we have \eqref{eq:large_sigma_potential_decrease}. 
Moreover, we have
\begin{equation}
\bar\sigma_t > \ell / \sqrt{2}\adown
= B_{\sup}^\mathrm{low}(q^\mathrm{low}).
\end{equation}
%\begin{equation}
%\sigma_t > \frac{\ell \norm{\nabla f(m_t)}}{\sqrt{2}\adown  E_z[\mathcal{Q}_z(\theta_t)]}
%= \frac{B_{\sup}^\mathrm{low}(q^\mathrm{low}) \norm{\nabla f(m_t)}}{E_z[\mathcal{Q}_z(\theta_t)]}.
%\end{equation}
Therefore, the condition of (vi) in \Cref{lemma:cases} holds and we have
$\E[\indup] = \mathrm{Pr}\left[f(m_t + \sigma_t z_t)\leq f(m_t) \mid \mathcal{F}_t\right] \leq q^\mathrm{low}$.
By taking the expectation on both sides of \eqref{eq:large_sigma_potential_decrease}, we obtain
\begin{multline}
%\MoveEqLeft[2] 
\hspace{-0.8em}
\E[ V(\theta_{t+1}) - V(\theta_t) \mid \mathcal{F}_t ]
%\\
\leq v\cdot\log\left(\frac{\aup}{\adown}\right)\cdot \left( q^\mathrm{low} - \frac{\log(1/\adown)}{\log(\aup/\adown)}\right)
\\
= 
\min\left\{ \frac{w}{4} , \log\left(\frac{\aup}{\adown}\right)\right\}
\cdot \left( q^\mathrm{low} - p_\mathrm{target}\right)
<0
\label{eq:large_expected_potential_decrease}
  .
\end{multline}

\paragraph*{(iii) $\frac{s\sqrt{L f(m_t)}}{\aup E_\mathcal{Q}}\leq\sigma_t\leq\frac{\ell \norm{\nabla f(m_t)}}{\sqrt{2}\adown \E[\mathcal{Q}_z(\theta_t)]}$}
Under this condition, we have
$\frac{\sigma_t \E[\mathcal{Q}_z(\theta_t)\cdot\ind{z_e\leq 0}]}{\norm{\nabla f(m_t)}}\leq \frac{\ell}{\sqrt{2}\adown \kappa_{\inf}}$
and
$\frac{\sigma_t \norm{\nabla f(m_t)}}{f(m_t)}\geq \frac{\sigma_t \sqrt{2L}}{\sqrt{f(m_t)}}\geq \frac{s\sqrt{2} L}{\aup E_\mathcal{Q}}$.
In addition, because
\begin{equation}
\bar\sigma_t
\leq
\ell / \sqrt{2}\adown = B_{\sup}^\mathrm{low}(q^\mathrm{low}),
\end{equation}
%\begin{equation}
%\sigma_t\leq\frac{\ell \norm{\nabla f(m_t)}}{\sqrt{2}\adown \E[\mathcal{Q}_z(\theta_t)]} = \frac{B_{\sup}^\mathrm{low}(q^\mathrm{low}) \norm{\nabla f(m_t)}}{E_z[\mathcal{Q}_z(\theta_t)]},
%\end{equation}
the condition of (iv) in \Cref{lemma:cases} holds and we have $\Pr[f(m + \sigma z) \leq f(m) ]\geq Q$.
Then,
\begin{equation}
\begin{split}
\MoveEqLeft[1] 
\E\left[ \log\left(\frac{f(m_{t+1})}{f(m_t)}\right)  \mid \mathcal{F}_t \right]
\\
&\leq
\frac{\sigma_t \norm{\nabla f(m_t)}}{f(m_t)}\cdot\left(
\frac{\sigma_t \cdot \E\left[\mathcal{Q}_z(\theta_t)\cdot \ind{z_e\leq 0} \right] }{2\norm{\nabla f(m_t)}}
- \frac{1}{\sqrt{2\pi}}
\right)
\\
&\qquad \cdot \Pr_{z_t}[f(m_t + \sigma_t z_t) \leq f(m_t) ]
\\
&\leq
\frac{s\sqrt{2} L}{\aup E_\mathcal{Q}}
\cdot\left(
- \frac{1}{\sqrt{2\pi}}
+ \frac{\ell}{2\sqrt{2}\adown \kappa_{\inf}}
\right)
\cdot Q
% \\
% &=
= 
- w 
.
\end{split}
\end{equation}
Taking the expectation on both sides of \eqref{eq:general_potential_decrease}, we obtain
\begin{multline}
%\MoveEqLeft[2]
\hspace{-0.8em}
\E[ V(\theta_{t+1}) - V(\theta_t)\mid \mathcal{F}_t ]
\leq - \left(1-\frac{v}{2}\right) \cdot w + v\cdot\log\left(\frac{\aup}{\adown}\right)
\\
\leq - \frac{w}{2} + v\cdot\log\left(\frac{\aup}{\adown}\right)
= - \frac{w}{2} + \min \left\{ \frac{w}{4}, \log\left(\frac{\aup}{\adown}\right) \right\}
\leq - \frac{w}{4} .
\label{eq:reasonable_expected_potential_decrease}
\end{multline}

\paragraph*{Altogater}
Finally, by taking a maximum among \eqref{eq:small_expected_potential_decrease}, \eqref{eq:large_expected_potential_decrease} and \eqref{eq:reasonable_expected_potential_decrease}, we obtain $\E[ V(\theta_{t+1}) - V(\theta_t)\mid \mathcal{F}_t ]\leq -B$.
The infimum in $B$ is taken over $q^\mathrm{low}, q^\mathrm{high}\in I_q$ such that $q^\mathrm{low}<p_\mathrm{target}<q^\mathrm{high}$.
The positivity $B$ is guaranteed with $w>0$ and $q^\mathrm{low}<p_\mathrm{target}<q^\mathrm{high}$.
This completes the proof.


% \subsection{Proof of \Cref{cor:expected_multiplicative_potential_decrease}}\label{apdx:cor:expected_multiplicative_potential_decrease}

%     \paragraph{For $\sigma_t < \frac{s \cdot \sqrt{ L f(m_t) }}{ \aup\cdot E_\mathcal{Q}}$.}
%     In light of \Cref{lemma:potential_decrease},
%     if $\sigma_t < \frac{s \cdot \sqrt{ L f(m_t) }}{ \aup\cdot E_\mathcal{Q}}$,
%     $V(\theta_{t+1}) - V(\theta_t)\leq v\cdot\log\left(\frac{\aup}{\adown}\right)\cdot\inddown - v\cdot\log(\aup)$.
%     Then,
%     \begin{equation}
%         \E[\exp\left(\eta\left(V(\theta_{t+1}) - V(\theta_t)\right)\right) \mid \mathcal{F}_t]
%         \leq
%         \E\left[\exp\left(v\eta\cdot\log\left(\frac{\aup}{\adown}\right)\cdot\inddown\right)\right]
%         \cdot\exp\left(-v\eta\cdot\log(\aup)\right)
%         .
%     \end{equation}
%     In the RHS, in light of \Cref{lemma:cases},
%     \begin{align}
%         &\E\left[\exp\left(v\eta\cdot\log\left(\frac{\aup}{\adown}\right)\cdot\inddown\right)\right]
%         \\
%         =&
%         \exp\left(v\eta\cdot\log\left(\frac{\aup}{\adown}\right)\right)\cdot \mathrm{Pr}\left[f(m_t+\sigma_t z_t) > f(m_t)\right]
%         +\mathrm{Pr}\left[f(m_t+\sigma_t z_t) \leq f(m_t)\right]
%         \\
%         =&
%         \left(\exp\left(v\eta\cdot\log\left(\frac{\aup}{\adown}\right)\right)-1\right)\cdot \mathrm{Pr}\left[f(m_t+\sigma_t z_t) > f(m_t)\right]
%         +1
%         \\
%         \leq&
%         \left(\exp\left(v\eta\cdot\log\left(\frac{\aup}{\adown}\right)\right)-1\right)\cdot(1-q^\mathrm{high}) + 1
%         \\
%         =&
%         \left(\exp\left(\eta
%             \cdot\min\left\{\frac{w}{4}, \log\left(\frac{\aup}{\adown}\right)\right\}
%         \right)\cdot(1-q^\mathrm{high}) + q^\mathrm{high}
%         .
%     \end{align}
%     In addition, as $v$ is defined as \eqref{eq:v},
%     \begin{align}
%         &\exp\left(-v\eta\cdot\log(\aup)\right)
%         \\
%         =&
%         \exp\left(-\eta
%             \cdot\left(1-\frac{\log\left(\frac{1}{\adown}\right)}{\log\left(\frac{\aup}{\adown}\right)}\right)
%             \cdot\min\left\{\frac{w}{4}, \log\left(\frac{\aup}{\adown}\right)\right\}
%         \right)
%         \\
%         =&
%         \exp\left(-\eta
%             \cdot\left(1-p_\mathrm{target}\right)
%             \cdot\min\left\{\frac{w}{4}, \log\left(\frac{\aup}{\adown}\right)\right\}
%         \right)
%         \enpace.
%     \end{align}
%     Then, \eqref{eq:small_sigma_expected_multiplicative_potential_decrease} is proved.
%     Let $\eta'=\eta\cdot\min\left\{\frac{w}{4}, \log\left(\frac{\aup}{\adown}\right)\right\}$ and choose $\eta>0$ such that $\eta'<1$.
%     The condition on $q^\mathrm{high}$ for $\rho_s$ to be smaller than 1 is
%     \begin{equation}
%         q^\mathrm{high}
%         >
%         \frac{
%             \exp(\eta') - \exp((1-p_\mathrm{target})\cdot\eta')
%         }{
%             \exp(\eta') - 1
%         }
%         .
%         \label{eq:condition_qhigh}
%     \end{equation}
%     %
%     Using a fact $x+1\leq \exp(x)$ for all $x\in\R$ and $\exp(y)\leq 1/(1-y)$ for all $y\in[0, 1)$,
%     with sufficiently small $\eta>0$,
%     the RHS of the above is upper-bounded as
%     \todo{WIP   }
%     \begin{align}
%         &\frac{
%             \exp\left(v\eta\cdot\log\left(\frac{\aup}{\adown}\right)\right) - \exp\left(v\eta\cdot\log(\aup)\right)
%         }{
%             \exp\left(v\eta\cdot\log\left(\frac{\aup}{\adown}\right)\right)-1
%         }
%         \\
%         \leq&
%         \frac{
%             \frac{1}{1-v\eta\cdot\log\left(\frac{\aup}{\adown}\right)} - \left(v\eta\cdot\log(\aup)+1\right)
%         }{
%             v\eta\cdot\log\left(\frac{\aup}{\adown}\right)
%         }
%         \\
%         =&
%         \frac{
%             v\eta\cdot\log\left(\frac{\aup}{\adown}\right) - v\eta\cdot\log(\aup) + v^2\eta^2\cdot\log(\aup)\cdot\log\left(\frac{\aup}{\adown}\right)
%         }{
%             v\eta\cdot\log\left(\frac{\aup}{\adown}\right)\cdot\left(1-v\eta\cdot\log\left(\frac{\aup}{\adown}\right)\right)
%         }
%         \\
%         =&
%         \frac{
%             \log\left(\frac{\aup}{\adown}\right) - \log(\aup) + v\eta\cdot\log(\aup)\cdot\log\left(\frac{\aup}{\adown}\right)
%         }{
%             \log\left(\frac{\aup}{\adown}\right)\cdot\left(1-v\eta\cdot\log\left(\frac{\aup}{\adown}\right)\right)
%         }
%         \\
%         \leq&
%         \frac{
%             \log\left(\frac{\aup}{\adown}\right) - \log(\aup) + \eta\cdot\log(\aup)\cdot\log\left(\frac{\aup}{\adown}\right)
%         }{
%             \log\left(\frac{\aup}{\adown}\right)\cdot\left(1-\eta\cdot\log\left(\frac{\aup}{\adown}\right)\right)
%         }
%         (\because v\in(0, 1))
%         \\           
%         \to&
%         \frac{
%             \log\left(\frac{1}{\adown}\right)
%         }{
%             \log\left(\frac{\aup}{\adown}\right)
%         }
%         =
%         p_\mathrm{target} \text{ as } \eta\to 0
%         .
%     \end{align}
%     %
%     Remember that $q^\mathrm{high}$ is set so that $q^\mathrm{high}>p_\mathrm{target}$.
%     Therefore, given any $v\in(0, 1)$, there exists $\eta_s>0$ such that if $\eta<\eta_s$, 
%     \eqref{eq:condition_qhigh} holds and then
%     $\rho_s<1$.

%     \paragraph{For $\sigma_t > \frac{\ell \cdot \norm{\nabla f(m_t)} }{ \sqrt{2}\adown\cdot E_{z_t}[\mathcal{Q}_{z_t}(\theta_t)]}$ .}

%     In light of \Cref{lemma:potential_decrease},
%     if $\sigma_t > \frac{\ell \cdot \norm{\nabla f(m_t)} }{ \sqrt{2}\adown\cdot E_{z_t}[\mathcal{Q}_{z_t}(\theta_t)]}$,
%     $V(\theta_{t+1}) - V(\theta_t)\leq v\cdot\log\left(\frac{\aup}{\adown}\right)\cdot\indup - v\cdot\log\left(\frac{1}{\adown}\right)$.
%     Then,
%     \begin{equation}
%         \E[\exp\left(\eta\left(V(\theta_{t+1}) - V(\theta_t)\right)\right) \mid \mathcal{F}_t]
%         \leq
%         \E\left[\exp\left(v\eta\cdot\log\left(\frac{\aup}{\adown}\right)\cdot\indup\right)\right]
%         \cdot\exp\left(-v\eta\cdot\log\left(\frac{1}{\adown}\right)\right)
%         .
%     \end{equation}
%     In the RHS, in light of \Cref{lemma:cases},
%     \begin{align}
%         &\E\left[\exp\left(v\eta\cdot\log\left(\frac{\aup}{\adown}\right)\cdot\indup\right)\right]
%         \\
%         =&
%         \exp\left(v\eta\cdot\log\left(\frac{\aup}{\adown}\right)\right)\cdot\mathrm{Pr}\left[f(m_t+\sigma_t z_t) \leq f(m_t)\right]
%         +\mathrm{Pr}\left[f(m_t+\sigma_t z_t) > f(m_t)\right]
%         \\
%         =&
%         \left(\exp\left(v\eta\cdot\log\left(\frac{\aup}{\adown}\right)\right)-1\right)\cdot\mathrm{Pr}\left[f(m_t+\sigma_t z_t) \leq f(m_t)\right]+1
%         \\
%         \leq&
%         \left(\exp\left(v\eta\cdot\log\left(\frac{\aup}{\adown}\right)\right)-1\right)\cdot q^\mathrm{low}+1
%         \\
%         =&
%         \left(\exp\left(\eta
%             \cdot\min\left\{
%             \frac{w}{4},
%             \log\left(\frac{\aup}{\adown}\right)
%             \right\}
%         \right)\cdot q^\mathrm{low}
%         + 1 - q^\mathrm{low}
%         .
%     \end{align}
%     In addition, with the definition of $v$ \eqref{eq:v},
%     \begin{align}
%         &\exp\left(-v\eta\cdot\log\left(\frac{1}{\adown}\right)\right)
%         \\
%         =&
%         \exp\left(-\eta\cdot
%             \frac{\log\left(\frac{1}{\adown}\right)}{\log\left(\frac{\aup}{\adown}\right)}
%             \cdot\min\left\{
%             \frac{w}{4},
%             \log\left(\frac{\aup}{\adown}\right)
%             \right\}
%         \right)
%         \\
%         =&
%         \exp\left(-\eta\cdot
%             p_\mathrm{target}
%             \cdot\min\left\{
%             \frac{w}{4},
%             \log\left(\frac{\aup}{\adown}\right)
%             \right\}
%         \right)       
%     \end{align}
%     Then, \eqref{eq:large_sigma_expected_multiplicative_potential_decrease} is proved.
%     The condition on $q^\mathrm{low}$ for $\rho_\ell$ to be smaller than 1 is
%     \begin{equation}
%         q^\mathrm{low}
%         <
%         \frac{
%             \left(\exp\left(v\eta\cdot\log\left(\frac{1}{\adown}\right)\right)-1\right)
%         }{
%             \left(\exp\left(v\eta\cdot\log\left(\frac{\aup}{\adown}\right)\right)-1\right)
%         }
%         .
%         \label{eq:condition_qlow}
%     \end{equation}
%     Using a fact $x+1\leq \exp(x)$ for all $x\in\R$ and $\exp(y)\leq 1/(1-y)$ for all $y\in[0, 1)$,
%     with sufficiently small $\eta>0$,
%     the RHS of the above is lower-bounded as
%     \begin{align}
%         &\frac{
%             \left(\exp\left(v\eta\cdot\log\left(\frac{1}{\adown}\right)\right)-1\right)
%         }{
%             \left(\exp\left(v\eta\cdot\log\left(\frac{\aup}{\adown}\right)\right)-1\right)
%         }
%         \\
%         \geq&
%         \frac{
%             v\eta\cdot\log\left(\frac{1}{\adown}\right) \cdot\left( 1-v\eta\cdot\log\left(\frac{\aup}{\adown}\right) \right)
%         }{
%             v\eta\cdot\log\left(\frac{\aup}{\adown}\right)
%         }
%         \\
%         \geq&
%         \frac{
%             \log\left(\frac{1}{\adown}\right) \cdot\left( 1-\eta\cdot\log\left(\frac{\aup}{\adown}\right) \right)
%         }{
%             \log\left(\frac{\aup}{\adown}\right)
%         }
%         \\           
%         \to&
%         p_\mathrm{target} \text{ as } \eta\to 0
%         .
%     \end{align}
%     Remember that $q^\mathrm{low}$ is set so that $q^\mathrm{low}<p_\mathrm{target}$.
%     Therefore, given any $v\in(0, 1)$, there exists $\eta_\ell>0$ such that if $\eta<\eta_\ell$, \eqref{eq:condition_qlow} is satisfied and then $\rho_\ell<1$.

%     \paragraph{For $ \frac{s \cdot \sqrt{ L f(m_t) }}{ \aup\cdot E_\mathcal{Q}} \leq \sigma_t \leq \frac{\ell \cdot \norm{\nabla f(m_t)} }{ \sqrt{2}\adown\cdot E_{z_t}[\mathcal{Q}_{z_t}(\theta_t)]}$ .}

%     In light of \Cref{lemma:potential_decrease},
%     if $ \frac{s \cdot \sqrt{ L f(m_t) }}{ \aup\cdot E_\mathcal{Q}} \leq \sigma_t \leq \frac{\ell \cdot \norm{\nabla f(m_t)} }{ \sqrt{2}\adown\cdot E_{z_t}[\mathcal{Q}_{z_t}(\theta_t)]}$,
%     $V(\theta_{t+1}) - V(\theta_t)\leq \left(1-\frac{v}{2}\right)\cdot\log\left(\frac{f(m_{t+1})}{f(m_t)}\right) + v\cdot \log\left(\frac{\aup}{\adown}\right)$.
%     Then,
%     \begin{multline}
%         \E[\exp\left(\eta\left(V(\theta_{t+1}) - V(\theta_t)\right)\right) \mid \mathcal{F}_t]
%         \\
%         \leq
%         \E\left[\exp\left(\eta\cdot\left(1-\frac{v}{2}\right)\cdot\log\left(\frac{f(m_{t+1})}{f(m_t)}\right)\right)\right]
%         \cdot\exp\left(v\eta\cdot\log\left(\frac{\aup}{\adown}\right)\right)
%         .
%     \end{multline}
%     Since $x\mapsto x^{\eta(1-v/2)}$ is concave over $x>0$ for any $\eta(1-v/2)\in (0, 1)$,
%     then, for any $\frac{f(m_{t+1})}{f(m_t)}>0$,
%     in light of Jensen's inequality,
%     \begin{equation}
%         E\left[\left(\frac{f(m_{t+1})}{f(m_t)}\right)^{\eta(1-v/2)}\right]
%         \leq
%         \left(E\left[\frac{f(m_{t+1})}{f(m_t)}\right]\right)^{\eta(1-v/2)}
%         .
%     \end{equation}
%     In light of \Cref{lemma:qualitygain},
%     if $ \frac{s \cdot \sqrt{ L f(m_t) }}{ \aup\cdot E_\mathcal{Q}} \leq \sigma_t \leq \frac{\ell \cdot \norm{\nabla f(m_t)} }{ \sqrt{2}\adown\cdot E_{z_t}[\mathcal{Q}_{z_t}(\theta_t)]}$,
%     \begin{equation}
%         E\left[\frac{f(m_{t+1})}{f(m_t)}\mid\mathcal{F}_t\right]
%         =
%         1+E\left[\frac{f(m_{t+1})-f(m_t)}{f(m_t)}\mid\mathcal{F}_t\right]
%         \leq
%         1-w
%         ,
%     \end{equation}
%     with $w$ defined as \eqref{eq:w}.
%     With the fact $x\leq\exp(x)-1$ for all $x\in\R$,
%     \begin{equation}
%         \left(1-w\right)^{\eta(1-v/2)}
%         \cdot\exp\left(v\eta\cdot\log\left(\frac{\aup}{\adown}\right)\right)
%         \leq
%         \exp\left(-w\eta\left(1-\frac{v}{2}\right) + v\eta\cdot\log\left(\frac{\aup}{\adown}\right)\right)
%         .
%     \end{equation}
%     As $v$ is defined as \eqref{eq:v}, the RHS is
%     \begin{align}
%         &\exp\left(-w\eta\left(1-\frac{v}{2}\right) + v\eta\cdot\log\left(\frac{\aup}{\adown}\right)\right)
%         \\
%         \leq&
%         \exp\left(-\frac12 w\eta 
%         + \min\left\{\frac{w\eta}{4},  \eta\cdot\log\left(\frac{\aup}{\adown}\right)\right\}\right)
%         \\
%         \leq&
%         \exp\left(-\frac14 w \eta\right)
%         \\
%         \leq&
%         \exp\left(-\eta\cdot\min\left\{\frac{w}{4}, \log\left(\frac{\aup}{\adown}\right)\right\}\right)
%         \\
%         \leq&
%         \exp\left(-\eta\cdot p_\mathrm{target}\cdot\min\left\{\frac{w}{4}, \log\left(\frac{\aup}{\adown}\right)\right\}\right)
%         \\       
%         =&
%         \rho_r
%         .
%     \end{align}
%     This completes the proof.
%     

\subsection{Proof of \Cref{lemma:v_variance_bound}}\label{apdx:lemma:v_variance_bound}
%\begin{proof}

Note first that
%\begin{equation}
${\color{blue}
\Var[V(\theta_{t+1})]
\leq
\E[\left(V(\theta_{t+1}) - V(\theta_t)\right)^2]}$
%\end{equation}
%
holds. 
% In light of \eqref{eq:general_potential_decrease},
% %
% % \begin{align}
% % V(\theta_{t+1}) - V(\theta_t) \leq v \cdot \log\left(\frac{\aup}{\adown}\right)
% % \label{eq:constant_upper_bound_of_V}
% % .
% % \end{align}
% we have $V(\theta_{t+1}) - V(\theta_t) \leq v \cdot \log\left(\frac{\aup}{\adown}\right)$.
% Combining this inequailty and \eqref{eq:constant_lower_bound_of_V},
From \eqref{eq:general_potential_decrease} and \eqref{eq:constant_lower_bound_of_V}, 
we obtain
% \begin{equation}
% \begin{aligned}[t]
% \MoveEqLeft[2]\abs{V(\theta_{t+1}) - V(\theta_t)}
% \\
% % &\leq
% % \max
% % \Biggl\{
% % v \cdot \log\left(\frac{\aup}{\adown}\right),
% % \left(1+v\right)\cdot\left|\log\left(\frac{f(m_{t+1})}{f(m_t)}\right)\right|
% % +2v\cdot\log\left(\frac{\aup}{\adown}\right)
% % \Biggr\}
% % \\
% &\leq
% \left(1+v\right)\cdot\left|\log\left(\frac{f(m_{t+1})}{f(m_t)}\right)\right|
% +2v\cdot\log\left(\frac{\aup}{\adown}\right)
% .
% \end{aligned}
% \end{equation}
\begin{equation}
\abs{V(\theta_{t+1}) - V(\theta_t)}
\leq
\left(1+v\right)\cdot\left|\log\left(\frac{f(m_{t+1})}{f(m_t)}\right)\right|
+2v\cdot\log\left(\frac{\aup}{\adown}\right)
.
\end{equation}
%
As the only random variable in the RHS of the expression above is $\left|\log\left(\frac{f(m_{t+1})}{f(m_t)}\right)\right|$,
to prove ${\color{blue}\E[\left(V(\theta_{t+1}) - V(\theta_t)\right)^2]} < \infty$, it suffices to prove that
% \begin{align}
% \E\left[\left|\log\left(\frac{f(m_{t+1})}{f(m_t)}\right)\right|\mid\mathcal{F}_t\right]
% <\infty
% ,
% \label{eq:abs_qualitygain_integrable}
% \end{align}
% and
% \begin{align}
% \E\left[\left(\log\left(\frac{f(m_{t+1})}{f(m_t)}\right)\right)^2\mid\mathcal{F}_t\right]
% <\infty
% .
% \label{eq:squared_qualitygain_integrable}
% \end{align}
%\begin{align}
${\color{blue}\E\left[\abs*{\log\left(\frac{f(m_{t+1})}{f(m_t)}\right)}\right]}
<\infty$
%\quad \text{and} \quad 
and
${\color{blue}\E\left[\left(\log\left(\frac{f(m_{t+1})}{f(m_t)}\right)\right)^2\right]}
<\infty$.
%\label{eq:qualitygain_integrable}
%\end{align}
%
By exploiting the fact that $x \leq \exp(x) - 1$ and $x^2 \leq 2 (\exp(x) - 1)$,
the above inequalities are straightforward for $d>3$ with \Cref{lemma:variancebound}.
% Namely,
% \begin{align}
% \E\left[\left|\log\left(\frac{f(m_{t+1})}{f(m_t)}\right)\right|\mid\mathcal{F}_t\right]
% \leq
% \frac{U}{L}\cdot\left(1+\frac{1}{d-3}\right) -1
% <\infty
% ,
% \end{align}
% and
% \begin{align}
% \E\left[\left(\log\left(\frac{f(m_{t+1})}{f(m_t)}\right)\right)^2\mid\mathcal{F}_t\right]
% \leq
% 2\cdot\left(\frac{U}{L}\cdot\left(1+\frac{1}{d-3}\right) -1\right)
% <\infty
% .
% \end{align}
This completes the proof. 
%\end{proof}



%{\appendices
%\section*{Proof of the First Zonklar Equation}
%Appendix one text goes here.
% You can choose not to have a title for an appendix if you want by leaving the argument blank
%\section*{Proof of the Second Zonklar Equation}
%Appendix two text goes here.}


% \bibliographystyle{plain}
% \bibliography{algorithmica}
%
\begin{thebibliography}{10}

\bibitem{akimoto2018drift}
Youhei Akimoto, Anne Auger, and Tobias Glasmachers.
\newblock Drift theory in continuous search spaces: expected hitting time of
  the {(1+1)-ES} with 1/5 success rule.
\newblock In {\em Proceedings of the Genetic and Evolutionary Computation
  Conference}, pages 801--808, 2018.

\bibitem{akimoto2020global}
Youhei Akimoto, Anne Auger, Tobias Glasmachers, and Daiki Morinaga.
\newblock Global linear convergence of evolution strategies on more than smooth
  strongly convex functions.
\newblock {\em SIAM Journal on Optimization}, 2022.

\bibitem{akimoto2020tcs}
Youhei Akimoto, Anne Auger, and Nikolaus Hansen.
\newblock Quality gain analysis of the weighted recombination evolution
  strategy on general convex quadratic functions.
\newblock {\em Theor. Comput. Sci.}, 832:42--67, 2020.

\bibitem{auger2013linear}
Anne Auger and Nikolaus Hansen.
\newblock Linear convergence on positively homogeneous functions of a
  comparison based step-size adaptive randomized search: the {(1+1) ES} with
  generalized one-fifth success rule.
\newblock {\em arXiv preprint arXiv:1310.8397}, 2013.

\bibitem{beyer2013dynamics}
Hans-Georg Beyer and Alexander Melkozerov.
\newblock The dynamics of self-adaptive multirecombinant evolution strategies
  on the general ellipsoid model.
\newblock {\em IEEE transactions on evolutionary computation}, 18(5):764--778,
  2013.

\bibitem{chen1982inequality}
Louis~HY Chen.
\newblock An inequality for the multivariate normal distribution.
\newblock {\em Journal of Multivariate Analysis}, 12(2):306--315, 1982.

\bibitem{chow1967strong}
Yuan~Shih Chow et~al.
\newblock On a strong law of large numbers for martingales.
\newblock {\em The Annals of Mathematical Statistics}, 38(2):610--610, 1967.

\bibitem{dong2019efficient}
Yinpeng Dong, Hang Su, Baoyuan Wu, Zhifeng Li, Wei Liu, Tong Zhang, and Jun
  Zhu.
\newblock Efficient decision-based black-box adversarial attacks on face
  recognition.
\newblock In {\em Proceedings of the IEEE Conference on Computer Vision and
  Pattern Recognition}, pages 7714--7722, 2019.

\bibitem{fujii2018cma}
Garuda Fujii, Masayuki Takahashi, and Youhei Akimoto.
\newblock Cma-es-based structural topology optimization using a level set
  boundary expression—application to optical and carpet cloaks.
\newblock {\em Computer Methods in Applied Mechanics and Engineering},
  332:624--643, 2018.

\bibitem{glasmachers2020global}
Tobias Glasmachers.
\newblock Global convergence of the {(1+1)} evolution strategy to a critical
  point.
\newblock {\em Evolutionary computation}, 28(1):27--53, 2020.

\bibitem{2019gradientlesgolovins}
Daniel Golovin, John Karro, Greg Kochanski, Chansoo Lee, Xingyou Song, and
  Qiuyi Zhang.
\newblock Gradientless descent: High-dimensional zeroth-order optimization.
\newblock In {\em International Conference on Learning Representations}, 2019.

\bibitem{hansen2014principled}
Nikolaus Hansen and Anne Auger.
\newblock Principled design of continuous stochastic search: From theory to
  practice.
\newblock In {\em Theory and principled methods for the design of
  metaheuristics}, pages 145--180. Springer, 2014.

\bibitem{hansen2003reducing}
Nikolaus Hansen, Sibylle~D M{\"u}ller, and Petros Koumoutsakos.
\newblock Reducing the time complexity of the derandomized evolution strategy
  with covariance matrix adaptation (cma-es).
\newblock {\em Evolutionary computation}, 11(1):1--18, 2003.

\bibitem{hansen2001completely}
Nikolaus Hansen and Andreas Ostermeier.
\newblock Completely derandomized self-adaptation in evolution strategies.
\newblock {\em Evolutionary computation}, 9(2):159--195, 2001.

\bibitem{hardy1952inequalities}
Godfrey~Harold Hardy, John~Edensor Littlewood, George P{\'o}lya, Gy{\"o}rgy
  P{\'o}lya, et~al.
\newblock {\em Inequalities}.
\newblock Cambridge university press, 1952.

\bibitem{jagerskupper2003analysis}
Jens J{\"a}gersk{\"u}pper.
\newblock Analysis of a simple evolutionary algorithm for minimization in
  euclidean spaces.
\newblock In {\em International Colloquium on Automata, Languages, and
  Programming}, pages 1068--1079. Springer, 2003.

\bibitem{jagerskupper20061+}
Jens J{\"a}gersk{\"u}pper.
\newblock How the {(1+1) ES} using isotropic mutations minimizes positive
  definite quadratic forms.
\newblock {\em Theoretical Computer Science}, 361(1):38--56, 2006.

\bibitem{jagerskupper2007algorithmic}
Jens J{\"a}gersk{\"u}pper.
\newblock Algorithmic analysis of a basic evolutionary algorithm for continuous
  optimization.
\newblock {\em Theoretical Computer Science}, 379(3):329--347, 2007.

\bibitem{Karimi2016}
Hamed Karimi, Julie Nutini, and Mark Schmidt.
\newblock Linear convergence of gradient and proximal-gradient methods under
  the polyak-{\l}ojasiewicz condition.
\newblock In {\em Machine Learning and Knowledge Discovery in Databases:
  European Conference, ECML PKDD 2016, Riva del Garda, Italy, September 19-23,
  2016, Proceedings, Part I 16}, pages 795--811. Springer, 2016.

\bibitem{kern2004learning}
Stefan Kern, Sibylle~D M{\"u}ller, Nikolaus Hansen, Dirk B{\"u}che, Jiri
  Ocenasek, and Petros Koumoutsakos.
\newblock Learning probability distributions in continuous evolutionary
  algorithms--a comparative review.
\newblock {\em Natural Computing}, 3(1):77--112, 2004.

\bibitem{kriest2017calibrating}
Iris Kriest, Volkmar Sauerland, Samar Khatiwala, Anand Srivastav, and Andreas
  Oschlies.
\newblock Calibrating a global three-dimensional biogeochemical ocean model
  (mops-1.0).
\newblock {\em Geoscientific Model Development}, 10:127--154, 2017.

\bibitem{morinaga2019generalized}
Daiki Morinaga and Youhei Akimoto.
\newblock Generalized drift analysis in continuous domain: linear convergence
  of {(1+ 1)-ES} on strongly convex functions with lipschitz continuous
  gradients.
\newblock In {\em Proceedings of the 15th ACM/SIGEVO Conference on Foundations
  of Genetic Algorithms}, pages 13--24, 2019.

\bibitem{morinaga2021convergence}
Daiki Morinaga, Kazuto Fukuchi, Jun Sakuma, and Youhei Akimoto.
\newblock Convergence rate of the {(1+1)}-evolution strategy with success-based
  step-size adaptation on convex quadratic functions.
\newblock In {\em Proceedings of the Genetic and Evolutionary Computation
  Conference}, pages 1169--1177, 2021.

\bibitem{nesterov2017random}
Yurii Nesterov and Vladimir Spokoiny.
\newblock Random gradient-free minimization of convex functions.
\newblock {\em Foundations of Computational Mathematics}, 17(2):527--566, 2017.

\bibitem{rechenberg1973evolution}
Ingo Rechenberg.
\newblock Evolution strategy: Optimization of technical systems by means of
  biological evolution.
\newblock {\em Fromman-Holzboog, Stuttgart}, 104:15--16, 1973.

\bibitem{uhlendorf2012long}
Jannis Uhlendorf, Agn{\`e}s Miermont, Thierry Delaveau, Gilles Charvin,
  Fran{\c{c}}ois Fages, Samuel Bottani, Gregory Batt, and Pascal Hersen.
\newblock Long-term model predictive control of gene expression at the
  population and single-cell levels.
\newblock {\em Proceedings of the National Academy of Sciences},
  109(35):14271--14276, 2012.

\bibitem{varelas2018comparative}
Konstantinos Varelas, Anne Auger, Dimo Brockhoff, Nikolaus Hansen, Ouassim~Ait
  ElHara, Yann Semet, Rami Kassab, and Fr{\'e}d{\'e}ric Barbaresco.
\newblock A comparative study of large-scale variants of cma-es.
\newblock In {\em International Conference on Parallel Problem Solving from
  Nature}, pages 3--15. Springer, 2018.

\end{thebibliography}


\section{Biography Section}
%If you have an EPS/PDF photo (graphicx package needed), extra braces are
% needed around the contents of the optional argument to biography to prevent
% the LaTeX parser from getting confused when it sees the complicated
% $\backslash${\tt{includegraphics}} command within an optional argument. (You can create
% your own custom macro containing the $\backslash${\tt{includegraphics}} command to make things
% simpler here.)
 
% \vspace{11pt}

%\bf{Daiki Morinaga}\vspace{-33pt}
\begin{IEEEbiography}[{\includegraphics[width=1in,height=1.25in,clip,keepaspectratio]{morinaga.jpeg}}]{Daiki Morinaga}
received the B.S. degree in information science and the M.S. degree in engineering from University of Tsukuba, Tsukuba, Japan, in 2020 and 2022, respectively.
He won the Best Paper Award at FOGA 2019.
His research interests include optimization techniques in engineering, such as black-box optimization and metaheuristics, and their theoretical analyses.
\end{IEEEbiography}

\begin{IEEEbiography}[{\includegraphics[width=1in,height=1.25in,clip,keepaspectratio]{fukuchi.jpg}}]{Kazuto Fukuchi} received a Ph.D. degree from University of Tsukuba, Japan, in 2018. Since 2019, he has been an Assistant Professor with the Faculty of Engineering, Information and Systems, University of Tsukuba, Japan, and also a Visiting Researcher with the Center for Advanced Intelligence Project, RIKEN, Japan. His research interests include mathematical statistics, machine learning, and their application.
\end{IEEEbiography}

\begin{IEEEbiography}[{\includegraphics[width=1in,height=1.25in,clip,keepaspectratio]{sakuma.png}}]{Jun Sakuma} received the Ph.D. degree in engineering from Tokyo
Institute of Technology, Tokyo, Japan, in 2003. Since 2016, he has been a Professor
with the Department of Computer Science, School of Systems and
Information Engineering, University of Tsukuba, Japan. He has also been Team Leader of the Artificial Intelligence Security and Privacy team in the Center for Advanced Intelligence
Project, RIKEN since 2016. From 2009 to 2016, he was an Associate Professor with the Department of Computer Science, University of Tsukuba. From 2004 to 2009, he was an Assistant Professor with the Department of Computational Intelligence and Systems Science, Interdisciplinary Graduate School of Science and Engineering, Tokyo Institute of
Technology, Tokyo, Japan. From 2003 to 2004, he was a Researcher with
Tokyo Research Laboratory, IBM, Tokyo, Japan. His research
interests include data mining, machine learning, data privacy, and
security. He is a member of the Institute of Electronics, Information
and Communication Engineers of Japan (IEICE).
\end{IEEEbiography}

\begin{IEEEbiography}[{\includegraphics[width=1in,height=1.25in,clip,keepaspectratio]{akimoto.jpg}}]{Youhei Akimoto} received the B.S. degree in computer science in 2007, and the M.S. and Ph.D. degrees in computational intelligence and systems science from Tokyo Institute of Technology, Japan, in 2008 and 2011, respectively. From 2010 to 2011, he was a Research Fellow of JSPS in Japan, and from 2011 to 2013, he was a Post-Doctoral Research Fellow with INRIA in France. From 2013 to 2018, he was an Assistant Professor with Shinshu University, Japan. Since 2018, he has been an Associate Professor with University of Tsukuba, Japan as well as a Visiting Researcher with the Center for Advanced Intelligence Projects, RIKEN. 
%He served as a Track Chair for the continuous optimization track of GECCO in 2015 and 2016. 
%He is an Associate Editor of ACM TELO and is on the editorial board of the ECJ. He won the Best Paper Award at GECCO 2018 and FOGA 2019. 
His research interests include design principles, theoretical analyses, and applications of stochastic search heuristics.
\end{IEEEbiography}



% \bf{If you include a photo:}\vspace{-33pt}
% \begin{IEEEbiography}[{\includegraphics[width=1in,height=1.25in,clip,keepaspectratio]{fig1}}]{Michael Shell}
% Use $\backslash${\tt{begin\{IEEEbiography\}}} and then for the 1st argument use $\backslash${\tt{includegraphics}} to declare and link the author photo.
% Use the author name as the 3rd argument followed by the biography text.
% \end{IEEEbiography}

%\vspace{11pt}

%\bf{If you will not include a photo:}\vspace{-33pt}
%\begin{IEEEbiographynophoto}{John Doe}
%Use $\backslash${\tt{begin\{IEEEbiographynophoto\}}} and the author name as the argument followed by the biography text.
%\end{IEEEbiographynophoto}
\vfill

\end{document}


